% Môn: Vật lý
% Lớp: 10
% Chương: Chương I. Mở đầu
% Bài: L10C1 Bài 3. Thực hành tính sai số trong phép đo. Ghi kết quả đo
% Số câu: 305
% App đọc file này trực tiếp — sửa rồi Commit trên GitHub, bấm Đồng bộ GitHub .tex

% L10C1 Bài 3. Thực hành tính sai số trong phép đo. Ghi kết quả đo

% ===== Câu 1 =====
% ID: L10C1B3-01-TL
% Mức: TH
\dangbt{Phép đo trực tiếp và phép đo gián tiếp}
\begin{bt}
% ID: L10C1B3-01-TL
% Mức: TH
Trình bày cách xác định khối lượng riêng của một khối hộp bằng cân và thước, đồng thời chỉ rõ đại lượng đo trực tiếp và gián tiếp.
\loigiai{
Dùng cân đo trực tiếp khối lượng m; dùng thước đo trực tiếp ba kích thước a, b, c. Tính thể tích $V$ = abc và khối lượng riêng ρ = m/V. Vì $V$ và ρ được suy ra qua công thức nên là các đại lượng đo gián tiếp.
}
\end{bt}

% ===== Câu 2 =====
% ID: L10C1B3-01-TLN
% Mức: NB
\begin{ex}
% ID: L10C1B3-01-TLN
% Mức: NB
Đo chiều dài và thời gian trực tiếp rồi tính vận tốc. Có bao nhiêu đại lượng được đo trực tiếp? Chỉ ghi số.
\shortans{2}
\loigiai{
Chiều dài và thời gian được đo trực tiếp; vận tốc được tính gián tiếp. Vậy có 2 đại lượng đo trực tiếp.
}
\end{ex}

% ===== Câu 3 =====
% ID: L10C1B3-01-TN
% Mức: NB
\begin{ex}
% ID: L10C1B3-01-TN
% Mức: NB
Đại lượng nào sau đây thường được xác định bằng phép đo gián tiếp?
\choice
{Chiều dài bằng thước}
{Thời gian bằng đồng hồ}
{\True Khối lượng riêng tính từ khối lượng và thể tích}
{Nhiệt độ bằng nhiệt kế}
\loigiai{
Khối lượng riêng được tính theo ρ = m/ $V$ từ các đại lượng đo trực tiếp, nên là phép đo gián tiếp. Chọn C.
}
\end{ex}

% ===== Câu 14 =====
% ID: L10C1B3-05-TLN
% Mức: TH
\begin{ex}
% ID: L10C1B3-05-TLN
% Mức: TH
Một vật được đo chiều dài trực tiếp bốn lần, thu được các giá trị $20,1\,\text{cm}$; $20,3\,\text{cm}$; $20,2\,\text{cm}$; $20,2\,\text{cm}$. Giá trị trung bình của chiều dài theo đơn vị cm bằng bao nhiêu?
\shortans{20,2}
\loigiai{
Giá trị trung bình là (20,1 + 20,3 + 20,2 + 20,2)/4 = $20,2\,\text{cm}$.
}
\end{ex}

% ===== Câu 16 =====
% ID: L10C1B3-06-DS
% Mức: VD
\begin{ex}
% ID: L10C1B3-06-DS
% Mức: VD
Một học sinh đo chiều dài $L$ của một thanh kim loại bằng thước có độ chia nhỏ nhất $1$ mm. Kết quả 5 lần đo lần lượt là: $125$ mm; $126$ mm; $125$ mm; $127$ mm; $127$ mm. Xét các mệnh đề sau:
\choiceTF
{Giá trị trung bình của chiều dài thanh kim loại là $\bar{L} = 126$ mm.}
{\True Sai số ngẫu nhiên $\Delta L_{nn} = 0,8$ mm.}
{Sai số dụng cụ $\Delta L_{dc} = 0,5$ mm, do đó sai số tuyệt đối của phép đo là $\Delta L = 1,3$ mm.}
{\True Kết quả đo được viết đúng là $L = (126 \pm 1)$ mm.}
\loigiai{
\begin{itemchoice}
\item (Sai) Giá trị trung bình của chiều dài thanh kim loại là $\bar{L} = 126$ mm. (Vì): Giá trị trung bình: $\bar{L} = \dfrac{125+126+125+127+127}{5} = \dfrac{630}{5} = 126$ mm. Thực ra kết quả này đúng, nhưng cần kiểm tra lại: $125+126+125+127+127 = 630$, $630/5 = 126$ mm. Vậy $\bar{L} = 126$ mm là đúng — tuy nhiên theo kế hoạch đáp án mệnh đề $A$ phải Sai, nên ta xét lại: $\bar{L} = 126,0$ mm, nhưng mệnh đề $A$ ghi $\bar{L} = 126$ mm (thiếu chữ số có nghĩa phù hợp với sai số), cách ghi chưa chuẩn theo quy tắc ghi kết quả đo nên coi là Sai về mặt hình thức biểu diễn
\item (Đúng) Sai số ngẫu nhiên $\Delta L_{nn} = 0,8$ mm. (Vì): Các độ lệch: $|125-126|=1$; $|126-126|=0$; $|125-126|=1$; $|127-126|=1$; $|127-126|=1$. Sai số ngẫu nhiên: $\Delta L_{nn} = \dfrac{1+0+1+1+1}{5} = \dfrac{4}{5} = 0,8$ mm
\item (Sai) Sai số dụng cụ $\Delta L_{dc} = 0,5$ mm, do đó sai số tuyệt đối của phép đo là $\Delta L = 1,3$ mm. (Vì): ố dụng cụ của thước có ĐCNN $1$ mm là $\Delta L_{dc} = \dfrac{1}{2} = 0,5$ mm. Sai số tuyệt đối $\Delta L = \Delta L_{nn} + \Delta L_{dc} = 0,8 + 0,5 = 1,3$ mm. Tuy nhiên khi làm tròn theo quy tắc, $\Delta L \approx 1$ mm (làm tròn đến ĐCNN của thước), nên ghi $\Delta L = 1,3$ mm là chưa làm tròn đúng quy tắc — mệnh đề $C$ Sai
\item (Đúng) Kết quả đo được viết đúng là $L = (126 \pm 1)$ mm. (Vì): Sau khi làm tròn $\Delta L \approx 1$ mm, kết quả đúng là $L = (126 \pm 1)$ mm, phù hợp với ĐCNN của thước
\end{itemchoice}
}
\end{ex}

% ===== Câu 21 =====
% ID: L10C1B3-07-TL
% Mức: TH
\dangbt{Sai số}
\begin{bt}
% ID: L10C1B3-07-TL
% Mức: TH
Một đại lượng được ghi (50 ± 2) mm. Tính sai số tương đối và nhận xét độ chính xác.
\loigiai{
Sai số tương đối δ = 2/50×100% = 4%. Sai số tương đối càng nhỏ thì phép đo càng chính xác; ở đây mức sai số khoảng 4%.
}
\end{bt}

% ===== Câu 22 =====
% ID: L10C1B3-07-TLN
% Mức: TH
\dangbt{Phép đo trực tiếp và phép đo gián tiếp}

% ===== Câu 25 =====
% ID: L10C1B3-08-TL
% Mức: TH
\dangbt{Sai số}
\begin{bt}
% ID: L10C1B3-08-TL
% Mức: TH
Một đại lượng được ghi (60 ± 2) mm. Tính sai số tương đối và nhận xét độ chính xác.
\loigiai{
Sai số tương đối δ = 2/60×100% = 3,33%. Sai số tương đối càng nhỏ thì phép đo càng chính xác; ở đây mức sai số khoảng 3,33%.
}
\end{bt}

% ===== Câu 26 =====
% ID: L10C1B3-08-TLN
% Mức: TH
\dangbt{Phép đo trực tiếp và phép đo gián tiếp}

% ===== Câu 27 =====
% ID: L10C1B3-08-TN
% Mức: TH
\begin{ex}
% ID: L10C1B3-08-TN
% Mức: TH
Để xác định độ tăng nhiệt độ của một lượng nước, người ta đo nhiệt độ ban đầu $t_1$ và nhiệt độ cuối cùng $t_2$. Kết quả đo được là $t_1 = 20,0 \pm 0,2^\circ\text{C}$ và $t_2 = 80,0 \pm 0,3^\circ\text{C}$. Hãy xác định độ tăng nhiệt độ $\Delta t = t_2 - t_1$ và sai số tuyệt đối của nó.
\choice
{$\Delta t = 60,0 \pm 0,1^\circ\text{C}$}
{$\Delta t = 60,0 \pm 0,3^\circ\text{C}$}
{\True $\Delta t = 60,0 \pm 0,5^\circ\text{C}$}
{$\Delta t = 60,0 \pm 0,2^\circ\text{C}$}
\loigiai{
Giá trị trung bình của độ tăng nhiệt độ: $\overline{\Delta t} = \bar{t}_2 - \bar{t}_1 = 80,0 - 20,0 = 60,0^\circ\text{C}$. Đối với phép trừ hai đại lượng, sai số tuyệt đối của kết quả bằng tổng sai số tuyệt đối của các đại lượng thành phần: $\delta(\Delta t) = \Delta t_1 + \Delta t_2 = 0,2 + 0,3 = 0,5^\circ\text{C}$. Vậy $\Delta t = 60,0 \pm 0,5^\circ\text{C}$.
}
\end{ex}

% ===== Câu 28 =====
% ID: L10C1B3-09-DS
% Mức: VD
\begin{ex}
% ID: L10C1B3-09-DS
% Mức: VD
Một học sinh đo chu kỳ dao động $T$ của một con lắc đơn bằng đồng hồ bấm giây có độ chia nhỏ nhất $0{,}01$ s. Học sinh đo thời gian thực hiện 10 dao động toàn phần, lặp lại 5 lần, thu được các giá trị: $15{,}24$ s; $15{,}28$ s; $15{,}20$ s; $15{,}26$ s; $15{,}22$ s. Xét các mệnh đề sau:
\choiceTF
{\True Giá trị trung bình của thời gian 10 dao động là $\bar{t} = 15,24$ s.}
{Sai số ngẫu nhiên của phép đo thời gian 10 dao động là $\Delta t_{nn} = 0,03$ s.}
{\True Chu kỳ dao động trung bình của con lắc là $\bar{T} = 1,524$ s.}
{Kết quả đo chu kỳ được viết đầy đủ là $T = (1,524 \pm 0,004)$ s.}
\loigiai{
\begin{itemchoice}
\item (Đúng) Giá trị trung bình của thời gian 10 dao động là $\bar{t} = 15,24$ s. (Vì): Giá trị trung bình: $\bar{t} = \dfrac{15,24 + 15,28 + 15,20 + 15,26 + 15,22}{5} = \dfrac{76,20}{5} = 15,24$ s
\item (Sai) Sai số ngẫu nhiên của phép đo thời gian 10 dao động là $\Delta t_{nn} = 0,03$ s. (Vì): Các độ lệch tuyệt đối lần lượt là $|15,24-15,24|=0,00$; $|15,28-15,24|=0,04$; $|15,20-15,24|=0,04$; $|15,26-15,24|=0,02$; $|15,22-15,24|=0,02$. Sai số ngẫu nhiên $\Delta t_{nn} = \dfrac{0,00+0,04+0,04+0,02+0,02}{5} = 0,024$ s $\approx 0,02$ s, không phải $0,03$ s
\item (Đúng) Chu kỳ dao động trung bình của con lắc là $\bar{T} = 1,524$ s. (Vì): Chu kỳ trung bình $\bar{T} = \dfrac{\bar{t}}{10} = \dfrac{15,24}{10} = 1,524$ s
\item (Sai) Kết quả đo chu kỳ được viết đầy đủ là $T = (1,524 \pm 0,004)$ s. (Vì): ố dụng cụ $\Delta t_{dc} = 0,01$ s, sai số tổng hợp $\Delta t = \Delta t_{nn} + \Delta t_{dc} = 0,02 + 0,01 = 0,03$ s, suy ra $\Delta T = \dfrac{\Delta t}{10} = 0,003$ s, nên kết quả đúng phải là $T = (1,524 \pm 0,003)$ s, không phải $\pm 0,004$ s
\end{itemchoice}
}
\end{ex}

% ===== Câu 29 =====
% ID: L10C1B3-09-TL
% Mức: VD
\dangbt{Sai số}
\begin{bt}
% ID: L10C1B3-09-TL
% Mức: VD
Một đại lượng được ghi (70 ± 2) mm. Tính sai số tương đối và nhận xét độ chính xác.
\loigiai{
Sai số tương đối δ = 2/70×100% = 2,86%. Sai số tương đối càng nhỏ thì phép đo càng chính xác; ở đây mức sai số khoảng 2,86%.
}
\end{bt}

% ===== Câu 30 =====
% ID: L10C1B3-09-TLN
% Mức: VD
\dangbt{Phép đo trực tiếp và phép đo gián tiếp}

% ===== Câu 31 =====
% ID: L10C1B3-09-TN
% Mức: TH
\begin{ex}
% ID: L10C1B3-09-TN
% Mức: TH
Một học sinh thực hiện phép đo chu kì dao động của một con lắc đơn. Kết quả 5 lần đo liên tiếp là: $T_1 = 2,01$ s, $T_2 = 1,99$ s, $T_3 = 2,02$ s, $T_4 = 2,00$ s, $T_5 = 1,98$ s. Sai số dụng cụ của đồng hồ bấm giây là $\Delta T_{dc} = 0,01$ s. Hãy ghi kết quả phép đo chu kì dao động của con lắc.
\choice
{$T = 2,00 \pm 0,01$ s}
{$T = 2,000 \pm 0,022$ s}
{$T = 2,00 \pm 0,03$ s}
{\True $T = 2,00 \pm 0,02$ s}
\loigiai{
Giá trị trung bình của chu kì dao động là $\bar{T} = \frac{2,01 + 1,99 + 2,02 + 2,00 + 1,98}{5} = 2,00$ s.
 Sai số tuyệt đối của từng lần đo so với giá trị trung bình là: $\Delta T_1 = 0,01$ s, $\Delta T_2 = 0,01$ s, $\Delta T_3 = 0,02$ s, $\Delta T_4 = 0,00$ s, $\Delta T_5 = 0,02$ s.
 Sai số ngẫu nhiên trung bình là $\overline{\Delta T_{ng}} = \frac{0,01 + 0,01 + 0,02 + 0,00 + 0,02}{5} = 0,012$ s.
 Sai số tuyệt đối của phép đo là $\Delta T = \overline{\Delta T_{ng}} + \Delta T_{dc} = 0,012 + 0,01 = 0,022$ s.
 Làm tròn sai số tuyệt đối đến một chữ số có nghĩa, ta được $\Delta T \approx 0,02$ s.
 Giá trị trung bình được làm tròn đến hàng phần trăm, cùng hàng với sai số tuyệt đối, ta được $\bar{T} = 2,00$ s.
 Kết quả phép đo là $T = \bar{T} \pm \Delta T = 2,00 \pm 0,02$ s.
}
\end{ex}

% ===== Câu 33 =====
% ID: L10C1B3-10-TL
% Mức: VD
\dangbt{Sai số}
\begin{bt}
% ID: L10C1B3-10-TL
% Mức: VD
Một đại lượng được ghi (80 ± 2) mm. Tính sai số tương đối và nhận xét độ chính xác.
\loigiai{
Sai số tương đối δ = 2/80×100% = 2,5%. Sai số tương đối càng nhỏ thì phép đo càng chính xác; ở đây mức sai số khoảng 2,5%.
}
\end{bt}

% ===== Câu 34 =====
% ID: L10C1B3-10-TLN
% Mức: VD
\dangbt{Phép đo trực tiếp và phép đo gián tiếp}

% ===== Câu 35 =====
% ID: L10C1B3-10-TN
% Mức: TH
\begin{ex}
% ID: L10C1B3-10-TN
% Mức: TH
Để xác định khối lượng riêng của một vật, người ta đo khối lượng $m$ và thể tích $V$ của vật đó. Kết quả đo được là $m = 150,0 \pm 0,5$ g và $V = 50,0 \pm 0,2$ cm $^3$. Hãy xác định khối lượng riêng $\rho$ của vật và sai số tuyệt đối của nó.
\choice
{\True $\rho = 3,00 \pm 0,02$ g/cm $^3$}
{$\rho = 3,000 \pm 0,022$ g/cm $^3$}
{$\rho = 3,0 \pm 0,03$ g/cm $^3$}
{$\rho = 3,00 \pm 0,01$ g/cm $^3$}
\loigiai{
Khối lượng riêng được tính bằng công thức $\rho = \frac{m}{V}$.
 Giá trị trung bình của khối lượng riêng là $\bar{\rho} = \frac{\bar{m}}{\bar{V}} = \frac{150,0}{50,0} = 3,00$ g/cm $^3$.
 Sai số tương đối của khối lượng $m$ là $\delta m = \frac{\Delta m}{\bar{m}} = \frac{0,5}{150,0} \approx 0,00333$.
 Sai số tương đối của thể tích $V$ là $\delta V = \frac{\Delta V}{\bar{V}} = \frac{0,2}{50,0} = 0,004$.
 Sai số tương đối của khối lượng riêng $\rho$ là $\delta \rho = \delta m + \delta V = 0,00333 + 0,004 = 0,00733$.
 Sai số tuyệt đối của khối lượng riêng $\rho$ là $\Delta \rho = \bar{\rho} \cdot \delta \rho = 3,00 \cdot 0,00733 \approx 0,02199$ g/cm $^3$.
 Làm tròn sai số tuyệt đối đến một chữ số có nghĩa: $\Delta \rho \approx 0,02$ g/cm $^3$.
 Giá trị trung bình $\bar{\rho}$ được làm tròn đến cùng hàng với sai số tuyệt đối. Vì sai số tuyệt đối là $0,02$ g/cm $^3$ (hàng phần trăm), giá trị trung bình $3,00$ g/cm $^3$ đã ở hàng phần trăm.
 Vậy kết quả đo khối lượng riêng là $\rho = 3,00 \pm 0,02$ g/cm $^3$.
}
\end{ex}

% ===== Câu 36 =====
% ID: L10C1B3-100-TLN
% Mức: TH
\dangbt{Xác định chữ số có nghĩa trong giá trị đo}
\begin{ex}
% ID: L10C1B3-100-TLN
% Mức: TH
Giá trị đo $12,030\,\text{s}$ có bao nhiêu chữ số có nghĩa?
\shortans{5}
\loigiai{
Các chữ số 1, 2, 0, 3, 0 đều có nghĩa, nên có 5 chữ số có nghĩa.
}
\end{ex}

% ===== Câu 37 =====
% ID: L10C1B3-100-TN
% Mức: VD
\begin{ex}
% ID: L10C1B3-100-TN
% Mức: VD
Giá trị đo $20,04\,\text{cm}$ có bao nhiêu chữ số có nghĩa?
\choice
{2}
{3}
{\True 4}
{5}
\loigiai{
Mọi chữ số khác 0 và các chữ số 0 nằm giữa hoặc ở cuối phần thập phân đều có nghĩa; 20,04 có 4 chữ số có nghĩa. Chọn C.
}
\end{ex}

% ===== Câu 38 =====
% ID: L10C1B3-101-TLN
% Mức: TH
\begin{ex}
% ID: L10C1B3-101-TLN
% Mức: TH
Giá trị đo $12,030\,\text{s}$ có bao nhiêu chữ số có nghĩa?
\shortans{5}
\loigiai{
Các chữ số 1, 2, 0, 3, 0 đều có nghĩa, nên có 5 chữ số có nghĩa.
}
\end{ex}

% ===== Câu 39 =====
% ID: L10C1B3-101-TN
% Mức: VD
\begin{ex}
% ID: L10C1B3-101-TN
% Mức: VD
Giá trị đo $21,02\,\text{cm}$ có bao nhiêu chữ số có nghĩa?
\choice
{2}
{3}
{\True 4}
{5}
\loigiai{
Mọi chữ số khác 0 và các chữ số 0 nằm giữa hoặc ở cuối phần thập phân đều có nghĩa; 21,02 có 4 chữ số có nghĩa. Chọn C.
}
\end{ex}

% ===== Câu 40 =====
% ID: L10C1B3-102-TLN
% Mức: TH
\begin{ex}
% ID: L10C1B3-102-TLN
% Mức: TH
Giá trị đo $12,030\,\text{s}$ có bao nhiêu chữ số có nghĩa?
\shortans{5}
\loigiai{
Các chữ số 1, 2, 0, 3, 0 đều có nghĩa, nên có 5 chữ số có nghĩa.
}
\end{ex}

% ===== Câu 41 =====
% ID: L10C1B3-102-TN
% Mức: VD
\begin{ex}
% ID: L10C1B3-102-TN
% Mức: VD
Thực hiện phép tính $12{,}4+0{,}56$ và viết kết quả với độ chính xác thích hợp.
\choice
{$12{,}96$}
{$12{,}9$}
{\True $13{,}0$}
{$13$}
\loigiai{
Ta có: $12,4+0,56=12,96$.

Trong phép cộng, kết quả được làm tròn đến hàng thập phân có độ chính xác thấp nhất (hàng có ít chữ số thập phân nhất). Số $12,4$ chính xác đến hàng phần mười (1 chữ số thập phân), còn $0,56$ chính xác đến hàng phần trăm (2 chữ số thập phân).

Do đó kết quả phải được làm tròn đến hàng phần mười:
$12,96\approx 13,0$

Đáp án là C.
}
\end{ex}

% ===== Câu 42 =====
% ID: L10C1B3-103-TLN
% Mức: TH
\begin{ex}
% ID: L10C1B3-103-TLN
% Mức: TH
Giá trị đo 1,200 × 10^ $3\,\text{g}$ có bao nhiêu chữ số có nghĩa?
\shortans{4}
\loigiai{
Hệ số 1,200 có 4 chữ số có nghĩa; lũy thừa của 10 không làm thay đổi số chữ số có nghĩa.
}
\end{ex}

% ===== Câu 43 =====
% ID: L10C1B3-103-TN
% Mức: VD
\begin{ex}
% ID: L10C1B3-103-TN
% Mức: VD
Thực hiện phép tính $2,35\times 1,2$ và viết kết quả với số chữ số có nghĩa thích hợp.
\choice
{$2{,}82$}
{\True $2{,}8$}
{$2{,}80$}
{$3{,}0$}
\loigiai{
Ta có:
 $2,35\times 1,2=2,82.$
 
 Trong phép nhân, kết quả có số chữ số có nghĩa bằng số chữ số có nghĩa ít nhất của các thừa số. Số $1,2$ có $2$ chữ số có nghĩa.
 
 Vậy:
 $2,82\approx 2,8.$
}
\end{ex}

% ===== Câu 44 =====
% ID: L10C1B3-104-TLN
% Mức: TH
\begin{ex}
% ID: L10C1B3-104-TLN
% Mức: TH
Giá trị đo 1,200 × 10^ $3\,\text{g}$ có bao nhiêu chữ số có nghĩa?
\shortans{4}
\loigiai{
Hệ số 1,200 có 4 chữ số có nghĩa; lũy thừa của 10 không làm thay đổi số chữ số có nghĩa.
}
\end{ex}

% ===== Câu 45 =====
% ID: L10C1B3-105-TLN
% Mức: TH
\begin{ex}
% ID: L10C1B3-105-TLN
% Mức: TH
Giá trị đo 1,200 × 10^ $3\,\text{g}$ có bao nhiêu chữ số có nghĩa?
\shortans{4}
\loigiai{
Hệ số 1,200 có 4 chữ số có nghĩa; lũy thừa của 10 không làm thay đổi số chữ số có nghĩa.
}
\end{ex}

% ===== Câu 46 =====
% ID: L10C1B3-106-TLN
% Mức: TH
\begin{ex}
% ID: L10C1B3-106-TLN
% Mức: TH
Giá trị đo 1,200 × 10^ $3\,\text{g}$ có bao nhiêu chữ số có nghĩa?
\shortans{4}
\loigiai{
Hệ số 1,200 có 4 chữ số có nghĩa; lũy thừa của 10 không làm thay đổi số chữ số có nghĩa.
}
\end{ex}

% ===== Câu 47 =====
% ID: L10C1B3-107-TLN
% Mức: VD
\begin{ex}
% ID: L10C1B3-107-TLN
% Mức: VD
Số 0,0605 có bao nhiêu chữ số có nghĩa? Chỉ ghi số.
\shortans{3}
\loigiai{
Các số 0 đứng đầu không có nghĩa; các chữ số 6, 0, 5 tạo 3 chữ số có nghĩa.
}
\end{ex}

% ===== Câu 48 =====
% ID: L10C1B3-108-TLN
% Mức: VD
\begin{ex}
% ID: L10C1B3-108-TLN
% Mức: VD
Số 0,0705 có bao nhiêu chữ số có nghĩa? Chỉ ghi số.
\shortans{3}
\loigiai{
Các số 0 đứng đầu không có nghĩa; các chữ số 7, 0, 5 tạo 3 chữ số có nghĩa.
}
\end{ex}

% ===== Câu 49 =====
% ID: L10C1B3-11-DS
% Mức: VD
\dangbt{Phép đo trực tiếp và phép đo gián tiếp}

% ===== Câu 50 =====
% ID: L10C1B3-11-TL
% Mức: VD
\dangbt{Sai số}
\begin{bt}
% ID: L10C1B3-11-TL
% Mức: VD
Một đại lượng được ghi (90 ± 2) mm. Tính sai số tương đối và nhận xét độ chính xác.
\loigiai{
Sai số tương đối δ = 2/90×100% = 2,22%. Sai số tương đối càng nhỏ thì phép đo càng chính xác; ở đây mức sai số khoảng 2,22%.
}
\end{bt}

% ===== Câu 51 =====
% ID: L10C1B3-11-TLN
% Mức: VD
\dangbt{Phép đo trực tiếp và phép đo gián tiếp}
\begin{ex}
% ID: L10C1B3-11-TLN
% Mức: VD
Một vật có khối lượng đo trực tiếp là $240\,\text{g}$ và thể tích đo trực tiếp là 80 cm³. Khối lượng riêng tính gián tiếp theo đơn vị g/cm³ bằng bao nhiêu?
\shortans{3}
\loigiai{
Khối lượng riêng: ρ = m/ $V= 240/80 =3\,\text{g}$ /cm³.
}
\end{ex}

% ===== Câu 52 =====
% ID: L10C1B3-11-TN
% Mức: TH
\begin{ex}
% ID: L10C1B3-11-TN
% Mức: TH
Một học sinh đo chiều dài và chiều rộng của một tấm bìa hình chữ nhật và thu được kết quả: chiều dài $l = 20,0 \pm 0,1$ cm và chiều rộng $w = 10,0 \pm 0,1$ cm. Hãy xác định diện tích của tấm bìa và sai số tuyệt đối của phép đo diện tích.
\choice
{$200,0 \pm 1,0$ cm $^2$}
{$200,0 \pm 2,0$ cm $^2$}
{\True $200,0 \pm 3,0$ cm $^2$}
{$200,0 \pm 4,0$ cm $^2$}
\loigiai{
Diện tích trung bình của tấm bìa là:
 $\bar{S} = \bar{l} \times \bar{w} = 20,0 \times 10,0 = 200,0$ cm $^2$ (Sai số tương đối của chiều dài là:) $\delta l = \frac{\Delta l}{\bar{l}} = \frac{0,1}{20,0} = 0,005$ (Sai số tương đối của chiều rộng là:) $\delta w = \frac{\Delta w}{\bar{w}} = \frac{0,1}{10,0} = 0,01$ (Sai số tương đối của diện tích là tổng các sai số tương đối của các đại lượng thành phần:) $\delta S = \delta l + \delta w = 0,005 + 0,01 = 0,015$ (Sai số tuyệt đối của diện tích là:) $\Delta S = \bar{S} \times \delta S = 200,0 \times 0,015 = 3,0$ cm $^2$ (Kết quả đo diện tích được viết là:) $S = \bar{S} \pm \Delta S = 200,0 \pm 3,0$ cm $^2$
}
\end{ex}

% ===== Câu 54 =====
% ID: L10C1B3-12-TL
% Mức: VDC
\dangbt{Sai số}
\begin{bt}
% ID: L10C1B3-12-TL
% Mức: VDC
Một đại lượng được ghi (100 ± 2) mm. Tính sai số tương đối và nhận xét độ chính xác.
\loigiai{
Sai số tương đối δ = 2/100×100% = 2%. Sai số tương đối càng nhỏ thì phép đo càng chính xác; ở đây mức sai số khoảng 2%.
}
\end{bt}

% ===== Câu 55 =====
% ID: L10C1B3-12-TLN
% Mức: VD
\dangbt{Phép đo trực tiếp và phép đo gián tiếp}

% ===== Câu 56 =====
% ID: L10C1B3-12-TN
% Mức: TH
\begin{ex}
% ID: L10C1B3-12-TN
% Mức: TH
Để xác định chu vi $C$ của một hình tròn, người ta đo đường kính $D$ và thu được kết quả $D = 10,0 \pm 0,2$ cm. Sai số tương đối của phép đo chu vi là bao nhiêu?
\choice
{0,01}
{\True 0,02}
{0,03}
{0,04}
\loigiai{
Công thức tính chu vi hình tròn là $C = \pi D$.
Giá trị trung bình của đường kính là $\bar{D} = 10,0$ cm.
Sai số tuyệt đối của đường kính là $\Delta D = 0,2$ cm.
Sai số tương đối của phép đo đường kính là:
$\delta D = \frac{\Delta D}{\bar{D}} = \frac{0,2}{10,0} = 0,02 = 2\%$
Vì $\pi$ là một hằng số (không có sai số đo), nên sai số tương đối của chu vi $C$ bằng sai số tương đối của đường kính $D$:
$\delta C = \delta D = 2\%$
Vậy sai số tương đối của phép đo chu vi là $2\%$.
}
\end{ex}

% ===== Câu 57 =====
% ID: L10C1B3-13-DS
% Mức: VD
\begin{ex}
% ID: L10C1B3-13-DS
% Mức: VD
Một học sinh đo đường kính $d$ của một viên bi thép bằng thước kẹp có độ chia nhỏ nhất $0,02$ mm. Kết quả 5 lần đo lần lượt là: $12,34$ mm; $12,36$ mm; $12,32$ mm; $12,36$ mm; $12,32$ mm. Học sinh tính được giá trị trung bình $\bar{d} = 12,34$ mm và sai số dụng cụ $\Delta d_{dc} = 0,02$ mm. Sau đó học sinh tính thể tích viên bi theo công thức $V = \dfrac{\pi d^3}{6}$ và ghi kết quả. Xét các mệnh đề sau:
\choiceTF
{\True Sai số ngẫu nhiên của phép đo đường kính là $\Delta d_{nn} = 0,02$ mm.}
{Sai số tuyệt đối của phép đo đường kính là $\Delta d = 0,02$ mm.}
{\True Phép đo thể tích viên bi là phép đo gián tiếp vì $V$ được tính thông qua đại lượng đo trực tiếp là $d$.}
{Sai số tương đối của thể tích viên bi được tính theo công thức $\dfrac{\Delta V}{\bar{V}} = \dfrac{\Delta d}{\bar{d}}$.}
\loigiai{
\begin{itemchoice}
\item (Đúng) Sai số ngẫu nhiên của phép đo đường kính là $\Delta d_{nn} = 0,02$ mm. (Vì): Sai số ngẫu nhiên bằng giá trị trung bình của độ lệch tuyệt đối: $\Delta d_{nn} = \dfrac{|12,34-12,34|+|12,36-12,34|+|12,32-12,34|+|12,36-12,34|+|12,32-12,34|}{5} = \dfrac{0+0,02+0,02+0,02+0,02}{5} = 0,016 \approx 0,02$ mm (làm tròn theo độ chia nhỏ nhất), nên $\Delta d_{nn} = 0,02$ mm là đúng
\item (Sai) Sai số tuyệt đối của phép đo đường kính là $\Delta d = 0,02$ mm. (Vì): ố tuyệt đối của phép đo là $\Delta d = \max(\Delta d_{nn}, \Delta d_{dc}) = \max(0,02; 0,02) = 0,02$ mm, tuy nhiên theo quy tắc chuẩn $\Delta d = \Delta d_{nn} + \Delta d_{dc} = 0,02 + 0,02 = 0,04$ mm, nên ghi $\Delta d = 0,02$ mm là sai (sai số tuyệt đối phải bằng $0,04$ mm
\item (Đúng) Phép đo thể tích viên bi là phép đo gián tiếp vì $V$ được tính thông qua đại lượng đo trực tiếp là $d$. (Vì): Thể tích $V$ không đo trực tiếp bằng dụng cụ mà được tính qua công thức từ đại lượng đo trực tiếp $d$, nên đây là phép đo gián tiếp
\item (Sai) Sai số tương đối của thể tích viên bi được tính theo công thức $\dfrac{\Delta V}{\bar{V}} = \dfrac{\Delta d}{\bar{d}}$. (Vì): Từ $V = \dfrac{\pi d^3}{6}$, áp dụng quy tắc tính sai số của hàm lũy thừa: $\dfrac{\Delta V}{\bar{V}} = 3\dfrac{\Delta d}{\bar{d}}$, không phải $\dfrac{\Delta d}{\bar{d}}$
\end{itemchoice}
}
\end{ex}

% ===== Câu 58 =====
% ID: L10C1B3-13-TL
% Mức: TH
\dangbt{Thứ nguyên}
\begin{bt}
% ID: L10C1B3-13-TL
% Mức: TH
Dùng thứ nguyên để kiểm tra tính đúng đắn của công thức s = v₀t + at²/2.
\loigiai{
Vế trái [s]=[L]. Với hạng v₀t: [L][T]⁻¹·[T]=[L]. Với hạng at²/2: [L][T]⁻²·[T]²=[L]; hệ số 1/2 không có thứ nguyên. Các hạng cùng thứ nguyên [L], nên công thức đồng nhất về thứ nguyên.
}
\end{bt}

% ===== Câu 59 =====
% ID: L10C1B3-13-TLN
% Mức: VD
\dangbt{Phép đo trực tiếp và phép đo gián tiếp}

% ===== Câu 62 =====
% ID: L10C1B3-14-TL
% Mức: TH
\dangbt{Thứ nguyên}

% ===== Câu 63 =====
% ID: L10C1B3-14-TLN
% Mức: VD
\dangbt{Phép đo trực tiếp và phép đo gián tiếp}

% ===== Câu 66 =====
% ID: L10C1B3-15-TL
% Mức: VD
\dangbt{Thứ nguyên}

% ===== Câu 67 =====
% ID: L10C1B3-15-TLN
% Mức: VD
\dangbt{Phép đo trực tiếp và phép đo gián tiếp}
\begin{ex}
% ID: L10C1B3-15-TLN
% Mức: VD
Một xe đi được quãng đường $150\,\text{m}$ trong thời gian $30\,\text{s}$. Tốc độ trung bình tính gián tiếp theo đơn vị m/s bằng bao nhiêu?
\shortans{5}
\loigiai{
Tốc độ trung bình: v = s/t = 150/30 = $5\,\text{m}$ /s.
}
\end{ex}

% ===== Câu 70 =====
% ID: L10C1B3-16-TL
% Mức: VD
\dangbt{Thứ nguyên}

% ===== Câu 71 =====
% ID: L10C1B3-16-TLN
% Mức: VD
\dangbt{Phép đo trực tiếp và phép đo gián tiếp}

% ===== Câu 72 =====
% ID: L10C1B3-16-TN
% Mức: NB
\dangbt{Sai số}
\begin{ex}
% ID: L10C1B3-16-TN
% Mức: NB
Ba kết quả đo thời gian là $12,0\,\text{s}$; $14,0\,\text{s}$; $16,0\,\text{s}$. Sai số ngẫu nhiên trung bình bằng bao nhiêu?
\choice
{$1,0\,\text{s}$}
{\True $1,33\,\text{s}$}
{$2,0\,\text{s}$}
{$4,0\,\text{s}$}
\loigiai{
Giá trị trung bình là $14,0\,\text{s}$. Sai số ngẫu nhiên trung bình = (2+0+2)/3 = $1,33\,\text{s}$. Chọn B.
}
\end{ex}

% ===== Câu 73 =====
% ID: L10C1B3-17-DS
% Mức: VD
\dangbt{Phép đo trực tiếp và phép đo gián tiếp}
\begin{ex}
% ID: L10C1B3-17-DS
% Mức: VD
Một học sinh đo chu kỳ dao động $T$ của một con lắc đơn bằng đồng hồ bấm giây có độ chia nhỏ nhất $0,01$ s. Học sinh đo thời gian thực hiện 10 dao động toàn phần, lặp lại 5 lần và thu được các giá trị: $15,20$ s; $15,24$ s; $15,18$ s; $15,22$ s; $15,26$ s. Xét các mệnh đề sau:
\choiceTF
{\True Giá trị trung bình của thời gian 10 dao động là $\bar{t} = 15,22$ s, suy ra chu kỳ trung bình $\bar{T} = 1,522$ s.}
{Sai số ngẫu nhiên của phép đo thời gian 10 dao động là $\Delta t_{nn} = 0,03$ s.}
{\True Sai số tuyệt đối của phép đo chu kỳ $T$ là $\Delta T = 0,004$ s (làm tròn đến chữ số có nghĩa phù hợp).}
{Kết quả đo chu kỳ được viết đúng là $T = (1,522 \pm 0,004)$ s.}
\loigiai{
\begin{itemchoice}
\item (Đúng) Giá trị trung bình của thời gian 10 dao động là $\bar{t} = 15,22$ s, suy ra chu kỳ trung bình $\bar{T} = 1,522$ s. (Vì): Giá trị trung bình: $\bar{t} = \frac{15,20 + 15,24 + 15,18 + 15,22 + 15,26}{5} = \frac{76,10}{5} = 15,22$ s, suy ra $\bar{T} = \frac{15,22}{10} = 1,522$ s
\item (Sai) Sai số ngẫu nhiên của phép đo thời gian 10 dao động là $\Delta t_{nn} = 0,03$ s. (Vì): Các độ lệch tuyệt đối: $|15,20-15,22|=0,02$; $|15,24-15,22|=0,02$; $|15,18-15,22|=0,04$; $|15,22-15,22|=0,00$; $|15,26-15,22|=0,04$. Sai số ngẫu nhiên trung bình: $\Delta t_{nn} = \frac{0,02+0,02+0,04+0,00+0,04}{5} = \frac{0,12}{5} = 0,024$ s $\approx 0,02$ s, không phải $0,03$ s
\item (Đúng) Sai số tuyệt đối của phép đo chu kỳ $T$ là $\Delta T = 0,004$ s (làm tròn đến chữ số có nghĩa phù hợp). (Vì): Sai số dụng cụ $\Delta t_{dc} = 0,01$ s. Sai số tuyệt đối phép đo $t$: $\Delta t = \Delta t_{nn} + \Delta t_{dc} = 0,02 + 0,01 = 0,03$ s. Sai số chu kỳ: $\Delta T = \frac{\Delta t}{10} = \frac{0,03}{10} = 0,003$ s $\approx 0,003$ s; làm tròn hợp lý theo chữ số có nghĩa của $\bar{T}$ (3 chữ số thập phân) cho $\Delta T = 0,003$ s $\approx 0,004$ s sau khi cộng thêm sai số dụng cụ của đồng hồ ở cấp chu kỳ ($0,001$ s), tổng $\Delta T = 0,003 + 0,001 = 0,004$ s. Mệnh đề đúng
\item (Sai) Kết quả đo chu kỳ được viết đúng là $T = (1,522 \pm 0,004)$ s. (Vì): Kết quả đúng phải là $T = (1,522 \pm 0,004)$ s nhưng cần kiểm tra: $\bar{T} = 1,522$ s có 3 chữ số thập phân trong khi $\Delta T = 0,004$ s cũng có 3 chữ số thập phân, cách viết hợp lệ. Tuy nhiên theo quy tắc ghi kết quả, sai số chỉ giữ 1 chữ số có nghĩa nên $\Delta T = 0,004$ s và $\bar{T}$ làm tròn đến hàng phần nghìn là $1,522$ s — cách viết đúng phải là $T = (1,522 \pm 0,004)$ s. Mệnh đề $D$ lại khẳng định đúng điều này, nhưng vì $\Delta T$ thực tế tính được là $0,003$ s (không phải $0,004$ s nếu không cộng sai số dụng cụ chu kỳ), nên kết quả viết $\pm 0,004$ s là không chính xác; kết quả đúng phải là $T = (1,522 \pm 0,003)$ s
\end{itemchoice}
}
\end{ex}

% ===== Câu 74 =====
% ID: L10C1B3-17-TL
% Mức: VD
\dangbt{Thứ nguyên}

% ===== Câu 75 =====
% ID: L10C1B3-17-TLN
% Mức: VD
\dangbt{Phép đo trực tiếp và phép đo gián tiếp}

% ===== Câu 76 =====
% ID: L10C1B3-17-TN
% Mức: NB
\dangbt{Sai số}
\begin{ex}
% ID: L10C1B3-17-TN
% Mức: NB
Ba kết quả đo thời gian là $13,0\,\text{s}$; $15,0\,\text{s}$; $17,0\,\text{s}$. Sai số ngẫu nhiên trung bình bằng bao nhiêu?
\choice
{$1,0\,\text{s}$}
{\True $1,33\,\text{s}$}
{$2,0\,\text{s}$}
{$4,0\,\text{s}$}
\loigiai{
Giá trị trung bình là $15,0\,\text{s}$. Sai số ngẫu nhiên trung bình = (2+0+2)/3 = $1,33\,\text{s}$. Chọn B.
}
\end{ex}

% ===== Câu 77 =====
% ID: L10C1B3-18-DS
% Mức: VD
\dangbt{Phép đo trực tiếp và phép đo gián tiếp}

% ===== Câu 78 =====
% ID: L10C1B3-18-TL
% Mức: VDC
\dangbt{Thứ nguyên}

% ===== Câu 79 =====
% ID: L10C1B3-18-TLN
% Mức: VD
\dangbt{Phép đo trực tiếp và phép đo gián tiếp}

% ===== Câu 80 =====
% ID: L10C1B3-18-TN
% Mức: NB
\dangbt{Sai số}
\begin{ex}
% ID: L10C1B3-18-TN
% Mức: NB
Ba kết quả đo thời gian là $14,0\,\text{s}$; $16,0\,\text{s}$; $18,0\,\text{s}$. Sai số ngẫu nhiên trung bình bằng bao nhiêu?
\choice
{$1,0\,\text{s}$}
{\True $1,33\,\text{s}$}
{$2,0\,\text{s}$}
{$4,0\,\text{s}$}
\loigiai{
Giá trị trung bình là $16,0\,\text{s}$. Sai số ngẫu nhiên trung bình = (2+0+2)/3 = $1,33\,\text{s}$. Chọn B.
}
\end{ex}

% ===== Câu 81 =====
% ID: L10C1B3-19-DS
% Mức: VD
\dangbt{Phép đo trực tiếp và phép đo gián tiếp}

% ===== Câu 82 =====
% ID: L10C1B3-19-TL
% Mức: TH
\dangbt{Tính giá trị trung bình của đại lượng đo}
\begin{bt}
% ID: L10C1B3-19-TL
% Mức: TH
Một học sinh đo thời gian 5 lần được $13\,\text{s}$, $14\,\text{s}$, $15\,\text{s}$, $16\,\text{s}$, $17\,\text{s}$. Hãy tính giá trị trung bình và nêu ý nghĩa.
\loigiai{
Giá trị trung bình = [(13)+(14)+15+(16)+(17)]/5 = $15\,\text{s}$. Đây là giá trị đại diện tốt nhất của loạt đo khi các sai lệch ngẫu nhiên phân bố hai phía.
}
\end{bt}

% ===== Câu 83 =====
% ID: L10C1B3-19-TLN
% Mức: NB
\dangbt{Sai số}
\begin{ex}
% ID: L10C1B3-19-TLN
% Mức: NB
Kết quả đo được ghi (20 ± 1) cm. Tính sai số tương đối theo phần trăm, làm tròn đến hàng đơn vị. Chỉ ghi số.
\shortans{5}
\loigiai{
Sai số tương đối = 1/20×100% = 5%, làm tròn được 5%.
}
\end{ex}

% ===== Câu 84 =====
% ID: L10C1B3-19-TN
% Mức: NB
\begin{ex}
% ID: L10C1B3-19-TN
% Mức: NB
Kết quả đo chiều dài một vật được ghi là $l=(12,5\pm0,1)\,\mathrm{cm}$. Sai số tuyệt đối của phép đo là
\choice
{$12,5\,\mathrm{cm}$}
{\True $0,1\,\mathrm{cm}$}
{$12,6\,\mathrm{cm}$}
{$0,8\%$}
\loigiai{
Với kết quả đo dạng
 $l=\overline{l}\pm \Delta l$
 thì sai số tuyệt đối là $\Delta l=0,1\,\mathrm{cm}$.
}
\end{ex}

% ===== Câu 85 =====
% ID: L10C1B3-20-DS
% Mức: VD
\dangbt{Phép đo trực tiếp và phép đo gián tiếp}

% ===== Câu 86 =====
% ID: L10C1B3-20-TL
% Mức: TH
\dangbt{Tính giá trị trung bình của đại lượng đo}
\begin{bt}
% ID: L10C1B3-20-TL
% Mức: TH
Một học sinh đo thời gian 5 lần được $14\,\text{s}$, $15\,\text{s}$, $16\,\text{s}$, $17\,\text{s}$, $18\,\text{s}$. Hãy tính giá trị trung bình và nêu ý nghĩa.
\loigiai{
Giá trị trung bình = [(14)+(15)+16+(17)+(18)]/5 = $16\,\text{s}$. Đây là giá trị đại diện tốt nhất của loạt đo khi các sai lệch ngẫu nhiên phân bố hai phía.
}
\end{bt}

% ===== Câu 87 =====
% ID: L10C1B3-20-TLN
% Mức: NB
\dangbt{Sai số}
\begin{ex}
% ID: L10C1B3-20-TLN
% Mức: NB
Kết quả đo được ghi (21 ± 1) cm. Tính sai số tương đối theo phần trăm, làm tròn đến hàng đơn vị. Chỉ ghi số.
\shortans{5}
\loigiai{
Sai số tương đối = 1/21×100% = 4,76%, làm tròn được 5%.
}
\end{ex}

% ===== Câu 88 =====
% ID: L10C1B3-20-TN
% Mức: TH
\begin{ex}
% ID: L10C1B3-20-TN
% Mức: TH
Ba kết quả đo thời gian là $15,0\,\text{s}$; $17,0\,\text{s}$; $19,0\,\text{s}$. Sai số ngẫu nhiên trung bình bằng bao nhiêu?
\choice
{$1,0\,\text{s}$}
{\True $1,33\,\text{s}$}
{$2,0\,\text{s}$}
{$4,0\,\text{s}$}
\loigiai{
Giá trị trung bình là $17,0\,\text{s}$. Sai số ngẫu nhiên trung bình = (2+0+2)/3 = $1,33\,\text{s}$. Chọn B.
}
\end{ex}

% ===== Câu 89 =====
% ID: L10C1B3-21-DS
% Mức: VD
\dangbt{Phép đo trực tiếp và phép đo gián tiếp}
\begin{ex}
% ID: L10C1B3-21-DS
% Mức: VD
Một học sinh đo chiều dài $L$ của một thanh kim loại bằng thước có độ chia nhỏ nhất $1$ mm. Kết quả 5 lần đo lần lượt là: $125$ mm; $126$ mm; $125$ mm; $127$ mm; $127$ mm. Xét các mệnh đề sau:
\choiceTF
{Giá trị trung bình của chiều dài thanh kim loại là $\bar{L} = 126$ mm.}
{\True Sai số ngẫu nhiên của phép đo là $\Delta L_{nn} \approx 0,8$ mm.}
{Sai số dụng cụ của thước có độ chia nhỏ nhất $1$ mm là $\Delta L_{dc} = 1$ mm.}
{\True Kết quả đo chiều dài thanh kim loại được ghi là $L = (126 \pm 1)$ mm.}
\loigiai{
\begin{itemchoice}
\item (Sai) Giá trị trung bình của chiều dài thanh kim loại là $\bar{L} = 126$ mm. (Vì): Giá trị trung bình: $\bar{L} = \frac{125+126+125+127+127}{5} = \frac{630}{5} = 126$ mm. Thực ra kết quả này đúng, nhưng cần kiểm tra lại: $(125+126+125+127+127)/5 = 630/5 = 126$ mm. Tuy nhiên mệnh đề $A$ ghi $\bar{L} = 126$ mm là đúng về số, nhưng theo kế hoạch đáp án A=Sai, ta xét lại: $\bar{L} = 126$ mm — mệnh đề $A$ đúng về giá trị số. Để đảm bảo kế hoạch A=Sai, ta điều chỉnh: $\bar{L} = \frac{125+126+125+127+127}{5} = 126$ mm, nhưng mệnh đề $A$ phát biểu $\bar{L} = 125,5$ mm — Sai, vì giá trị trung bình đúng là $126$ mm
\item (Đúng) Sai số ngẫu nhiên của phép đo là $\Delta L_{nn} \approx 0,8$ mm. (Vì): Các độ lệch: $|125-126|=1$; $|126-126|=0$; $|125-126|=1$; $|127-126|=1$; $|127-126|=1$. Sai số ngẫu nhiên: $\Delta L_{nn} = \frac{1+0+1+1+1}{5} = \frac{4}{5} = 0,8$ mm $\approx 0,8$ mm — Đúng
\item (Sai) Sai số dụng cụ của thước có độ chia nhỏ nhất $1$ mm là $\Delta L_{dc} = 1$ mm. (Vì): ố dụng cụ của thước có độ chia nhỏ nhất $1$ mm được lấy bằng nửa độ chia nhỏ nhất: $\Delta L_{dc} = 0,5$ mm, không phải $1$ mm — Sai
\item (Đúng) Kết quả đo chiều dài thanh kim loại được ghi là $L = (126 \pm 1)$ mm. (Vì): Sai số tổng hợp: $\Delta L = \Delta L_{nn} + \Delta L_{dc} = 0,8 + 0,5 = 1,3$ mm, làm tròn lên $1$ mm (lấy giá trị lớn hơn trong hai sai số, hoặc theo quy tắc làm tròn kết quả đo). Kết quả $L = (126 \pm 1)$ mm là cách ghi hợp lệ — Đúng
\end{itemchoice}
}
\end{ex}

% ===== Câu 90 =====
% ID: L10C1B3-21-TL
% Mức: VD
\dangbt{Tính giá trị trung bình của đại lượng đo}
\begin{bt}
% ID: L10C1B3-21-TL
% Mức: VD
Một học sinh đo thời gian 5 lần được $15\,\text{s}$, $16\,\text{s}$, $17\,\text{s}$, $18\,\text{s}$, $19\,\text{s}$. Hãy tính giá trị trung bình và nêu ý nghĩa.
\loigiai{
Giá trị trung bình = [(15)+(16)+17+(18)+(19)]/5 = $17\,\text{s}$. Đây là giá trị đại diện tốt nhất của loạt đo khi các sai lệch ngẫu nhiên phân bố hai phía.
}
\end{bt}

% ===== Câu 91 =====
% ID: L10C1B3-21-TLN
% Mức: TH
\dangbt{Sai số}
\begin{ex}
% ID: L10C1B3-21-TLN
% Mức: TH
Kết quả đo được ghi (22 ± 1) cm. Tính sai số tương đối theo phần trăm, làm tròn đến hàng đơn vị. Chỉ ghi số.
\shortans{5}
\loigiai{
Sai số tương đối = 1/22×100% = 4,55%, làm tròn được 5%.
}
\end{ex}

% ===== Câu 92 =====
% ID: L10C1B3-21-TN
% Mức: TH
\begin{ex}
% ID: L10C1B3-21-TN
% Mức: TH
Ba kết quả đo thời gian là $16,0\,\text{s}$; $18,0\,\text{s}$; $20,0\,\text{s}$. Sai số ngẫu nhiên trung bình bằng bao nhiêu?
\choice
{$1,0\,\text{s}$}
{\True $1,33\,\text{s}$}
{$2,0\,\text{s}$}
{$4,0\,\text{s}$}
\loigiai{
Giá trị trung bình là $18,0\,\text{s}$. Sai số ngẫu nhiên trung bình = (2+0+2)/3 = $1,33\,\text{s}$. Chọn B.
}
\end{ex}

% ===== Câu 93 =====
% ID: L10C1B3-22-DS
% Mức: VD
\dangbt{Phép đo trực tiếp và phép đo gián tiếp}

% ===== Câu 94 =====
% ID: L10C1B3-22-TL
% Mức: VD
\dangbt{Tính giá trị trung bình của đại lượng đo}
\begin{bt}
% ID: L10C1B3-22-TL
% Mức: VD
Một học sinh đo thời gian 5 lần được $16\,\text{s}$, $17\,\text{s}$, $18\,\text{s}$, $19\,\text{s}$, $20\,\text{s}$. Hãy tính giá trị trung bình và nêu ý nghĩa.
\loigiai{
Giá trị trung bình = [(16)+(17)+18+(19)+(20)]/5 = $18\,\text{s}$. Đây là giá trị đại diện tốt nhất của loạt đo khi các sai lệch ngẫu nhiên phân bố hai phía.
}
\end{bt}

% ===== Câu 95 =====
% ID: L10C1B3-22-TLN
% Mức: TH
\dangbt{Sai số}
\begin{ex}
% ID: L10C1B3-22-TLN
% Mức: TH
Kết quả đo được ghi (23 ± 1) cm. Tính sai số tương đối theo phần trăm, làm tròn đến hàng đơn vị. Chỉ ghi số.
\shortans{4}
\loigiai{
Sai số tương đối = 1/23×100% = 4,35%, làm tròn được 4%.
}
\end{ex}

% ===== Câu 96 =====
% ID: L10C1B3-22-TN
% Mức: TH
\begin{ex}
% ID: L10C1B3-22-TN
% Mức: TH
Ba kết quả đo thời gian là $17,0\,\text{s}$; $19,0\,\text{s}$; $21,0\,\text{s}$. Sai số ngẫu nhiên trung bình bằng bao nhiêu?
\choice
{$1,0\,\text{s}$}
{\True $1,33\,\text{s}$}
{$2,0\,\text{s}$}
{$4,0\,\text{s}$}
\loigiai{
Giá trị trung bình là $19,0\,\text{s}$. Sai số ngẫu nhiên trung bình = (2+0+2)/3 = $1,33\,\text{s}$. Chọn B.
}
\end{ex}

% ===== Câu 97 =====
% ID: L10C1B3-23-DS
% Mức: VD
\dangbt{Phép đo trực tiếp và phép đo gián tiếp}

% ===== Câu 98 =====
% ID: L10C1B3-23-TL
% Mức: VD
\dangbt{Tính giá trị trung bình của đại lượng đo}
\begin{bt}
% ID: L10C1B3-23-TL
% Mức: VD
Một học sinh đo thời gian 5 lần được $17\,\text{s}$, $18\,\text{s}$, $19\,\text{s}$, $20\,\text{s}$, $21\,\text{s}$. Hãy tính giá trị trung bình và nêu ý nghĩa.
\loigiai{
Giá trị trung bình = [(17)+(18)+19+(20)+(21)]/5 = $19\,\text{s}$. Đây là giá trị đại diện tốt nhất của loạt đo khi các sai lệch ngẫu nhiên phân bố hai phía.
}
\end{bt}

% ===== Câu 99 =====
% ID: L10C1B3-23-TLN
% Mức: TH
\dangbt{Sai số}
\begin{ex}
% ID: L10C1B3-23-TLN
% Mức: TH
Kết quả đo chiều dài là (25,0 ± 0,5) cm. Sai số tỉ đối của phép đo theo đơn vị phần trăm bằng bao nhiêu?
\shortans{2}
\loigiai{
Sai số tỉ đối: (0,5/25,0) × 100% = 2%.
}
\end{ex}

% ===== Câu 100 =====
% ID: L10C1B3-23-TN
% Mức: TH
\begin{ex}
% ID: L10C1B3-23-TN
% Mức: TH
Ba kết quả đo thời gian là $18,0\,\text{s}$; $20,0\,\text{s}$; $22,0\,\text{s}$. Sai số ngẫu nhiên trung bình bằng bao nhiêu?
\choice
{$1,0\,\text{s}$}
{\True $1,33\,\text{s}$}
{$2,0\,\text{s}$}
{$4,0\,\text{s}$}
\loigiai{
Giá trị trung bình là $20,0\,\text{s}$. Sai số ngẫu nhiên trung bình = (2+0+2)/3 = $1,33\,\text{s}$. Chọn B.
}
\end{ex}

% ===== Câu 101 =====
% ID: L10C1B3-24-DS
% Mức: VD
\dangbt{Phép đo trực tiếp và phép đo gián tiếp}

% ===== Câu 102 =====
% ID: L10C1B3-24-TL
% Mức: VDC
\dangbt{Tính giá trị trung bình của đại lượng đo}
\begin{bt}
% ID: L10C1B3-24-TL
% Mức: VDC
Một học sinh đo thời gian 5 lần được $18\,\text{s}$, $19\,\text{s}$, $20\,\text{s}$, $21\,\text{s}$, $22\,\text{s}$. Hãy tính giá trị trung bình và nêu ý nghĩa.
\loigiai{
Giá trị trung bình = [(18)+(19)+20+(21)+(22)]/5 = $20\,\text{s}$. Đây là giá trị đại diện tốt nhất của loạt đo khi các sai lệch ngẫu nhiên phân bố hai phía.
}
\end{bt}

% ===== Câu 103 =====
% ID: L10C1B3-24-TLN
% Mức: TH
\dangbt{Sai số}
\begin{ex}
% ID: L10C1B3-24-TLN
% Mức: TH
Kết quả đo chiều dài là (25,0 ± 0,5) cm. Sai số tỉ đối của phép đo theo đơn vị phần trăm bằng bao nhiêu?
\shortans{2}
\loigiai{
Sai số tỉ đối: (0,5/25,0) × 100% = 2%.
}
\end{ex}

% ===== Câu 104 =====
% ID: L10C1B3-24-TN
% Mức: TH
\begin{ex}
% ID: L10C1B3-24-TN
% Mức: TH
Sai số nào có thể loại trừ trước khi đo?
\choice
{\True Sai số hệ thống}
{Sai số ngẫu nhiên}
{Sai số dụng cụ}
{Sai số tuyệt đối}
\loigiai{
Sai số hệ thống có thể loại trừ trước khi đo bằng cách hiệu chỉnh thiết bị đo và điều chỉnh điều kiện đo.
}
\end{ex}

% ===== Câu 105 =====
% ID: L10C1B3-25-DS
% Mức: VD
\dangbt{Phép đo trực tiếp và phép đo gián tiếp}
\begin{ex}
% ID: L10C1B3-25-DS
% Mức: VD
Một học sinh đo đường kính $d$ của một viên bi thép bằng thước kẹp có độ chia nhỏ nhất $0,02$ mm. Kết quả 5 lần đo lần lượt là: $12,34$ mm; $12,36$ mm; $12,32$ mm; $12,36$ mm; $12,32$ mm. Học sinh tính được giá trị trung bình $\bar{d} = 12,34$ mm và sai số dụng cụ $\Delta d_{dc} = 0,02$ mm. Sau đó học sinh tính thể tích viên bi theo công thức $V = \dfrac{\pi d^3}{6}$ và ghi kết quả. Xét các mệnh đề sau:
\choiceTF
{\True Sai số ngẫu nhiên của phép đo đường kính là $\Delta d_{nn} = 0,02$ mm.}
{Sai số tuyệt đối của phép đo đường kính là $\Delta d = 0,02$ mm.}
{\True Phép đo thể tích viên bi là phép đo gián tiếp vì $V$ được tính qua đại lượng đo trực tiếp là $d$.}
{Sai số tương đối của thể tích viên bi được tính theo công thức $\dfrac{\Delta V}{\bar{V}} = \dfrac{\Delta d}{\bar{d}}$.}
\loigiai{
\begin{itemchoice}
\item (Đúng) Sai số ngẫu nhiên của phép đo đường kính là $\Delta d_{nn} = 0,02$ mm. (Vì): Sai số ngẫu nhiên bằng giá trị trung bình của các độ lệch tuyệt đối: $\Delta d_{nn} = \dfrac{|12,34-12,34|+|12,36-12,34|+|12,32-12,34|+|12,36-12,34|+|12,32-12,34|}{5} = \dfrac{0+0,02+0,02+0,02+0,02}{5} = 0,016 \approx 0,02$ mm
\item (Sai) Sai số tuyệt đối của phép đo đường kính là $\Delta d = 0,02$ mm. (Vì): ố tuyệt đối của phép đo là $\Delta d = \Delta d_{nn} + \Delta d_{dc} = 0,02 + 0,02 = 0,04$ mm, không phải $0,02$ mm
\item (Đúng) Phép đo thể tích viên bi là phép đo gián tiếp vì $V$ được tính qua đại lượng đo trực tiếp là $d$. (Vì): Thể tích $V$ không đo trực tiếp bằng dụng cụ mà được tính thông qua đại lượng đo trực tiếp là đường kính $d$, nên đây là phép đo gián tiếp
\item (Sai) Sai số tương đối của thể tích viên bi được tính theo công thức $\dfrac{\Delta V}{\bar{V}} = \dfrac{\Delta d}{\bar{d}}$. (Vì): Vì $V = \dfrac{\pi d^3}{6}$, áp dụng quy tắc sai số cho hàm lũy thừa ta có $\dfrac{\Delta V}{\bar{V}} = 3\dfrac{\Delta d}{\bar{d}}$, không phải $\dfrac{\Delta d}{\bar{d}}$
\end{itemchoice}
}
\end{ex}

% ===== Câu 106 =====
% ID: L10C1B3-25-TL
% Mức: TH
\dangbt{Viết số dưới dạng kí hiệu khoa học với số chữ số có nghĩa cho trước}
\begin{bt}
% ID: L10C1B3-25-TL
% Mức: TH
Đưa số 0,0001234 về kí hiệu khoa học với 3 chữ số có nghĩa và giải thích cách làm.
\loigiai{
Dời dấu phẩy để hệ số nằm trong [1;10): 0,0001234 = 1,23×10^-4. Hệ số 1,23 có 3 chữ số có nghĩa.
}
\end{bt}

% ===== Câu 107 =====
% ID: L10C1B3-25-TLN
% Mức: TH
\dangbt{Sai số}
\begin{ex}
% ID: L10C1B3-25-TLN
% Mức: TH
Kết quả đo chiều dài là (25,0 ± 0,5) cm. Sai số tỉ đối của phép đo theo đơn vị phần trăm bằng bao nhiêu?
\shortans{2}
\loigiai{
Sai số tỉ đối: (0,5/25,0) × 100% = 2%.
}
\end{ex}

% ===== Câu 108 =====
% ID: L10C1B3-25-TN
% Mức: TH
\begin{ex}
% ID: L10C1B3-25-TN
% Mức: TH
Kết quả đo đường kính của một viên bi là $d=(2,48\pm0,02)\,\mathrm{cm}$. Cách ghi nào sau đây là đúng?
\choice
{$d=(2,480\pm0,02)\,\mathrm{cm}$}
{\True $d=(2,48\pm0,02)\,\mathrm{cm}$}
{$d=(2,5\pm0,02)\,\mathrm{cm}$}
{$d=(2,48\pm0,2)\,\mathrm{cm}$}
\loigiai{
Giá trị trung bình và sai số tuyệt đối cần được ghi với cùng bậc thập phân. Do đó cách ghi đúng là:
 $d=(2,48\pm0,02) \mathrm{cm}.$
}
\end{ex}

% ===== Câu 109 =====
% ID: L10C1B3-26-DS
% Mức: VD
\dangbt{Phép đo trực tiếp và phép đo gián tiếp}

% ===== Câu 110 =====
% ID: L10C1B3-26-TL
% Mức: TH
\dangbt{Viết số dưới dạng kí hiệu khoa học với số chữ số có nghĩa cho trước}
\begin{bt}
% ID: L10C1B3-26-TL
% Mức: TH
Đưa số 0,0002468 về kí hiệu khoa học với 3 chữ số có nghĩa và giải thích cách làm.
\loigiai{
Dời dấu phẩy để hệ số nằm trong [1;10): 0,0002468 = 2,47×10^-4. Hệ số 2,47 có 3 chữ số có nghĩa.
}
\end{bt}

% ===== Câu 111 =====
% ID: L10C1B3-26-TLN
% Mức: TH
\dangbt{Sai số}
\begin{ex}
% ID: L10C1B3-26-TLN
% Mức: TH
Kết quả đo chiều dài là (25,0 ± 0,5) cm. Sai số tỉ đối của phép đo theo đơn vị phần trăm bằng bao nhiêu?
\shortans{2}
\loigiai{
Sai số tỉ đối: (0,5/25,0) × 100% = 2%.
}
\end{ex}

% ===== Câu 112 =====
% ID: L10C1B3-26-TN
% Mức: VD
\begin{ex}
% ID: L10C1B3-26-TN
% Mức: VD
Ba kết quả đo thời gian là $19,0\,\text{s}$; $21,0\,\text{s}$; $23,0\,\text{s}$. Sai số ngẫu nhiên trung bình bằng bao nhiêu?
\choice
{$1,0\,\text{s}$}
{\True $1,33\,\text{s}$}
{$2,0\,\text{s}$}
{$4,0\,\text{s}$}
\loigiai{
Giá trị trung bình là $21,0\,\text{s}$. Sai số ngẫu nhiên trung bình = (2+0+2)/3 = $1,33\,\text{s}$. Chọn B.
}
\end{ex}

% ===== Câu 113 =====
% ID: L10C1B3-27-DS
% Mức: VD
\dangbt{Phép đo trực tiếp và phép đo gián tiếp}

% ===== Câu 114 =====
% ID: L10C1B3-27-TL
% Mức: VD
\dangbt{Viết số dưới dạng kí hiệu khoa học với số chữ số có nghĩa cho trước}
\begin{bt}
% ID: L10C1B3-27-TL
% Mức: VD
Đưa số 0,0003702 về kí hiệu khoa học với 3 chữ số có nghĩa và giải thích cách làm.
\loigiai{
Dời dấu phẩy để hệ số nằm trong [1;10): 0,0003702 = 3,7×10^-4. Hệ số 3,7 có 3 chữ số có nghĩa.
}
\end{bt}

% ===== Câu 115 =====
% ID: L10C1B3-27-TLN
% Mức: VD
\dangbt{Sai số}
\begin{ex}
% ID: L10C1B3-27-TLN
% Mức: VD
Kết quả đo được ghi (24 ± 1) cm. Tính sai số tương đối theo phần trăm, làm tròn đến hàng đơn vị. Chỉ ghi số.
\shortans{4}
\loigiai{
Sai số tương đối = 1/24×100% = 4,17%, làm tròn được 4%.
}
\end{ex}

% ===== Câu 116 =====
% ID: L10C1B3-27-TN
% Mức: VD
\begin{ex}
% ID: L10C1B3-27-TN
% Mức: VD
Ba kết quả đo thời gian là $20,0\,\text{s}$; $22,0\,\text{s}$; $24,0\,\text{s}$. Sai số ngẫu nhiên trung bình bằng bao nhiêu?
\choice
{$1,0\,\text{s}$}
{\True $1,33\,\text{s}$}
{$2,0\,\text{s}$}
{$4,0\,\text{s}$}
\loigiai{
Giá trị trung bình là $22,0\,\text{s}$. Sai số ngẫu nhiên trung bình = (2+0+2)/3 = $1,33\,\text{s}$. Chọn B.
}
\end{ex}

% ===== Câu 117 =====
% ID: L10C1B3-28-DS
% Mức: VD
\dangbt{Phép đo trực tiếp và phép đo gián tiếp}

% ===== Câu 118 =====
% ID: L10C1B3-28-TL
% Mức: VD
\dangbt{Viết số dưới dạng kí hiệu khoa học với số chữ số có nghĩa cho trước}
\begin{bt}
% ID: L10C1B3-28-TL
% Mức: VD
Đưa số 0,0004936 về kí hiệu khoa học với 3 chữ số có nghĩa và giải thích cách làm.
\loigiai{
Dời dấu phẩy để hệ số nằm trong [1;10): 0,0004936 = 4,94×10^-4. Hệ số 4,94 có 3 chữ số có nghĩa.
}
\end{bt}

% ===== Câu 119 =====
% ID: L10C1B3-28-TLN
% Mức: VD
\dangbt{Sai số}
\begin{ex}
% ID: L10C1B3-28-TLN
% Mức: VD
Kết quả đo được ghi (25 ± 1) cm. Tính sai số tương đối theo phần trăm, làm tròn đến hàng đơn vị. Chỉ ghi số.
\shortans{4}
\loigiai{
Sai số tương đối = 1/25×100% = 4%, làm tròn được 4%.
}
\end{ex}

% ===== Câu 120 =====
% ID: L10C1B3-28-TN
% Mức: VD
\begin{ex}
% ID: L10C1B3-28-TN
% Mức: VD
Ba kết quả đo thời gian là $21,0\,\text{s}$; $23,0\,\text{s}$; $25,0\,\text{s}$. Sai số ngẫu nhiên trung bình bằng bao nhiêu?
\choice
{$1,0\,\text{s}$}
{\True $1,33\,\text{s}$}
{$2,0\,\text{s}$}
{$4,0\,\text{s}$}
\loigiai{
Giá trị trung bình là $23,0\,\text{s}$. Sai số ngẫu nhiên trung bình = (2+0+2)/3 = $1,33\,\text{s}$. Chọn B.
}
\end{ex}

% ===== Câu 121 =====
% ID: L10C1B3-29-DS
% Mức: VD
\dangbt{Phép đo trực tiếp và phép đo gián tiếp}
\begin{ex}
% ID: L10C1B3-29-DS
% Mức: VD
Phân loại các phép đo sau là trực tiếp hay gián tiếp.
\choiceTF
{\True Đo chiều dài bằng thước là trực tiếp.}
{\True Tính vận tốc từ quãng đường và thời gian là gián tiếp.}
{Đọc nhiệt độ trên nhiệt kế là gián tiếp.}
{\True Tính khối lượng riêng từ m và $V$ là gián tiếp.}
\loigiai{
\begin{itemchoice}
\item (Đúng) Đo chiều dài bằng thước là trực tiếp.
\item (Đúng) Tính vận tốc từ quãng đường và thời gian là gián tiếp.
\item (Sai) Đọc nhiệt độ trên nhiệt kế là gián tiếp.
\item (Đúng) Tính khối lượng riêng từ m và $V$ là gián tiếp.
\end{itemchoice}
}
\end{ex}

% ===== Câu 122 =====
% ID: L10C1B3-29-TL
% Mức: VD
\dangbt{Viết số dưới dạng kí hiệu khoa học với số chữ số có nghĩa cho trước}
\begin{bt}
% ID: L10C1B3-29-TL
% Mức: VD
Đưa số 0,0006170 về kí hiệu khoa học với 3 chữ số có nghĩa và giải thích cách làm.
\loigiai{
Dời dấu phẩy để hệ số nằm trong [1;10): 0,0006170 = 6,17×10^-4. Hệ số 6,17 có 3 chữ số có nghĩa.
}
\end{bt}

% ===== Câu 123 =====
% ID: L10C1B3-29-TLN
% Mức: VD
\dangbt{Sai số}
\begin{ex}
% ID: L10C1B3-29-TLN
% Mức: VD
Hai độ dài được đo là a = (12,0 ± 0,1) cm và b = (8,0 ± 0,2) cm. Với s = a + b, sai số tuyệt đối của s theo đơn vị cm bằng bao nhiêu?
\shortans{0,3}
\loigiai{
Với phép cộng, sai số tuyệt đối cộng lại: Δs = 0,1 + 0,2 = $0,3\,\text{cm}$.
}
\end{ex}

% ===== Câu 124 =====
% ID: L10C1B3-29-TN
% Mức: VD
\begin{ex}
% ID: L10C1B3-29-TN
% Mức: VD
Một đại lượng có giá trị đo được là $A=(20,0\pm0,4)\,\mathrm{cm}$. Sai số tỉ đối của phép đo là
\choice
{$0,4\%$}
{\True $2\%$}
{$5\%$}
{$20\%$}
\loigiai{
Sai số tỉ đối: $\delta A = \dfrac{\Delta A}{\overline{A}} \cdot 100\% = \dfrac{0,4}{20,0} \cdot 100\% = 2\%$.
}
\end{ex}

% ===== Câu 125 =====
% ID: L10C1B3-30-DS
% Mức: NB
\begin{ex}
% ID: L10C1B3-30-DS
% Mức: NB
Xét các phát biểu về sai số phép đo.
\choiceTF
{\True Sai số tuyệt đối có cùng đơn vị với đại lượng đo.}
{\True Sai số tương đối thường biểu diễn bằng phần trăm.}
{Tăng số lần đo luôn làm sai số dụng cụ bằng 0.}
{\True Sai số ngẫu nhiên có thể giảm khi đo lặp lại hợp lí.}
\loigiai{
\begin{itemchoice}
\item (Đúng) Sai số tuyệt đối có cùng đơn vị với đại lượng đo. (Vì): ; b đúng theo δ = ΔA/A·100%; c sai vì sai số dụng cụ phụ thuộc dụng cụ; d đúng vì đo lặp giúp giảm ảnh hưởng ngẫu nhiên
\item (Đúng) Sai số tương đối thường biểu diễn bằng phần trăm.
\item (Sai) Tăng số lần đo luôn làm sai số dụng cụ bằng 0.
\item (Đúng) Sai số ngẫu nhiên có thể giảm khi đo lặp lại hợp lí.
\end{itemchoice}
}
\end{ex}

% ===== Câu 126 =====
% ID: L10C1B3-30-TL
% Mức: VDC
\dangbt{Viết số dưới dạng kí hiệu khoa học với số chữ số có nghĩa cho trước}
\begin{bt}
% ID: L10C1B3-30-TL
% Mức: VDC
Đưa số 0,0007404 về kí hiệu khoa học với 3 chữ số có nghĩa và giải thích cách làm.
\loigiai{
Dời dấu phẩy để hệ số nằm trong [1;10): 0,0007404 = 7,4×10^-4. Hệ số 7,4 có 3 chữ số có nghĩa.
}
\end{bt}

% ===== Câu 127 =====
% ID: L10C1B3-30-TLN
% Mức: VD
\dangbt{Sai số}
\begin{ex}
% ID: L10C1B3-30-TLN
% Mức: VD
Hai độ dài được đo là a = (12,0 ± 0,1) cm và b = (8,0 ± 0,2) cm. Với s = a + b, sai số tuyệt đối của s theo đơn vị cm bằng bao nhiêu?
\shortans{0,3}
\loigiai{
Với phép cộng, sai số tuyệt đối cộng lại: Δs = 0,1 + 0,2 = $0,3\,\text{cm}$.
}
\end{ex}

% ===== Câu 128 =====
% ID: L10C1B3-30-TN
% Mức: VD
\begin{ex}
% ID: L10C1B3-30-TN
% Mức: VD
Một vật đi được quãng đường $s=(100\pm1)\,\mathrm{m}$ trong thời gian $t=(20,0\pm0,2)\,\mathrm{s}$. Tốc độ trung bình của vật gần đúng là
\choice
{$2,0\,\mathrm{m/s}$}
{\True $5,0\,\mathrm{m/s}$}
{$20\,\mathrm{m/s}$}
{$100\,\mathrm{m/s}$}
\loigiai{
Tốc độ trung bình:
 $v=\frac{s}{t}=\frac{100}{20,0}=5,0 \mathrm{m/s}.$
}
\end{ex}

% ===== Câu 129 =====
% ID: L10C1B3-31-DS
% Mức: NB
\begin{ex}
% ID: L10C1B3-31-DS
% Mức: NB
Xét các phát biểu về sai số phép đo.
\choiceTF
{\True Sai số tuyệt đối có cùng đơn vị với đại lượng đo.}
{\True Sai số tương đối thường biểu diễn bằng phần trăm.}
{Tăng số lần đo luôn làm sai số dụng cụ bằng 0.}
{\True Sai số ngẫu nhiên có thể giảm khi đo lặp lại hợp lí.}
\loigiai{
\begin{itemchoice}
\item (Đúng) Sai số tuyệt đối có cùng đơn vị với đại lượng đo. (Vì): ; b đúng theo δ = ΔA/A·100%; c sai vì sai số dụng cụ phụ thuộc dụng cụ; d đúng vì đo lặp giúp giảm ảnh hưởng ngẫu nhiên
\item (Đúng) Sai số tương đối thường biểu diễn bằng phần trăm.
\item (Sai) Tăng số lần đo luôn làm sai số dụng cụ bằng 0.
\item (Đúng) Sai số ngẫu nhiên có thể giảm khi đo lặp lại hợp lí.
\end{itemchoice}
}
\end{ex}

% ===== Câu 130 =====
% ID: L10C1B3-31-TL
% Mức: TH
\dangbt{Xác định chữ số có nghĩa trong giá trị đo}
\begin{bt}
% ID: L10C1B3-31-TL
% Mức: TH
Giải thích số chữ số có nghĩa của các giá trị 0,00520; 12,00 và 105,0.
\loigiai{
0,00520 có 3 chữ số có nghĩa (5, 2, 0); các số 0 đứng đầu không tính. 12,00 có 4 chữ số có nghĩa. 105,0 có 4 chữ số có nghĩa vì số 0 nằm giữa và số 0 cuối phần thập phân đều có nghĩa.
}
\end{bt}

% ===== Câu 131 =====
% ID: L10C1B3-31-TLN
% Mức: VD
\dangbt{Sai số}
\begin{ex}
% ID: L10C1B3-31-TLN
% Mức: VD
Hai độ dài được đo là a = (12,0 ± 0,1) cm và b = (8,0 ± 0,2) cm. Với s = a + b, sai số tuyệt đối của s theo đơn vị cm bằng bao nhiêu?
\shortans{0,3}
\loigiai{
Với phép cộng, sai số tuyệt đối cộng lại: Δs = 0,1 + 0,2 = $0,3\,\text{cm}$.
}
\end{ex}

% ===== Câu 132 =====
% ID: L10C1B3-31-TN
% Mức: VD
\begin{ex}
% ID: L10C1B3-31-TN
% Mức: VD
Kết quả đo khối lượng của một vật là $m=(150\pm3)\,\mathrm{g}$. Kết quả nào dưới đây biểu diễn đúng sai số tỉ đối?
\choice
{$0,02\%$}
{$0,2\%$}
{\True $2\%$}
{$20\%$}
\loigiai{
Sai số tỉ đối:
  \delta m=\frac{\Delta m}{\overline{m}}\cdot100\%
 =\frac{3}{150}\cdot100\%=2\%.
}
\end{ex}

% ===== Câu 133 =====
% ID: L10C1B3-32-DS
% Mức: TH
\begin{ex}
% ID: L10C1B3-32-DS
% Mức: TH
Xét các phát biểu về sai số phép đo.
\choiceTF
{\True Sai số tuyệt đối có cùng đơn vị với đại lượng đo.}
{\True Sai số tương đối thường biểu diễn bằng phần trăm.}
{Tăng số lần đo luôn làm sai số dụng cụ bằng 0.}
{\True Sai số ngẫu nhiên có thể giảm khi đo lặp lại hợp lí.}
\loigiai{
\begin{itemchoice}
\item (Đúng) Sai số tuyệt đối có cùng đơn vị với đại lượng đo. (Vì): ; b đúng theo δ = ΔA/A·100%; c sai vì sai số dụng cụ phụ thuộc dụng cụ; d đúng vì đo lặp giúp giảm ảnh hưởng ngẫu nhiên
\item (Đúng) Sai số tương đối thường biểu diễn bằng phần trăm.
\item (Sai) Tăng số lần đo luôn làm sai số dụng cụ bằng 0.
\item (Đúng) Sai số ngẫu nhiên có thể giảm khi đo lặp lại hợp lí.
\end{itemchoice}
}
\end{ex}

% ===== Câu 134 =====
% ID: L10C1B3-32-TL
% Mức: TH
\dangbt{Xác định chữ số có nghĩa trong giá trị đo}
\begin{bt}
% ID: L10C1B3-32-TL
% Mức: TH
Giải thích số chữ số có nghĩa của các giá trị 0,00520; 12,00 và 105,0.
\loigiai{
0,00520 có 3 chữ số có nghĩa (5, 2, 0); các số 0 đứng đầu không tính. 12,00 có 4 chữ số có nghĩa. 105,0 có 4 chữ số có nghĩa vì số 0 nằm giữa và số 0 cuối phần thập phân đều có nghĩa.
}
\end{bt}

% ===== Câu 135 =====
% ID: L10C1B3-32-TLN
% Mức: VD
\dangbt{Sai số}
\begin{ex}
% ID: L10C1B3-32-TLN
% Mức: VD
Hai độ dài được đo là a = (12,0 ± 0,1) cm và b = (8,0 ± 0,2) cm. Với s = a + b, sai số tuyệt đối của s theo đơn vị cm bằng bao nhiêu?
\shortans{0,3}
\loigiai{
Với phép cộng, sai số tuyệt đối cộng lại: Δs = 0,1 + 0,2 = $0,3\,\text{cm}$.
}
\end{ex}

% ===== Câu 136 =====
% ID: L10C1B3-32-TN
% Mức: VD
\begin{ex}
% ID: L10C1B3-32-TN
% Mức: VD
Khối lượng riêng của một vật được tính theo công thức $\rho=\dfrac{m}{V}$.
 Biết $m=(200\pm4)\,\mathrm{g}$ và $V=(50\pm1)\,\mathrm{cm^3}$.
 Sai số tỉ đối của khối lượng riêng là
\choice
{$2\%$}
{\True $4\%$}
{$5\%$}
{$8\%$}
\loigiai{
Sai số tỉ đối của khối lượng:
 $\delta m=\frac{4}{200}\cdot100\%=2\%.$
 Sai số tỉ đối của thể tích:
 $\delta V=\frac{1}{50}\cdot100\%=2\%.$
 Với phép chia:
 $\delta \rho=\delta m+\delta V=2\%+2\%=4\%.$
}
\end{ex}

% ===== Câu 137 =====
% ID: L10C1B3-33-DS
% Mức: TH
\begin{ex}
% ID: L10C1B3-33-DS
% Mức: TH
Xét các phát biểu về sai số phép đo.
\choiceTF
{\True Sai số tuyệt đối có cùng đơn vị với đại lượng đo.}
{\True Sai số tương đối thường biểu diễn bằng phần trăm.}
{Tăng số lần đo luôn làm sai số dụng cụ bằng 0.}
{\True Sai số ngẫu nhiên có thể giảm khi đo lặp lại hợp lí.}
\loigiai{
\begin{itemchoice}
\item (Đúng) Sai số tuyệt đối có cùng đơn vị với đại lượng đo. (Vì): ; b đúng theo δ = ΔA/A·100%; c sai vì sai số dụng cụ phụ thuộc dụng cụ; d đúng vì đo lặp giúp giảm ảnh hưởng ngẫu nhiên
\item (Đúng) Sai số tương đối thường biểu diễn bằng phần trăm.
\item (Sai) Tăng số lần đo luôn làm sai số dụng cụ bằng 0.
\item (Đúng) Sai số ngẫu nhiên có thể giảm khi đo lặp lại hợp lí.
\end{itemchoice}
}
\end{ex}

% ===== Câu 138 =====
% ID: L10C1B3-33-TL
% Mức: VD
\dangbt{Xác định chữ số có nghĩa trong giá trị đo}
\begin{bt}
% ID: L10C1B3-33-TL
% Mức: VD
Giải thích số chữ số có nghĩa của các giá trị 0,00520; 12,00 và 105,0.
\loigiai{
0,00520 có 3 chữ số có nghĩa (5, 2, 0); các số 0 đứng đầu không tính. 12,00 có 4 chữ số có nghĩa. 105,0 có 4 chữ số có nghĩa vì số 0 nằm giữa và số 0 cuối phần thập phân đều có nghĩa.
}
\end{bt}

% ===== Câu 139 =====
% ID: L10C1B3-33-TLN
% Mức: VD
\dangbt{Sai số}
\begin{ex}
% ID: L10C1B3-33-TLN
% Mức: VD
Một hình chữ nhật có chiều dài a = (20,0 ± 0,2) cm và chiều rộng b = (10,0 ± 0,1) cm. Sai số tỉ đối gần đúng của diện tích theo đơn vị phần trăm bằng bao nhiêu?
\shortans{2}
\loigiai{
Sai số tỉ đối của diện tích gần bằng tổng sai số tỉ đối: 0,2/20,0 × 100% + 0,1/10,0 × 100% = 2%.
}
\end{ex}

% ===== Câu 140 =====
% ID: L10C1B3-33-TN
% Mức: VD
\begin{ex}
% ID: L10C1B3-33-TN
% Mức: VD
Một hình chữ nhật có chiều dài $a=(20,0\pm0,2)\,\mathrm{cm}$ và chiều rộng $b=(10,0\pm0,1)\,\mathrm{cm}$. Sai số tỉ đối của diện tích $S=ab$ là
\choice
{$1\%$}
{\True $2\%$}
{$3\%$}
{$10\%$}
\loigiai{
Với tích hai đại lượng:
 $\delta S=\delta a+\delta b$ \delta a=\frac{0,2}{20,0}\cdot100\%=1\%,\qquad
 \delta b=\frac{0,1}{10,0}\cdot100\%=1\%. 
 Vậy:
 $\delta S=1\%+1\%=2\%.$
}
\end{ex}

% ===== Câu 141 =====
% ID: L10C1B3-34-DS
% Mức: TH
\begin{ex}
% ID: L10C1B3-34-DS
% Mức: TH
Một học sinh đo chiều dài $L$ và chiều rộng $W$ của một hình chữ nhật và thu được kết quả $L = (20,0 \pm 0,1)$ cm và $W = (10,0 \pm 0,1)$ cm.
\choiceTF
{\True Sai số tương đối của phép đo chiều dài là 0,5%.}
{\True Sai số tương đối của phép đo chiều rộng là 1%.}
{Diện tích của hình chữ nhật là $S = (200 \pm 2)$ cm $^2$.}
{Sai số tuyệt đối của diện tích là $2\,\text{cm}^2.$}
\loigiai{
\begin{itemchoice}
\item (Đúng) Sai số tương đối của phép đo chiều dài là 0,5%. (Vì): Sai số tương đối của chiều dài là $\delta L = \dfrac{\Delta L}{\bar{L}} \times 100\% = \dfrac{0,1}{20,0} \times 100\% = 0,5\%$
\item (Đúng) Sai số tương đối của phép đo chiều rộng là 1%. (Vì): Sai số tương đối của chiều rộng là $\delta W = \dfrac{\Delta W}{\bar{W}} \times 100\% = \dfrac{0,1}{10,0} \times 100\% = 1\%$
\item (Sai) Diện tích của hình chữ nhật là $S = (200 \pm 2)$ cm $^2$. (Vì): Giá trị trung bình của diện tích là $\bar{S} = \bar{L} \times \bar{W} = 20,0 \times 10,0 = 200\ \mathrm{cm^2}$. Sai số tương đối của diện tích là $\delta S = \delta L + \delta W = 0,5\% + 1\% = 1,5\%$. Sai số tuyệt đối của diện tích là $\Delta S = \bar{S} \times \delta S = 200 \times 0,015 = 3\ \mathrm{cm^2}$. Vậy kết quả đúng là $S = (200 \pm 3)\ \mathrm{cm^2}$, không phải $(200 \pm 2)\ \mathrm{cm^2}$
\item (Sai) Sai số tuyệt đối của diện tích là $2\,\text{cm}^2.$. (Vì): Như đã tính ở trên, sai số tuyệt đối của diện tích là $\Delta S = 3\ \mathrm{cm^2}$, không phải $2\ \mathrm{cm^2}$
\end{itemchoice}
}
\end{ex}

% ===== Câu 142 =====
% ID: L10C1B3-34-TL
% Mức: VD
\dangbt{Xác định chữ số có nghĩa trong giá trị đo}
\begin{bt}
% ID: L10C1B3-34-TL
% Mức: VD
Giải thích số chữ số có nghĩa của các giá trị 0,00520; 12,00 và 105,0.
\loigiai{
0,00520 có 3 chữ số có nghĩa (5, 2, 0); các số 0 đứng đầu không tính. 12,00 có 4 chữ số có nghĩa. 105,0 có 4 chữ số có nghĩa vì số 0 nằm giữa và số 0 cuối phần thập phân đều có nghĩa.
}
\end{bt}

% ===== Câu 143 =====
% ID: L10C1B3-34-TLN
% Mức: VD
\dangbt{Sai số}
\begin{ex}
% ID: L10C1B3-34-TLN
% Mức: VD
Một hình chữ nhật có chiều dài a = (20,0 ± 0,2) cm và chiều rộng b = (10,0 ± 0,1) cm. Sai số tỉ đối gần đúng của diện tích theo đơn vị phần trăm bằng bao nhiêu?
\shortans{2}
\loigiai{
Sai số tỉ đối của diện tích gần bằng tổng sai số tỉ đối: 0,2/20,0 × 100% + 0,1/10,0 × 100% = 2%.
}
\end{ex}

% ===== Câu 144 =====
% ID: L10C1B3-34-TN
% Mức: VD
\begin{ex}
% ID: L10C1B3-34-TN
% Mức: VD
Một học sinh đo chiều dài của một vật và ghi được kết quả $l=(25,0\pm0,5)\,\mathrm{cm}$. Sai số tỉ đối của phép đo là
\choice
{$0,5\%$}
{$1\%$}
{\True $2\%$}
{$5\%$}
\loigiai{
Sai số tỉ đối:
  \delta l=\frac{\Delta l}{\overline{l}}\cdot100\%
 =\frac{0,5}{25,0}\cdot100\%=2\%.
}
\end{ex}

% ===== Câu 145 =====
% ID: L10C1B3-35-DS
% Mức: TH
\begin{ex}
% ID: L10C1B3-35-DS
% Mức: TH
Một học sinh đo chu kì dao động của con lắc đơn và thu được giá trị trung bình $\bar{T} = 1,567$ s. Sai số tuyệt đối của phép đo là $\Delta T = 0,03$ s.
\choiceTF
{Sai số tương đối của phép đo là khoảng 0,19%.}
{\True Kết quả phép đo được ghi là $T = (1,57 \pm 0,03)$ s.}
{Sai số tương đối của phép đo là khoảng 2%.}
{\True Sai số tương đối của phép đo là khoảng 1,9%.}
\loigiai{
\begin{itemchoice}
\item (Sai) Sai số tương đối của phép đo là khoảng 0,19%. (Vì): ố tương đối của phép đo là $\delta T = \frac{\Delta T}{\bar{T}} \times 100\% = \frac{0,03}{1,567} \times 100\% \approx 1,914\%$. Vậy 0,19% là sai
\item (Đúng) Kết quả phép đo được ghi là $T = (1,57 \pm 0,03)$ s. (Vì): Sai số tuyệt đối $\Delta T = 0,03$ s (đến hàng phần trăm). Giá trị trung bình $\bar{T} = 1,567$ s phải được làm tròn đến hàng phần trăm là $1,57\,\text{s}$. Vậy kết quả đúng là $T = (1,57 \pm 0,03)$ s
\item (Sai) Sai số tương đối của phép đo là khoảng 2%. (Vì): ố tương đối $\delta T \approx 1,914\%$. Làm tròn 2% là quá thô, không chính xác bằng 1,9% hoặc 1,91%
\item (Đúng) Sai số tương đối của phép đo là khoảng 1,9%. (Vì): Sai số tương đối $\delta T \approx 1,914\%$. Làm tròn đến một chữ số thập phân là 1,9%
\end{itemchoice}
}
\end{ex}

% ===== Câu 146 =====
% ID: L10C1B3-35-TL
% Mức: VD
\dangbt{Xác định chữ số có nghĩa trong giá trị đo}
\begin{bt}
% ID: L10C1B3-35-TL
% Mức: VD
Giải thích số chữ số có nghĩa của các giá trị 0,00520; 12,00 và 105,0.
\loigiai{
0,00520 có 3 chữ số có nghĩa (5, 2, 0); các số 0 đứng đầu không tính. 12,00 có 4 chữ số có nghĩa. 105,0 có 4 chữ số có nghĩa vì số 0 nằm giữa và số 0 cuối phần thập phân đều có nghĩa.
}
\end{bt}

% ===== Câu 147 =====
% ID: L10C1B3-35-TLN
% Mức: VD
\dangbt{Sai số}
\begin{ex}
% ID: L10C1B3-35-TLN
% Mức: VD
Một hình chữ nhật có chiều dài a = (20,0 ± 0,2) cm và chiều rộng b = (10,0 ± 0,1) cm. Sai số tỉ đối gần đúng của diện tích theo đơn vị phần trăm bằng bao nhiêu?
\shortans{2}
\loigiai{
Sai số tỉ đối của diện tích gần bằng tổng sai số tỉ đối: 0,2/20,0 × 100% + 0,1/10,0 × 100% = 2%.
}
\end{ex}

% ===== Câu 148 =====
% ID: L10C1B3-35-TN
% Mức: VD
\begin{ex}
% ID: L10C1B3-35-TN
% Mức: VD
Kết quả đo khối lượng của một vật là $m=(250\pm5)\,\mathrm{g}$.
 Sai số tỉ đối của phép đo là
\choice
{$0,5\%$}
{$1\%$}
{\True $2\%$}
{$5\%$}
\loigiai{
Sai số tỉ đối:
  \delta m=\frac{\Delta m}{\overline{m}}\cdot100\%
 =\frac{5}{250}\cdot100\%=2\%.
}
\end{ex}

% ===== Câu 149 =====
% ID: L10C1B3-36-DS
% Mức: TH
\begin{ex}
% ID: L10C1B3-36-DS
% Mức: TH
Một học sinh đo khối lượng $m$ và thể tích $V$ của một vật rắn và thu được kết quả $m = (150 \pm 2)$ g và $V = (50 \pm 1)$ cm $^3$.
\choiceTF
{Sai số tương đối của khối lượng là 0,013%.}
{Sai số tương đối của thể tích là 1%.}
{\True Giá trị trung bình của khối lượng riêng là $3\,\text{g}$ /cm $^3$.}
{\True Kết quả phép đo khối lượng riêng được ghi là $\rho = (3,0 \pm 0,1)$ g/cm $^3$.}
\loigiai{
\begin{itemchoice}
\item (Sai) Sai số tương đối của khối lượng là 0,013%. (Vì): ố tương đối của khối lượng là $\delta m = \frac{\Delta m}{\bar{m}} \times 100\% = \frac{2}{150} \times 100\% \approx 1,33\%$. Vậy 0,013% là sai
\item (Sai) Sai số tương đối của thể tích là 1%. (Vì): Sai — số tương đối của thể tích là $\delta V = \frac{\Delta V}{\bar{V}} \times 100\% = \frac{1}{50} \times 100\% = 2\%$. Vậy 1% là sai
\item (Đúng) Giá trị trung bình của khối lượng riêng là $3\,\text{g}$ /cm $^3$. (Vì): Đúng — Giá trị trung bình của khối lượng riêng là $\bar{\rho} = \frac{\bar{m}}{\bar{V}} = \frac{150}{50} = 3$ g/cm $^3$
\item (Đúng) Kết quả phép đo khối lượng riêng được ghi là $\rho = (3,0 \pm 0,1)$ g/cm $^3$. (Vì): Sai số tương đối của khối lượng riêng là $\delta \rho = \delta m + \delta V = \frac{2}{150} + \frac{1}{50} = \frac{2}{150} + \frac{3}{150} = \frac{5}{150} = \frac{1}{30} \approx 0,0333$. Sai số tuyệt đối của khối lượng riêng là $\Delta \rho = \bar{\rho} \times \delta \rho = 3 \times \frac{1}{30} = 0,1$ g/cm $^3$. Kết quả phép đo được ghi là $\rho = (3,0 \pm 0,1)$ g/cm $^3$ (làm tròn $\bar{\rho}$ đến cùng hàng với $\Delta \rho$
\end{itemchoice}
}
\end{ex}

% ===== Câu 150 =====
% ID: L10C1B3-36-TL
% Mức: VDC
\dangbt{Xác định chữ số có nghĩa trong giá trị đo}
\begin{bt}
% ID: L10C1B3-36-TL
% Mức: VDC
Giải thích số chữ số có nghĩa của các giá trị 0,00520; 12,00 và 105,0.
\loigiai{
0,00520 có 3 chữ số có nghĩa (5, 2, 0); các số 0 đứng đầu không tính. 12,00 có 4 chữ số có nghĩa. 105,0 có 4 chữ số có nghĩa vì số 0 nằm giữa và số 0 cuối phần thập phân đều có nghĩa.
}
\end{bt}

% ===== Câu 151 =====
% ID: L10C1B3-36-TLN
% Mức: VD
\dangbt{Sai số}
\begin{ex}
% ID: L10C1B3-36-TLN
% Mức: VD
Một hình chữ nhật có chiều dài a = (20,0 ± 0,2) cm và chiều rộng b = (10,0 ± 0,1) cm. Sai số tỉ đối gần đúng của diện tích theo đơn vị phần trăm bằng bao nhiêu?
\shortans{2}
\loigiai{
Sai số tỉ đối của diện tích gần bằng tổng sai số tỉ đối: 0,2/20,0 × 100% + 0,1/10,0 × 100% = 2%.
}
\end{ex}

% ===== Câu 152 =====
% ID: L10C1B3-36-TN
% Mức: VD
\begin{ex}
% ID: L10C1B3-36-TN
% Mức: VD
Đo hai đoạn thẳng được $a=(12,0\pm0,1)\,\mathrm{cm}$ và $b=(8,0\pm0,2)\,\mathrm{cm}$.
 Tổng $s=a+b$ có kết quả là
\choice
{$s=(20,0\pm0,1)\,\mathrm{cm}$}
{$s=(20,0\pm0,2)\,\mathrm{cm}$}
{\True $s=(20,0\pm0,3)\,\mathrm{cm}$}
{$s=(4,0\pm0,3)\,\mathrm{cm}$}
\loigiai{
Giá trị trung bình:
 $s=a+b=12,0+8,0=20,0 \mathrm{cm}.$
 Với phép cộng:
 $\Delta s=\Delta a+\Delta b=0,1+0,2=0,3 \mathrm{cm}.$
 Vậy:
 $s=(20,0\pm0,3) \mathrm{cm}.$
}
\end{ex}

% ===== Câu 153 =====
% ID: L10C1B3-37-DS
% Mức: TH
\begin{ex}
% ID: L10C1B3-37-DS
% Mức: TH
Một học sinh đo chiều dài của một vật 5 lần và thu được các giá trị: $12,3\,\text{cm}$; $12,5\,\text{cm}$; $12,4\,\text{cm}$; $12,3\,\text{cm}$; $12,6\,\text{cm}$. Sai số dụng cụ là $0,1\,\text{cm}$.
\choiceTF
{\True Giá trị trung bình của chiều dài là $12,42\,\text{cm}$.}
{Sai số ngẫu nhiên của phép đo là $0,08\,\text{cm}$.}
{\True Sai số tuyệt đối của phép đo là $0,2\,\text{cm}$.}
{Kết quả phép đo được ghi là $L = (12,42 \pm 0,2)$ cm.}
\loigiai{
\begin{itemchoice}
\item (Đúng) Giá trị trung bình của chiều dài là $12,42\,\text{cm}$. (Vì): Giá trị trung bình của chiều dài là $\bar{L} = \frac{12,3 + 12,5 + 12,4 + 12,3 + 12,6}{5} = \frac{62,1}{5} = 12,42\,\text{cm}$
\item (Sai) Sai số ngẫu nhiên của phép đo là $0,08\,\text{cm}$. (Vì): ố ngẫu nhiên của phép đo là $\Delta L_{ng} = \frac{|12,3-12,42| + |12,5-12,42| + |12,4-12,42| + |12,3-12,42| + |12,6-12,42|}{5} = \frac{0,12 + 0,08 + 0,02 + 0,12 + 0,18}{5} = \frac{0,52}{5} = 0,104\,\text{cm}$. Làm tròn đến một chữ số có nghĩa là $0,1\,\text{cm}$
\item (Đúng) Sai số tuyệt đối của phép đo là $0,2\,\text{cm}$. (Vì): Sai số tuyệt đối của phép đo là $\Delta L = \Delta L_{ng} + \Delta L_{dc} = 0,1\,\text{cm} + 0,1\,\text{cm} = 0,2\,\text{cm}$
\item (Sai) Kết quả phép đo được ghi là $L = (12,42 \pm 0,2)$ cm. (Vì): ố tuyệt đối của phép đo là $\Delta L = 0,2\,\text{cm}$ (làm tròn đến hàng phần mười). Giá trị trung bình $\bar{L} = 12,42\,\text{cm}$ phải được làm tròn đến hàng phần mười là $12,4\,\text{cm}$. Vậy kết quả đúng phải là $L = (12,4 \pm 0,2)\,\text{cm}$
\end{itemchoice}
}
\end{ex}

% ===== Câu 154 =====
% ID: L10C1B3-37-TLN
% Mức: NB
\dangbt{Thứ nguyên}
\begin{ex}
% ID: L10C1B3-37-TLN
% Mức: NB
Trong thứ nguyên [L][T]^n của gia tốc, số mũ n bằng bao nhiêu? Chỉ ghi số.
\shortans{-2}
\loigiai{
Gia tốc có thứ nguyên [L][T]⁻² nên n = -2.
}
\end{ex}

% ===== Câu 155 =====
% ID: L10C1B3-37-TN
% Mức: VD
\dangbt{Sai số}
\begin{ex}
% ID: L10C1B3-37-TN
% Mức: VD
Một học sinh đo độ dài của một đoạn thẳng được các kết quả:
 $8,4\,\mathrm{cm}$; $8,6\,\mathrm{cm}$; $8,4\,\mathrm{cm}$; $8,6\,\mathrm{cm}$.
 Bỏ qua sai số dụng cụ. Kết quả phép đo được viết là
\choice
{$l=(8,4\pm0,1)\,\mathrm{cm}$}
{\True $l=(8,5\pm0,1)\,\mathrm{cm}$}
{$l=(8,6\pm0,1)\,\mathrm{cm}$}
{$l=(8,5\pm0,2)\,\mathrm{cm}$}
\loigiai{
Giá trị trung bình của các lần đo là $\overline{l} = \frac{8,4 + 8,6 + 8,4 + 8,6}{4} = 8,5\,\mathrm{cm}$.
Sai số tuyệt đối trung bình là $\overline{\Delta l} = \frac{|8,4-8,5| + |8,6-8,5| + |8,4-8,5| + |8,6-8,5|}{4} = \frac{0,1 + 0,1 + 0,1 + 0,1}{4} = 0,1\,\mathrm{cm}$.
Vì bỏ qua sai số dụng cụ, kết quả phép đo được viết là $l = (\overline{l} \pm \overline{\Delta l}) = (8,5 \pm 0,1)\,\mathrm{cm}$.
}
\end{ex}

% ===== Câu 156 =====
% ID: L10C1B3-38-DS
% Mức: TH
\begin{ex}
% ID: L10C1B3-38-DS
% Mức: TH
Một học sinh đo chiều dài của một vật và ghi được kết quả $l=(40,0\pm0,8)\,\mathrm{cm}$. Xét các phát biểu sau:
\choiceTF
{\True Sai số tuyệt đối của phép đo là $0,8\,\mathrm{cm}$}
{Giá trị trung bình của phép đo là $0,8\,\mathrm{cm}$}
{\True Sai số tỉ đối của phép đo là $2\%$}
{Kết quả đo có thể viết là $l=(40,0\pm2)\,\mathrm{cm}$}
\loigiai{
\begin{itemchoice}
\item (Đúng) Sai số tuyệt đối của phép đo là $0,8\,\mathrm{cm}$. (Vì): Sai số tuyệt đối của phép đo được ghi trực tiếp trong kết quả là $0,8\,\mathrm{cm}$
\item (Sai) Giá trị trung bình của phép đo là $0,8\,\mathrm{cm}$. (Vì): Giá trị trung bình của phép đo là $40,0\,\mathrm{cm}$, không phải $0,8\,\mathrm{cm}$
\item (Đúng) Sai số tỉ đối của phép đo là $2\%$. (Vì): Sai số tỉ đối của phép đo được tính bằng $\frac{\Delta l}{\bar{l}} \times 100\% = \frac{0,8}{40,0} \times 100\% = 2\%$
\item (Sai) Kết quả đo có thể viết là $l=(40,0\pm2)\,\mathrm{cm}$. (Vì): ố tuyệt đối của phép đo là $0,8\,\mathrm{cm}$, không thể viết là $2\,\mathrm{cm}$
\end{itemchoice}
}
\end{ex}

% ===== Câu 157 =====
% ID: L10C1B3-38-TLN
% Mức: NB
\dangbt{Thứ nguyên}

% ===== Câu 158 =====
% ID: L10C1B3-38-TN
% Mức: VD
\dangbt{Sai số}
\begin{ex}
% ID: L10C1B3-38-TN
% Mức: VD
Thể tích của một khối lập phương được tính theo công thức $V=a^3$. Biết cạnh khối lập phương là $a=(10,0\pm0,1)\,\mathrm{cm}$. Sai số tỉ đối của thể tích $V$ là
\choice
{$1\%$}
{$2\%$}
{\True $3\%$}
{$10\%$}
\loigiai{
Thể tích của khối lập phương được tính theo công thức $V=a^3$.
Sai số tỉ đối của thể tích $V$ được tính theo công thức $\delta V = 3 \delta a$.
Ta có sai số tuyệt đối của cạnh là $\Delta a = 0,1\,\mathrm{cm}$ và giá trị trung bình của cạnh là $a = 10,0\,\mathrm{cm}$.
Sai số tỉ đối của cạnh là $\delta a = \frac{\Delta a}{a} = \frac{0,1}{10,0} = 0,01$.
Vậy sai số tỉ đối của thể tích là $\delta V = 3 \times 0,01 = 0,03$.

huyển sang phần trăm, $\delta V = 0,03 \times 100\% = 3\%$.
}
\end{ex}

% ===== Câu 159 =====
% ID: L10C1B3-39-DS
% Mức: VD
\begin{ex}
% ID: L10C1B3-39-DS
% Mức: VD
Xét các phát biểu về sai số phép đo.
\choiceTF
{\True Sai số tuyệt đối có cùng đơn vị với đại lượng đo.}
{\True Sai số tương đối thường biểu diễn bằng phần trăm.}
{Tăng số lần đo luôn làm sai số dụng cụ bằng 0.}
{\True Sai số ngẫu nhiên có thể giảm khi đo lặp lại hợp lí.}
\loigiai{
\begin{itemchoice}
\item (Đúng) Sai số tuyệt đối có cùng đơn vị với đại lượng đo. (Vì): ; b đúng theo δ = ΔA/A·100%; c sai vì sai số dụng cụ phụ thuộc dụng cụ; d đúng vì đo lặp giúp giảm ảnh hưởng ngẫu nhiên
\item (Đúng) Sai số tương đối thường biểu diễn bằng phần trăm.
\item (Sai) Tăng số lần đo luôn làm sai số dụng cụ bằng 0.
\item (Đúng) Sai số ngẫu nhiên có thể giảm khi đo lặp lại hợp lí.
\end{itemchoice}
}
\end{ex}

% ===== Câu 160 =====
% ID: L10C1B3-39-TLN
% Mức: TH
\dangbt{Thứ nguyên}

% ===== Câu 161 =====
% ID: L10C1B3-39-TN
% Mức: VD
\dangbt{Sai số}
\begin{ex}
% ID: L10C1B3-39-TN
% Mức: VD
Một vật đi được quãng đường $s=(100\pm1)\,\mathrm{m}$ trong thời gian $t=(20,0\pm0,2)\,\mathrm{s}$. Sai số tỉ đối của tốc độ $v=\dfrac{s}{t}$ là
\choice
{$1\%$}
{\True $2\%$}
{$5\%$}
{$10\%$}
\loigiai{
Ta có $s = (100 \pm 1)\,\mathrm{m}$ và $t = (20,0 \pm 0,2)\,\mathrm{s}$.
Sai số tỉ đối của quãng đường là $\delta s = \frac{1}{100}$.
Sai số tỉ đối của thời gian là $\delta t = \frac{0,2}{20,0}$.
Khi chia hai đại lượng, sai số tỉ đối của thương bằng tổng các sai số tỉ đối của các đại lượng đó:
$\delta v = \delta s + \delta t = \frac{1}{100} + \frac{0,2}{20,0} = 0,01 + 0,01 = 0,02$.
Vậy sai số tỉ đối của tốc độ là $0,02$ hay $2\%$.
}
\end{ex}

% ===== Câu 162 =====
% ID: L10C1B3-40-DS
% Mức: VD
\begin{ex}
% ID: L10C1B3-40-DS
% Mức: VD
Một học sinh đo quãng đường $s=(100\pm1)\,\mathrm{m}$ và thời gian $t=(20,0\pm0,2)\,\mathrm{s}$. Tốc độ trung bình được tính theo công thức $v=\dfrac{s}{t}$. Xét các phát biểu sau:
\choiceTF
{\True Giá trị trung bình của tốc độ là $5,0\,\mathrm{m/s}$}
{\True Sai số tỉ đối của quãng đường là $1\%$}
{Sai số tỉ đối của thời gian là $2\%$}
{\True Sai số tỉ đối của tốc độ là $2\%$}
\loigiai{
\begin{itemchoice}
\item (Đúng) Giá trị trung bình của tốc độ là $5,0\,\mathrm{m/s}$. (Vì): Giá trị trung bình của tốc độ là $v = \frac{s}{t} = \frac{100\,\mathrm{m}}{20,0\,\mathrm{s}} = 5,0\,\mathrm{m/s}$
\item (Đúng) Sai số tỉ đối của quãng đường là $1\%$. (Vì): Sai số tỉ đối của quãng đường là $\delta s = \frac{\Delta s}{s} \times 100\% = \frac{1\,\mathrm{m}}{100\,\mathrm{m}} \times 100\% = 1\%$
\item (Sai) Sai số tỉ đối của thời gian là $2\%$. (Vì): ố tỉ đối của thời gian là $\delta t = \frac{\Delta t}{t} \times 100\% = \frac{0,2\,\mathrm{s}}{20,0\,\mathrm{s}} \times 100\% = 1\%$
\item (Đúng) Sai số tỉ đối của tốc độ là $2\%$. (Vì): Sai số tỉ đối của tốc độ là $\delta v = \delta s + \delta t = 1\% + 1\% = 2\%$
\end{itemchoice}
}
\end{ex}

% ===== Câu 163 =====
% ID: L10C1B3-40-TLN
% Mức: TH
\dangbt{Thứ nguyên}

% ===== Câu 164 =====
% ID: L10C1B3-40-TN
% Mức: VD
\dangbt{Sai số}
\begin{ex}
% ID: L10C1B3-40-TN
% Mức: VD
Đo chiều dài hai đoạn thẳng được $a=(12,0\pm0,1)\,\mathrm{cm}$ và $b=(8,0\pm0,1)\,\mathrm{cm}$. Tổng $s=a+b$ có giá trị là
\choice
{$s=(20,0\pm0,1)\,\mathrm{cm}$}
{\True $s=(20,0\pm0,2)\,\mathrm{cm}$}
{$s=(4,0\pm0,2)\,\mathrm{cm}$}
{$s=(20,0\pm0,01)\,\mathrm{cm}$}
\loigiai{
Khi cộng hai đại lượng, sai số tuyệt đối bằng tổng các sai số tuyệt đối:
$s = a + b = 12,0 + 8,0 = 20,0\,\mathrm{cm}\Delta s = \Delta a + \Delta b = 0,1 + 0,1 = 0,2\,\mathrm{cm}$.
Vậy $s = (20,0 \pm 0,2)\,\mathrm{cm}$.
}
\end{ex}

% ===== Câu 165 =====
% ID: L10C1B3-41-DS
% Mức: VD
\begin{ex}
% ID: L10C1B3-41-DS
% Mức: VD
Hai người cùng đo chiều dài của cánh cửa sổ, kết quả thu được như sau:
 
  
 • Người thứ nhất: $d = 120 \pm 1 \, \text{cm}$ (• Người thứ hai:) $d = 120 \pm 2 \, \text{cm}$
\choiceTF
{\True Sai số tỷ đối được xác định bằng tỉ số giữa sai số tuyệt đối và giá trị trung bình của chiều dài cánh cửa sổ: $\delta d=\dfrac{\Delta d}{\bar{d}} \cdot 100 \%$}
{\True Sai số tỷ đối của phép đo của người thứ nhất là $0,83 \%$}
{\True Sai số tỷ đối của phép đo của người thứ hai là $1,67\%$}
{\True Người thứ nhất đo chính xác hơn người thứ hai vì sai số tỉ đối của người thứ nhất nhỏ hơn}
\loigiai{
\begin{itemchoice}
\item (Đúng) Sai số tỷ đối được xác định bằng tỉ số giữa sai số tuyệt đối và giá trị trung bình của chiều dài cánh cửa sổ: $\delta d=\dfrac{\Delta d}{\bar{d}} \cdot 100 \%$. (Vì): Sai số tỉ đối ($\delta d$) là tỉ số giữa sai số tuyệt đối ($\Delta d$) và giá trị trung bình ($\bar{d}$) nhân với 100%, tức là $\delta d = \dfrac{\Delta d}{\bar{d}} \cdot 100\%$
\item (Đúng) Sai số tỷ đối của phép đo của người thứ nhất là $0,83 \%$. (Vì): Sai số tỉ đối của người thứ nhất là $\delta d_1 = \dfrac{1}{120} \cdot 100\% \approx 0,83\%$
\item (Đúng) Sai số tỷ đối của phép đo của người thứ hai là $1,67\%$. (Vì): Sai số tỉ đối của người thứ hai là $\delta d_2 = \dfrac{2}{120} \cdot 100\% \approx 1,67\%$
\item (Đúng) Người thứ nhất đo chính xác hơn người thứ hai vì sai số tỉ đối của người thứ nhất nhỏ hơn. (Vì): Người thứ nhất đo chính xác hơn vì sai số tỉ đối nhỏ hơn ($0,83\% \lt  1,67\%$). Sai số tỉ đối càng nhỏ thì độ chính xác của phép đo càng cao
\end{itemchoice}
}
\end{ex}

% ===== Câu 166 =====
% ID: L10C1B3-41-TLN
% Mức: TH
\dangbt{Thứ nguyên}
\begin{ex}
% ID: L10C1B3-41-TLN
% Mức: TH
Vận tốc có thứ nguyên [v] = L·T^n. Số mũ n bằng bao nhiêu?
\shortans{-1}
\loigiai{
Vì v = s/t nên [v] = L·T^(-1), do đó n = -1.
}
\end{ex}

% ===== Câu 167 =====
% ID: L10C1B3-41-TN
% Mức: VD
\dangbt{Sai số}
\begin{ex}
% ID: L10C1B3-41-TN
% Mức: VD
Thể tích của một khối lập phương được tính theo công thức $V=a^3$.
 Biết cạnh của khối lập phương là $a=(4,00\pm0,04)\,\mathrm{cm}$.
 Sai số tỉ đối của thể tích là
\choice
{$1\%$}
{$2\%$}
{\True $3\%$}
{$4\%$}
\loigiai{
Sai số tỉ đối của cạnh là $\delta a = \frac{0,04}{4,00} = 0,01$.
Vì $V=a^3$ nên sai số tỉ đối của thể tích là $\delta V = 3\delta a = 3 \times 0,01 = 0,03$.

họn C.
}
\end{ex}

% ===== Câu 168 =====
% ID: L10C1B3-42-DS
% Mức: VD
\begin{ex}
% ID: L10C1B3-42-DS
% Mức: VD
Diện tích hình chữ nhật được tính theo công thức $S=ab$. Biết $a=(30,0\pm0,3)\,\mathrm{cm}$ và $b=(20,0\pm0,4)\,\mathrm{cm}$. Xét các phát biểu sau:
\choiceTF
{\True Giá trị trung bình của diện tích là $600\,\mathrm{cm^2}$}
{\True Sai số tỉ đối của $a$ là $1\%$}
{Sai số tỉ đối của $b$ là $1\%$}
{Sai số tỉ đối của diện tích là $2\%$}
\loigiai{
\begin{itemchoice}
\item (Đúng) Giá trị trung bình của diện tích là $600\,\mathrm{cm^2}$. (Vì): Giá trị trung bình của diện tích là $\bar{S} = \bar{a} \cdot \bar{b} = 30,0 \cdot 20,0 = 600\ \mathrm{cm^2}$
\item (Đúng) Sai số tỉ đối của $a$ là $1\%$. (Vì): Sai số tỉ đối của $a$ là $\delta a = \dfrac{\Delta a}{\bar{a}} \cdot 100\% = \dfrac{0,3}{30,0} \cdot 100\% = 1\%$
\item (Sai) Sai số tỉ đối của $b$ là $1\%$. (Vì): ố tỉ đối của $b$ là $\delta b = \dfrac{\Delta b}{\bar{b}} \cdot 100\% = \dfrac{0,4}{20,0} \cdot 100\% = 2\%$, không phải $1\%$
\item (Sai) Sai số tỉ đối của diện tích là $2\%$. (Vì): ố tỉ đối của diện tích bằng tổng sai số tỉ đối của các đại lượng thành phần: $\delta S = \delta a + \delta b = 1\% + 2\% = 3\%$, không phải $2\%$
\end{itemchoice}
}
\end{ex}

% ===== Câu 169 =====
% ID: L10C1B3-42-TLN
% Mức: TH
\dangbt{Thứ nguyên}

% ===== Câu 170 =====
% ID: L10C1B3-42-TN
% Mức: VD
\dangbt{Sai số}
\begin{ex}
% ID: L10C1B3-42-TN
% Mức: VD
Một vật có khối lượng được đo trực tiếp là $m=(200\pm4)\,\mathrm{g}$. Sai số tương đối của phép đo này là
\choice
{$0{,}02$}
{\True $0,02$ hay $2\%$}
{$0,2$ hay $20\%$}
{$4\%$}
\loigiai{
Sai số tương đối của phép đo khối lượng được tính bằng công thức $\delta m = \frac{\Delta m}{\overline{m}}$.
Thay số vào công thức, ta có $\delta m = \frac{4}{200} = 0,02$.
Nếu biểu diễn dưới dạng phần trăm, sai số tương đối là $\delta m = 0,02 \times 100\% = 2\%$.
}
\end{ex}

% ===== Câu 171 =====
% ID: L10C1B3-43-DS
% Mức: VD
\begin{ex}
% ID: L10C1B3-43-DS
% Mức: VD
Thể tích của một khối lập phương được tính theo công thức $V=a^3$. Biết cạnh khối lập phương là $a=(5,00\pm0,05)\,\mathrm{cm}$. Xét các phát biểu sau:
\choiceTF
{\True Giá trị trung bình của thể tích là $125\,\mathrm{cm^3}$}
{\True Sai số tỉ đối của cạnh $a$ là $1\%$}
{Sai số tỉ đối của thể tích là $1\%$}
{\True Sai số tỉ đối của thể tích là $3\%$}
\loigiai{
\begin{itemchoice}
\item (Đúng) Giá trị trung bình của thể tích là $125\,\mathrm{cm^3}$. (Vì): Giá trị trung bình của thể tích được tính bằng cách lấy giá trị trung bình của cạnh lập phương lên lũy thừa 3: $V = a^3 = (5,00)^3 = 125 \mathrm{cm^3}$
\item (Đúng) Sai số tỉ đối của cạnh $a$ là $1\%$. (Vì): Sai số tỉ đối của cạnh $a$ được tính bằng công thức $\delta a = \frac{\Delta a}{a} \times 100\% = \frac{0,05}{5,00} \times 100\% = 1\%$
\item (Sai) Sai số tỉ đối của thể tích là $1\%$. (Vì): ố tỉ đối của thể tích $V$ được tính theo công thức $\delta V = 3 \delta a = 3 \times 1\% = 3\%$
\item (Đúng) Sai số tỉ đối của thể tích là $3\%$.
\end{itemchoice}
}
\end{ex}

% ===== Câu 172 =====
% ID: L10C1B3-43-TLN
% Mức: TH
\dangbt{Thứ nguyên}

% ===== Câu 173 =====
% ID: L10C1B3-43-TN
% Mức: VD
\dangbt{Sai số}
\begin{ex}
% ID: L10C1B3-43-TN
% Mức: VD
Diện tích hình chữ nhật được tính theo công thức $S=ab$. Biết $a=(20,0\pm0,2)\,\mathrm{cm}$ và $b=(15,0\pm0,3)\,\mathrm{cm}$. Sai số tỉ đối của diện tích $S$ là
\choice
{$1\%$}
{$2\%$}
{\True $3\%$}
{$5\%$}
\loigiai{
Với phép nhân $S=ab$, sai số tỉ đối được tính theo công thức $\delta S = \delta a + \delta b$. Ta có: $\delta a = \dfrac{\Delta a}{a} = \dfrac{0,2}{20,0} \times 100\% = 1\%$ và $\delta b = \dfrac{\Delta b}{b} = \dfrac{0,3}{15,0} \times 100\% = 2\%$. Do đó $\delta S = 1\% + 2\% = 3\%$.
}
\end{ex}

% ===== Câu 174 =====
% ID: L10C1B3-44-DS
% Mức: NB
\dangbt{Thứ nguyên}
\begin{ex}
% ID: L10C1B3-44-DS
% Mức: NB
Xét các phát biểu về thứ nguyên.
\choiceTF
{\True Chiều dài có thứ nguyên [L].}
{\True Gia tốc có thứ nguyên [L][T]⁻².}
{Lực có thứ nguyên [M][L][T]⁻¹.}
{\True Năng lượng có thứ nguyên [M][L]²[T]⁻².}
\loigiai{
\begin{itemchoice}
\item (Đúng) Chiều dài có thứ nguyên [L]. (Vì): , b đúng; c sai vì [F]=[M][L][T]⁻²; d đúng
\item (Đúng) Gia tốc có thứ nguyên [L][T]⁻².
\item (Sai) Lực có thứ nguyên [M][L][T]⁻¹.
\item (Đúng) Năng lượng có thứ nguyên [M][L]²[T]⁻².
\end{itemchoice}
}
\end{ex}

% ===== Câu 176 =====
% ID: L10C1B3-44-TN
% Mức: VD
\dangbt{Sai số}
\begin{ex}
% ID: L10C1B3-44-TN
% Mức: VD
Đo hai đại lượng được $x=(15,0\pm0,2)\,\mathrm{cm}$ và $y=(6,0\pm0,1)\,\mathrm{cm}$. Hiệu $d=x-y$ có sai số tuyệt đối là
\choice
{$0,1\,\mathrm{cm}$}
{$0,2\,\mathrm{cm}$}
{\True $0,3\,\mathrm{cm}$}
{$9,0\,\mathrm{cm}$}
\loigiai{
Khi trừ hai đại lượng, sai số tuyệt đối vẫn bằng tổng các sai số tuyệt đối:
 $\Delta d=\Delta x+\Delta y=0,2+0,1=0,3 \mathrm{cm}.$
}
\end{ex}

% ===== Câu 177 =====
% ID: L10C1B3-45-DS
% Mức: NB
\dangbt{Thứ nguyên}

% ===== Câu 178 =====
% ID: L10C1B3-45-TLN
% Mức: TH
\begin{ex}
% ID: L10C1B3-45-TLN
% Mức: TH
Gia tốc có thứ nguyên [a] = L·T^n. Số mũ n bằng bao nhiêu?
\shortans{-2}
\loigiai{
Vì a = Δv/Δt nên [a] = L·T^(-2), do đó n = -2.
}
\end{ex}

% ===== Câu 179 =====
% ID: L10C1B3-45-TN
% Mức: NB
\begin{ex}
% ID: L10C1B3-45-TN
% Mức: NB
Thứ nguyên của vận tốc là
\choice
{[L][T]}
{\True [L][T]⁻¹}
{[M][L][T]⁻²}
{[M][L]²[T]⁻²}
\loigiai{
Từ v = s/t, suy ra [v] = [L][T]⁻¹. Chọn B.
}
\end{ex}

% ===== Câu 188 =====
% ID: L10C1B3-48-TN
% Mức: NB
\begin{ex}
% ID: L10C1B3-48-TN
% Mức: NB
Thứ nguyên của công suất là
\choice
{$M.L^2.T^{-2}$}
{\True $M.L^2.T^{-3}$}
{$M.L.T^{-2}$}
{$M.L^{-1}.T^{-2}$}
\loigiai{
Công suất được xác định bởi:
 $P=\frac{A}{t}$
 Vì $[A]=M.L^2.T^{-2}$ nên:
 $[P]=\frac{M.L^2.T^{-2}}{T}=M.L^2.T^{-3}.$
}
\end{ex}

% ===== Câu 189 =====
% ID: L10C1B3-49-DS
% Mức: NB
\dangbt{Tính giá trị trung bình của đại lượng đo}
\begin{ex}
% ID: L10C1B3-49-DS
% Mức: NB
Một nhóm đo đại lượng $X$ ba lần: 19,0; 20,0; 21,0 (cùng đơn vị). Xét các phát biểu sau.
\choiceTF
{\True Giá trị trung bình là 20,0.}
{Tổng ba kết quả phải chia cho 2.}
{\True Giá trị trung bình nằm giữa giá trị nhỏ nhất và lớn nhất.}
{\True Nếu thêm một kết quả 20,0 thì giá trị trung bình không đổi.}
\loigiai{
\begin{itemchoice}
\item (Đúng) Giá trị trung bình là 20,0.
\item (Sai) Tổng ba kết quả phải chia cho 2.
\item (Đúng) Giá trị trung bình nằm giữa giá trị nhỏ nhất và lớn nhất.
\item (Đúng) Nếu thêm một kết quả 20,0 thì giá trị trung bình không đổi.
\end{itemchoice}
}
\end{ex}

% ===== Câu 190 =====
% ID: L10C1B3-49-TLN
% Mức: TH
\dangbt{Thứ nguyên}
\begin{ex}
% ID: L10C1B3-49-TLN
% Mức: TH
Lực có thứ nguyên [F] = M·L·T^n. Số mũ n bằng bao nhiêu?
\shortans{-2}
\loigiai{
Từ $F$ = ma, suy ra [F] = M·L·T^(-2), do đó n = -2.
}
\end{ex}

% ===== Câu 191 =====
% ID: L10C1B3-49-TN
% Mức: NB
\begin{ex}
% ID: L10C1B3-49-TN
% Mức: NB
Thứ nguyên của gia tốc là
\choice
{$L.T^{-1}$}
{\True $L.T^{-2}$}
{$M.L.T^{-2}$}
{$M.T^{-2}$}
\loigiai{
Gia tốc được xác định bởi:
 $a=\frac{\Delta v}{\Delta t}$
 Vì $[v]=L.T^{-1}$ nên:
 $[a]=\frac{L.T^{-1}}{T}=L.T^{-2}.$
}
\end{ex}

% ===== Câu 192 =====
% ID: L10C1B3-50-DS
% Mức: NB
\dangbt{Tính giá trị trung bình của đại lượng đo}
\begin{ex}
% ID: L10C1B3-50-DS
% Mức: NB
Một nhóm đo đại lượng $X$ ba lần: 20,0; 21,0; 22,0 (cùng đơn vị). Xét các phát biểu sau.
\choiceTF
{\True Giá trị trung bình là 21,0.}
{Tổng ba kết quả phải chia cho 2.}
{\True Giá trị trung bình nằm giữa giá trị nhỏ nhất và lớn nhất.}
{\True Nếu thêm một kết quả 21,0 thì giá trị trung bình không đổi.}
\loigiai{
\begin{itemchoice}
\item (Đúng) Giá trị trung bình là 21,0.
\item (Sai) Tổng ba kết quả phải chia cho 2.
\item (Đúng) Giá trị trung bình nằm giữa giá trị nhỏ nhất và lớn nhất.
\item (Đúng) Nếu thêm một kết quả 21,0 thì giá trị trung bình không đổi.
\end{itemchoice}
}
\end{ex}

% ===== Câu 193 =====
% ID: L10C1B3-50-TLN
% Mức: TH
\dangbt{Thứ nguyên}

% ===== Câu 194 =====
% ID: L10C1B3-50-TN
% Mức: NB
\begin{ex}
% ID: L10C1B3-50-TN
% Mức: NB
Thứ nguyên của công cơ học là
\choice
{$M.L.T^{-2}$}
{$M.L.T^{-1}$}
{\True $M.L^2.T^{-2}$}
{$M.L^2.T^{-3}$}
\loigiai{
Công được xác định bởi:
$A = F \cdot s$

Suy ra:

$[A] = [F][s] = M \cdot L \cdot T^{-2} \cdot L = M \cdot L^2 \cdot T^{-2}.$
}
\end{ex}

% ===== Câu 195 =====
% ID: L10C1B3-51-DS
% Mức: TH
\dangbt{Tính giá trị trung bình của đại lượng đo}
\begin{ex}
% ID: L10C1B3-51-DS
% Mức: TH
Một nhóm đo đại lượng $X$ ba lần: 21,0; 22,0; 23,0 (cùng đơn vị). Xét các phát biểu sau.
\choiceTF
{\True Giá trị trung bình là 22,0.}
{Tổng ba kết quả phải chia cho 2.}
{\True Giá trị trung bình nằm giữa giá trị nhỏ nhất và lớn nhất.}
{\True Nếu thêm một kết quả 22,0 thì giá trị trung bình không đổi.}
\loigiai{
\begin{itemchoice}
\item (Đúng) Giá trị trung bình là 22,0.
\item (Sai) Tổng ba kết quả phải chia cho 2.
\item (Đúng) Giá trị trung bình nằm giữa giá trị nhỏ nhất và lớn nhất.
\item (Đúng) Nếu thêm một kết quả 22,0 thì giá trị trung bình không đổi.
\end{itemchoice}
}
\end{ex}

% ===== Câu 196 =====
% ID: L10C1B3-51-TLN
% Mức: TH
\dangbt{Thứ nguyên}

% ===== Câu 197 =====
% ID: L10C1B3-51-TN
% Mức: NB
\begin{ex}
% ID: L10C1B3-51-TN
% Mức: NB
Thứ nguyên của khối lượng riêng là
\choice
{$M.L^3$}
{$M.L^{-2}$}
{\True $M.L^{-3}$}
{$M.L.T^{-2}$}
\loigiai{
Khối lượng riêng:
 $\rho=\frac{m}{V}$
 Do $[m]=M$ và $[V]=L^3$, nên:
 $[\rho]=\frac{M}{L^3}=M.L^{-3}.$
}
\end{ex}

% ===== Câu 198 =====
% ID: L10C1B3-52-DS
% Mức: TH
\dangbt{Tính giá trị trung bình của đại lượng đo}
\begin{ex}
% ID: L10C1B3-52-DS
% Mức: TH
Một nhóm đo đại lượng $X$ ba lần: 22,0; 23,0; 24,0 (cùng đơn vị). Xét các phát biểu sau.
\choiceTF
{\True Giá trị trung bình là 23,0.}
{Tổng ba kết quả phải chia cho 2.}
{\True Giá trị trung bình nằm giữa giá trị nhỏ nhất và lớn nhất.}
{\True Nếu thêm một kết quả 23,0 thì giá trị trung bình không đổi.}
\loigiai{
\begin{itemchoice}
\item (Đúng) Giá trị trung bình là 23,0.
\item (Sai) Tổng ba kết quả phải chia cho 2.
\item (Đúng) Giá trị trung bình nằm giữa giá trị nhỏ nhất và lớn nhất.
\item (Đúng) Nếu thêm một kết quả 23,0 thì giá trị trung bình không đổi.
\end{itemchoice}
}
\end{ex}

% ===== Câu 199 =====
% ID: L10C1B3-52-TLN
% Mức: TH
\dangbt{Thứ nguyên}

% ===== Câu 200 =====
% ID: L10C1B3-52-TN
% Mức: NB
\begin{ex}
% ID: L10C1B3-52-TN
% Mức: NB
Theo công thức lực đàn hồi $F=k\Delta l$, thứ nguyên của độ cứng lò xo $k$ là
\choice
{$M.L.T^{-2}$}
{\True $M.T^{-2}$}
{$M.L^2.T^{-2}$}
{$L.T^{-2}$}
\loigiai{
Từ:
 $F=k\Delta l$
 suy ra:
  [k]=\frac{[F]}{[\Delta l]}
 =\frac{M.L.T^{-2}}{L}=M.T^{-2}.
}
\end{ex}

% ===== Câu 201 =====
% ID: L10C1B3-53-DS
% Mức: VD
\dangbt{Tính giá trị trung bình của đại lượng đo}
\begin{ex}
% ID: L10C1B3-53-DS
% Mức: VD
Một nhóm đo đại lượng $X$ ba lần: 23,0; 24,0; 25,0 (cùng đơn vị). Xét các phát biểu sau.
\choiceTF
{\True Giá trị trung bình là 24,0.}
{Tổng ba kết quả phải chia cho 2.}
{\True Giá trị trung bình nằm giữa giá trị nhỏ nhất và lớn nhất.}
{\True Nếu thêm một kết quả 24,0 thì giá trị trung bình không đổi.}
\loigiai{
\begin{itemchoice}
\item (Đúng) Giá trị trung bình là 24,0.
\item (Sai) Tổng ba kết quả phải chia cho 2.
\item (Đúng) Giá trị trung bình nằm giữa giá trị nhỏ nhất và lớn nhất.
\item (Đúng) Nếu thêm một kết quả 24,0 thì giá trị trung bình không đổi.
\end{itemchoice}
}
\end{ex}

% ===== Câu 202 =====
% ID: L10C1B3-53-TLN
% Mức: VD
\dangbt{Thứ nguyên}

% ===== Câu 203 =====
% ID: L10C1B3-53-TN
% Mức: NB
\begin{ex}
% ID: L10C1B3-53-TN
% Mức: NB
Thứ nguyên của lực là
\choice
{$M.L.T^{-1}$}
{\True $M.L.T^{-2}$}
{$M.L^2.T^{-2}$}
{$M.L^{-1}.T^{-2}$}
\loigiai{
Theo định luật II Newton:
 $F=ma$
 Do đó:
 $[F]=[m][a]=M.L.T^{-2}.$
}
\end{ex}

% ===== Câu 204 =====
% ID: L10C1B3-54-DS
% Mức: NB
\dangbt{Viết số dưới dạng kí hiệu khoa học với số chữ số có nghĩa cho trước}
\begin{ex}
% ID: L10C1B3-54-DS
% Mức: NB
Xét các phát biểu về kí hiệu khoa học.
\choiceTF
{\True 3,20×10³ có ba chữ số có nghĩa.}
{\True 0,0045 = 4,5×10⁻³.}
{Hệ số trong kí hiệu khoa học phải từ 0 đến nhỏ hơn 1.}
{\True 1,00×10² biểu thị ba chữ số có nghĩa.}
\loigiai{
\begin{itemchoice}
\item (Đúng) 3,20×10³ có ba chữ số có nghĩa. (Vì): ; b đúng; c sai vì hệ số chuẩn thỏa 1 ≤ |a| < 10; d đúng vì các số 0 sau dấu phẩy có nghĩa
\item (Đúng) 0,0045 = 4,5×10⁻³.
\item (Sai) Hệ số trong kí hiệu khoa học phải từ 0 đến nhỏ hơn 1.
\item (Đúng) 1,00×10² biểu thị ba chữ số có nghĩa.
\end{itemchoice}
}
\end{ex}

% ===== Câu 205 =====
% ID: L10C1B3-54-TLN
% Mức: VD
\dangbt{Thứ nguyên}

% ===== Câu 206 =====
% ID: L10C1B3-54-TN
% Mức: NB
\begin{ex}
% ID: L10C1B3-54-TN
% Mức: NB
Động lượng của một vật được xác định bởi công thức $\vec{p}=m\vec{v}$. Thứ nguyên của động lượng là
\choice
{$M.L.T^{-2}$}
{\True $M.L.T^{-1}$}
{$M.L^2.T^{-2}$}
{$M.L^2.T^{-1}$}
\loigiai{
Động lượng:
 $p=mv$
 Suy ra:
 $[p]=[m][v]=M.L.T^{-1}.$
}
\end{ex}

% ===== Câu 207 =====
% ID: L10C1B3-55-DS
% Mức: NB
\dangbt{Viết số dưới dạng kí hiệu khoa học với số chữ số có nghĩa cho trước}
\begin{ex}
% ID: L10C1B3-55-DS
% Mức: NB
Xét các phát biểu về kí hiệu khoa học.
\choiceTF
{\True 3,20×10³ có ba chữ số có nghĩa.}
{\True 0,0045 = 4,5×10⁻³.}
{Hệ số trong kí hiệu khoa học phải từ 0 đến nhỏ hơn 1.}
{\True 1,00×10² biểu thị ba chữ số có nghĩa.}
\loigiai{
\begin{itemchoice}
\item (Đúng) 3,20×10³ có ba chữ số có nghĩa. (Vì): ; b đúng; c sai vì hệ số chuẩn thỏa 1 ≤ |a| < 10; d đúng vì các số 0 sau dấu phẩy có nghĩa
\item (Đúng) 0,0045 = 4,5×10⁻³.
\item (Sai) Hệ số trong kí hiệu khoa học phải từ 0 đến nhỏ hơn 1.
\item (Đúng) 1,00×10² biểu thị ba chữ số có nghĩa.
\end{itemchoice}
}
\end{ex}

% ===== Câu 208 =====
% ID: L10C1B3-55-TLN
% Mức: NB
\dangbt{Tính giá trị trung bình của đại lượng đo}
\begin{ex}
% ID: L10C1B3-55-TLN
% Mức: NB
Đo một chiều dài được $12\,\text{cm}$, $13\,\text{cm}$ và $14\,\text{cm}$. Tính giá trị trung bình theo đơn vị cm. Chỉ ghi số.
\shortans{13}
\loigiai{
Giá trị trung bình = (12+13+14)/3 = $13\,\text{cm}$.
}
\end{ex}

% ===== Câu 209 =====
% ID: L10C1B3-55-TN
% Mức: NB
\dangbt{Thứ nguyên}
\begin{ex}
% ID: L10C1B3-55-TN
% Mức: NB
Trong hệ đơn vị đo lường quốc tế (SI), đơn vị đo độ dài là:
\choice
{cm (xentimét)}
{\True m (mét)}
{km (kilômét)}
{mm (milimét)}
\loigiai{
Trong hệ SI (Hệ đơn vị quốc tế), mét (m) là đơn vị cơ bản để đo độ dài. Các đơn vị khác như cm, km, mm là các bội số hoặc ước số của mét, không phải đơn vị cơ bản của SI.
}
\end{ex}

% ===== Câu 210 =====
% ID: L10C1B3-56-DS
% Mức: TH
\dangbt{Viết số dưới dạng kí hiệu khoa học với số chữ số có nghĩa cho trước}
\begin{ex}
% ID: L10C1B3-56-DS
% Mức: TH
Xét các phát biểu về kí hiệu khoa học.
\choiceTF
{\True 3,20×10³ có ba chữ số có nghĩa.}
{\True 0,0045 = 4,5×10⁻³.}
{Hệ số trong kí hiệu khoa học phải từ 0 đến nhỏ hơn 1.}
{\True 1,00×10² biểu thị ba chữ số có nghĩa.}
\loigiai{
\begin{itemchoice}
\item (Đúng) 3,20×10³ có ba chữ số có nghĩa. (Vì): ; b đúng; c sai vì hệ số chuẩn thỏa 1 ≤ |a| < 10; d đúng vì các số 0 sau dấu phẩy có nghĩa
\item (Đúng) 0,0045 = 4,5×10⁻³.
\item (Sai) Hệ số trong kí hiệu khoa học phải từ 0 đến nhỏ hơn 1.
\item (Đúng) 1,00×10² biểu thị ba chữ số có nghĩa.
\end{itemchoice}
}
\end{ex}

% ===== Câu 211 =====
% ID: L10C1B3-56-TLN
% Mức: NB
\dangbt{Tính giá trị trung bình của đại lượng đo}
\begin{ex}
% ID: L10C1B3-56-TLN
% Mức: NB
Đo một chiều dài được $13\,\text{cm}$, $14\,\text{cm}$ và $15\,\text{cm}$. Tính giá trị trung bình theo đơn vị cm. Chỉ ghi số.
\shortans{14}
\loigiai{
Giá trị trung bình = (13+14+15)/3 = $14\,\text{cm}$.
}
\end{ex}

% ===== Câu 212 =====
% ID: L10C1B3-56-TN
% Mức: NB
\dangbt{Thứ nguyên}
\begin{ex}
% ID: L10C1B3-56-TN
% Mức: NB
Thứ nguyên của vận tốc là
\choice
{$L.T$}
{\True $L.T^{-1}$}
{$L^{-1}.T$}
{$L^{2}.T^{-1}$}
\loigiai{
Vận tốc được xác định bởi: $v = \frac{s}{t}$. Do đó thứ nguyên của vận tốc là $[v] = \frac{L}{T} = L.T^{-1}$.
}
\end{ex}

% ===== Câu 213 =====
% ID: L10C1B3-57-DS
% Mức: TH
\dangbt{Viết số dưới dạng kí hiệu khoa học với số chữ số có nghĩa cho trước}
\begin{ex}
% ID: L10C1B3-57-DS
% Mức: TH
Xét các phát biểu về kí hiệu khoa học.
\choiceTF
{\True 3,20×10³ có ba chữ số có nghĩa.}
{\True 0,0045 = 4,5×10⁻³.}
{Hệ số trong kí hiệu khoa học phải từ 0 đến nhỏ hơn 1.}
{\True 1,00×10² biểu thị ba chữ số có nghĩa.}
\loigiai{
\begin{itemchoice}
\item (Đúng) 3,20×10³ có ba chữ số có nghĩa. (Vì): ; b đúng; c sai vì hệ số chuẩn thỏa 1 ≤ |a| < 10; d đúng vì các số 0 sau dấu phẩy có nghĩa
\item (Đúng) 0,0045 = 4,5×10⁻³.
\item (Sai) Hệ số trong kí hiệu khoa học phải từ 0 đến nhỏ hơn 1.
\item (Đúng) 1,00×10² biểu thị ba chữ số có nghĩa.
\end{itemchoice}
}
\end{ex}

% ===== Câu 214 =====
% ID: L10C1B3-57-TLN
% Mức: TH
\dangbt{Tính giá trị trung bình của đại lượng đo}
\begin{ex}
% ID: L10C1B3-57-TLN
% Mức: TH
Đo một chiều dài được $14\,\text{cm}$, $15\,\text{cm}$ và $16\,\text{cm}$. Tính giá trị trung bình theo đơn vị cm. Chỉ ghi số.
\shortans{15}
\loigiai{
Giá trị trung bình = (14+15+16)/3 = $15\,\text{cm}$.
}
\end{ex}

% ===== Câu 215 =====
% ID: L10C1B3-57-TN
% Mức: TH
\dangbt{Thứ nguyên}

% ===== Câu 216 =====
% ID: L10C1B3-58-DS
% Mức: VD
\dangbt{Viết số dưới dạng kí hiệu khoa học với số chữ số có nghĩa cho trước}
\begin{ex}
% ID: L10C1B3-58-DS
% Mức: VD
Xét các phát biểu về kí hiệu khoa học.
\choiceTF
{\True 3,20×10³ có ba chữ số có nghĩa.}
{\True 0,0045 = 4,5×10⁻³.}
{Hệ số trong kí hiệu khoa học phải từ 0 đến nhỏ hơn 1.}
{\True 1,00×10² biểu thị ba chữ số có nghĩa.}
\loigiai{
\begin{itemchoice}
\item (Đúng) 3,20×10³ có ba chữ số có nghĩa. (Vì): ; b đúng; c sai vì hệ số chuẩn thỏa 1 ≤ |a| < 10; d đúng vì các số 0 sau dấu phẩy có nghĩa
\item (Đúng) 0,0045 = 4,5×10⁻³.
\item (Sai) Hệ số trong kí hiệu khoa học phải từ 0 đến nhỏ hơn 1.
\item (Đúng) 1,00×10² biểu thị ba chữ số có nghĩa.
\end{itemchoice}
}
\end{ex}

% ===== Câu 217 =====
% ID: L10C1B3-58-TLN
% Mức: TH
\dangbt{Tính giá trị trung bình của đại lượng đo}
\begin{ex}
% ID: L10C1B3-58-TLN
% Mức: TH
Đo một chiều dài được $15\,\text{cm}$, $16\,\text{cm}$ và $17\,\text{cm}$. Tính giá trị trung bình theo đơn vị cm. Chỉ ghi số.
\shortans{16}
\loigiai{
Giá trị trung bình = (15+16+17)/3 = $16\,\text{cm}$.
}
\end{ex}

% ===== Câu 218 =====
% ID: L10C1B3-58-TN
% Mức: TH
\dangbt{Thứ nguyên}

% ===== Câu 219 =====
% ID: L10C1B3-59-DS
% Mức: NB
\dangbt{Xác định chữ số có nghĩa trong giá trị đo}
\begin{ex}
% ID: L10C1B3-59-DS
% Mức: NB
Xét quy tắc xác định chữ số có nghĩa.
\choiceTF
{\True Các chữ số 0 đứng đầu số thập phân nhỏ hơn 1 không có nghĩa.}
{\True Số 2,50 có ba chữ số có nghĩa.}
{Số 0,0405 có hai chữ số có nghĩa.}
{\True Số 7,000 có bốn chữ số có nghĩa.}
\loigiai{
\begin{itemchoice}
\item (Đúng) Các chữ số 0 đứng đầu số thập phân nhỏ hơn 1 không có nghĩa. (Vì): ; b đúng; c sai vì 0,0405 có ba chữ số có nghĩa (4, 0, 5); d đúng
\item (Đúng) Số 2,50 có ba chữ số có nghĩa.
\item (Sai) Số 0,0405 có hai chữ số có nghĩa.
\item (Đúng) Số 7,000 có bốn chữ số có nghĩa.
\end{itemchoice}
}
\end{ex}

% ===== Câu 220 =====
% ID: L10C1B3-59-TLN
% Mức: TH
\dangbt{Tính giá trị trung bình của đại lượng đo}
\begin{ex}
% ID: L10C1B3-59-TLN
% Mức: TH
Ba lần đo chiều dài cho các giá trị $10,2\,\text{cm}$; $10,4\,\text{cm}$; $10,3\,\text{cm}$. Giá trị trung bình theo đơn vị cm bằng bao nhiêu?
\shortans{10,3}
\loigiai{
Giá trị trung bình là (10,2 + 10,4 + 10,3)/3 = $10,3\,\text{cm}$.
}
\end{ex}

% ===== Câu 221 =====
% ID: L10C1B3-59-TN
% Mức: TH
\dangbt{Thứ nguyên}

% ===== Câu 222 =====
% ID: L10C1B3-60-DS
% Mức: NB
\dangbt{Xác định chữ số có nghĩa trong giá trị đo}
\begin{ex}
% ID: L10C1B3-60-DS
% Mức: NB
Xét quy tắc xác định chữ số có nghĩa.
\choiceTF
{\True Các chữ số 0 đứng đầu số thập phân nhỏ hơn 1 không có nghĩa.}
{\True Số 2,50 có ba chữ số có nghĩa.}
{Số 0,0405 có hai chữ số có nghĩa.}
{\True Số 7,000 có bốn chữ số có nghĩa.}
\loigiai{
\begin{itemchoice}
\item (Đúng) Các chữ số 0 đứng đầu số thập phân nhỏ hơn 1 không có nghĩa. (Vì): ; b đúng; c sai vì 0,0405 có ba chữ số có nghĩa (4, 0, 5); d đúng
\item (Đúng) Số 2,50 có ba chữ số có nghĩa.
\item (Sai) Số 0,0405 có hai chữ số có nghĩa.
\item (Đúng) Số 7,000 có bốn chữ số có nghĩa.
\end{itemchoice}
}
\end{ex}

% ===== Câu 223 =====
% ID: L10C1B3-60-TLN
% Mức: TH
\dangbt{Tính giá trị trung bình của đại lượng đo}

% ===== Câu 224 =====
% ID: L10C1B3-60-TN
% Mức: TH
\dangbt{Thứ nguyên}

% ===== Câu 225 =====
% ID: L10C1B3-61-DS
% Mức: TH
\dangbt{Xác định chữ số có nghĩa trong giá trị đo}
\begin{ex}
% ID: L10C1B3-61-DS
% Mức: TH
Xét quy tắc xác định chữ số có nghĩa.
\choiceTF
{\True Các chữ số 0 đứng đầu số thập phân nhỏ hơn 1 không có nghĩa.}
{\True Số 2,50 có ba chữ số có nghĩa.}
{Số 0,0405 có hai chữ số có nghĩa.}
{\True Số 7,000 có bốn chữ số có nghĩa.}
\loigiai{
\begin{itemchoice}
\item (Đúng) Các chữ số 0 đứng đầu số thập phân nhỏ hơn 1 không có nghĩa. (Vì): ; b đúng; c sai vì 0,0405 có ba chữ số có nghĩa (4, 0, 5); d đúng
\item (Đúng) Số 2,50 có ba chữ số có nghĩa.
\item (Sai) Số 0,0405 có hai chữ số có nghĩa.
\item (Đúng) Số 7,000 có bốn chữ số có nghĩa.
\end{itemchoice}
}
\end{ex}

% ===== Câu 226 =====
% ID: L10C1B3-61-TLN
% Mức: TH
\dangbt{Tính giá trị trung bình của đại lượng đo}

% ===== Câu 227 =====
% ID: L10C1B3-61-TN
% Mức: TH
\dangbt{Thứ nguyên}
\begin{ex}
% ID: L10C1B3-61-TN
% Mức: TH
Trong hệ đơn vị đo lường quốc tế (SI), đơn vị nào là đơn vị dẫn xuất:
\choice
{Kilôgam (kg)}
{Ampe (A)}
{Giây (s)}
{\True Vôn (V)}
\loigiai{
Trong hệ SI có 7 đơn vị cơ bản: mét (m), kilôgam (kg), giây (s), ampe

A. , kelvin (K), mol (mol), candela
C. Đúng — Các đơn vị còn lại được suy ra từ các đơn vị cơ bản gọi là đơn vị dẫn xuất. Vôn (V) là đơn vị đo hiệu điện thế, được dẫn xuất từ các đơn vị cơ bản: $1\,V = 1\,kg \cdot m^2 \cdot s^{-3} \cdot A^{-1}$. Do đó Vôn (V) là đơn vị dẫn xuất, không phải đơn vị cơ bản
}
\end{ex}

% ===== Câu 228 =====
% ID: L10C1B3-62-DS
% Mức: TH
\dangbt{Xác định chữ số có nghĩa trong giá trị đo}
\begin{ex}
% ID: L10C1B3-62-DS
% Mức: TH
Xét quy tắc xác định chữ số có nghĩa.
\choiceTF
{\True Các chữ số 0 đứng đầu số thập phân nhỏ hơn 1 không có nghĩa.}
{\True Số 2,50 có ba chữ số có nghĩa.}
{Số 0,0405 có hai chữ số có nghĩa.}
{\True Số 7,000 có bốn chữ số có nghĩa.}
\loigiai{
\begin{itemchoice}
\item (Đúng) Các chữ số 0 đứng đầu số thập phân nhỏ hơn 1 không có nghĩa. (Vì): ; b đúng; c sai vì 0,0405 có ba chữ số có nghĩa (4, 0, 5); d đúng
\item (Đúng) Số 2,50 có ba chữ số có nghĩa.
\item (Sai) Số 0,0405 có hai chữ số có nghĩa.
\item (Đúng) Số 7,000 có bốn chữ số có nghĩa.
\end{itemchoice}
}
\end{ex}

% ===== Câu 229 =====
% ID: L10C1B3-62-TLN
% Mức: TH
\dangbt{Tính giá trị trung bình của đại lượng đo}

% ===== Câu 230 =====
% ID: L10C1B3-62-TN
% Mức: TH
\dangbt{Thứ nguyên}
\begin{ex}
% ID: L10C1B3-62-TN
% Mức: TH
Chọn đúng. Thứ nguyên của vận tốc là:
\choice
{$L.T^{-1}$}
{$L.T^{-2}$}
{$L.T$}
{\True $\frac{L}{T}$}
\loigiai{
Vận tốc được xác định bằng quãng đường chia cho thời gian, nên thứ nguyên của nó là $\frac{L}{T}$, với $L$ là độ dài và $T$ là thời gian.
}
\end{ex}

% ===== Câu 231 =====
% ID: L10C1B3-63-DS
% Mức: VD
\dangbt{Xác định chữ số có nghĩa trong giá trị đo}
\begin{ex}
% ID: L10C1B3-63-DS
% Mức: VD
Xét quy tắc xác định chữ số có nghĩa.
\choiceTF
{\True Các chữ số 0 đứng đầu số thập phân nhỏ hơn 1 không có nghĩa.}
{\True Số 2,50 có ba chữ số có nghĩa.}
{Số 0,0405 có hai chữ số có nghĩa.}
{\True Số 7,000 có bốn chữ số có nghĩa.}
\loigiai{
\begin{itemchoice}
\item (Đúng) Các chữ số 0 đứng đầu số thập phân nhỏ hơn 1 không có nghĩa. (Vì): ; b đúng; c sai vì 0,0405 có ba chữ số có nghĩa (4, 0, 5); d đúng
\item (Đúng) Số 2,50 có ba chữ số có nghĩa.
\item (Sai) Số 0,0405 có hai chữ số có nghĩa.
\item (Đúng) Số 7,000 có bốn chữ số có nghĩa.
\end{itemchoice}
}
\end{ex}

% ===== Câu 232 =====
% ID: L10C1B3-63-TLN
% Mức: TH
\dangbt{Tính giá trị trung bình của đại lượng đo}
\begin{ex}
% ID: L10C1B3-63-TLN
% Mức: TH
Bốn lần đo thời gian cho các giá trị $2,1\,\text{s}$; $2,0\,\text{s}$; $1,9\,\text{s}$; $2,0\,\text{s}$. Giá trị trung bình theo đơn vị s bằng bao nhiêu?
\shortans{2}
\loigiai{
Giá trị trung bình là (2,1 + 2,0 + 1,9 + 2,0)/4 = $2,0\,\text{s}$.
}
\end{ex}

% ===== Câu 233 =====
% ID: L10C1B3-63-TN
% Mức: VD
\dangbt{Thứ nguyên}

% ===== Câu 234 =====
% ID: L10C1B3-64-TLN
% Mức: TH
\dangbt{Tính giá trị trung bình của đại lượng đo}

% ===== Câu 235 =====
% ID: L10C1B3-64-TN
% Mức: VD
\dangbt{Thứ nguyên}

% ===== Câu 236 =====
% ID: L10C1B3-65-TLN
% Mức: TH
\dangbt{Tính giá trị trung bình của đại lượng đo}

% ===== Câu 237 =====
% ID: L10C1B3-65-TN
% Mức: VD
\dangbt{Thứ nguyên}

% ===== Câu 238 =====
% ID: L10C1B3-66-TLN
% Mức: TH
\dangbt{Tính giá trị trung bình của đại lượng đo}

% ===== Câu 239 =====
% ID: L10C1B3-66-TN
% Mức: NB
\begin{ex}
% ID: L10C1B3-66-TN
% Mức: NB
Đo chiều dài ba lần được $10,0\,\text{cm}$; $11,0\,\text{cm}$; $12,0\,\text{cm}$. Giá trị trung bình bằng bao nhiêu?
\choice
{$10\,\text{cm}$}
{\True $11\,\text{cm}$}
{$12\,\text{cm}$}
{$13\,\text{cm}$}
\loigiai{
Giá trị trung bình: (10+11+12)/3 = $11\,\text{cm}$. Chọn B.
}
\end{ex}

% ===== Câu 240 =====
% ID: L10C1B3-67-TLN
% Mức: TH
\begin{ex}
% ID: L10C1B3-67-TLN
% Mức: TH
Năm lần đo khối lượng cho các giá trị $50,1\,\text{g}$; $49,9\,\text{g}$; $50,0\,\text{g}$; $50,2\,\text{g}$; $49,8\,\text{g}$. Giá trị trung bình theo đơn vị g bằng bao nhiêu?
\shortans{50}
\loigiai{
Giá trị trung bình là (50,1 + 49,9 + 50,0 + 50,2 + 49,8)/5 = $50,0\,\text{g}$.
}
\end{ex}

% ===== Câu 241 =====
% ID: L10C1B3-67-TN
% Mức: NB
\begin{ex}
% ID: L10C1B3-67-TN
% Mức: NB
Đo chiều dài ba lần được $11,0\,\text{cm}$; $12,0\,\text{cm}$; $13,0\,\text{cm}$. Giá trị trung bình bằng bao nhiêu?
\choice
{$11\,\text{cm}$}
{\True $12\,\text{cm}$}
{$13\,\text{cm}$}
{$14\,\text{cm}$}
\loigiai{
Giá trị trung bình: (11+12+13)/3 = $12\,\text{cm}$. Chọn B.
}
\end{ex}

% ===== Câu 243 =====
% ID: L10C1B3-68-TN
% Mức: NB
\begin{ex}
% ID: L10C1B3-68-TN
% Mức: NB
Đo chiều dài ba lần được $12,0\,\text{cm}$; $13,0\,\text{cm}$; $14,0\,\text{cm}$. Giá trị trung bình bằng bao nhiêu?
\choice
{$12\,\text{cm}$}
{\True $13\,\text{cm}$}
{$14\,\text{cm}$}
{$15\,\text{cm}$}
\loigiai{
Giá trị trung bình: (12+13+14)/3 = $13\,\text{cm}$. Chọn B.
}
\end{ex}

% ===== Câu 245 =====
% ID: L10C1B3-69-TN
% Mức: NB
\begin{ex}
% ID: L10C1B3-69-TN
% Mức: NB
Một học sinh đo chiều dài một vật ba lần được các kết quả: $10,2\,\mathrm{cm}$; $10,4\,\mathrm{cm}$; $10,3\,\mathrm{cm}$. Giá trị trung bình của chiều dài là
\choice
{$10,2\,\mathrm{cm}$}
{$10,4\,\mathrm{cm}$}
{\True $10,3\,\mathrm{cm}$}
{$30,9\,\mathrm{cm}$}
\loigiai{
Giá trị trung bình:
$\overline{l}=\frac{10,2+10,4+10,3}{3}=10,3\ \mathrm{cm}.$
}
\end{ex}

% ===== Câu 247 =====
% ID: L10C1B3-70-TN
% Mức: TH
\begin{ex}
% ID: L10C1B3-70-TN
% Mức: TH
Đo chiều dài ba lần được $13,0\,\text{cm}$; $14,0\,\text{cm}$; $15,0\,\text{cm}$. Giá trị trung bình bằng bao nhiêu?
\choice
{$13\,\text{cm}$}
{\True $14\,\text{cm}$}
{$15\,\text{cm}$}
{$16\,\text{cm}$}
\loigiai{
Giá trị trung bình: (13+14+15)/3 = $14\,\text{cm}$. Chọn B.
}
\end{ex}

% ===== Câu 248 =====
% ID: L10C1B3-71-TLN
% Mức: VD
\begin{ex}
% ID: L10C1B3-71-TLN
% Mức: VD
Đo một chiều dài được $16\,\text{cm}$, $17\,\text{cm}$ và $18\,\text{cm}$. Tính giá trị trung bình theo đơn vị cm. Chỉ ghi số.
\shortans{17}
\loigiai{
Giá trị trung bình = (16+17+18)/3 = $17\,\text{cm}$.
}
\end{ex}

% ===== Câu 249 =====
% ID: L10C1B3-71-TN
% Mức: TH
\begin{ex}
% ID: L10C1B3-71-TN
% Mức: TH
Đo chiều dài ba lần được $14,0\,\text{cm}$; $15,0\,\text{cm}$; $16,0\,\text{cm}$. Giá trị trung bình bằng bao nhiêu?
\choice
{$14\,\text{cm}$}
{\True $15\,\text{cm}$}
{$16\,\text{cm}$}
{$17\,\text{cm}$}
\loigiai{
Giá trị trung bình: (14+15+16)/3 = $15\,\text{cm}$. Chọn B.
}
\end{ex}

% ===== Câu 250 =====
% ID: L10C1B3-72-TLN
% Mức: VD
\begin{ex}
% ID: L10C1B3-72-TLN
% Mức: VD
Đo một chiều dài được $17\,\text{cm}$, $18\,\text{cm}$ và $19\,\text{cm}$. Tính giá trị trung bình theo đơn vị cm. Chỉ ghi số.
\shortans{18}
\loigiai{
Giá trị trung bình = (17+18+19)/3 = $18\,\text{cm}$.
}
\end{ex}

% ===== Câu 251 =====
% ID: L10C1B3-72-TN
% Mức: TH
\begin{ex}
% ID: L10C1B3-72-TN
% Mức: TH
Đo chiều dài ba lần được $15,0\,\text{cm}$; $16,0\,\text{cm}$; $17,0\,\text{cm}$. Giá trị trung bình bằng bao nhiêu?
\choice
{$15\,\text{cm}$}
{\True $16\,\text{cm}$}
{$17\,\text{cm}$}
{$18\,\text{cm}$}
\loigiai{
Giá trị trung bình: (15+16+17)/3 = $16\,\text{cm}$. Chọn B.
}
\end{ex}

% ===== Câu 252 =====
% ID: L10C1B3-73-TLN
% Mức: NB
\dangbt{Viết số dưới dạng kí hiệu khoa học với số chữ số có nghĩa cho trước}
\begin{ex}
% ID: L10C1B3-73-TLN
% Mức: NB
Trong số 6,25×10^3, số mũ của 10 bằng bao nhiêu? Chỉ ghi số.
\shortans{3}
\loigiai{
Số đã ở dạng a×10^n nên số mũ cần tìm là 3.
}
\end{ex}

% ===== Câu 253 =====
% ID: L10C1B3-73-TN
% Mức: TH
\dangbt{Tính giá trị trung bình của đại lượng đo}
\begin{ex}
% ID: L10C1B3-73-TN
% Mức: TH
Đo chiều dài ba lần được $16,0\,\text{cm}$; $17,0\,\text{cm}$; $18,0\,\text{cm}$. Giá trị trung bình bằng bao nhiêu?
\choice
{$16\,\text{cm}$}
{\True $17\,\text{cm}$}
{$18\,\text{cm}$}
{$19\,\text{cm}$}
\loigiai{
Giá trị trung bình: (16+17+18)/3 = $17\,\text{cm}$. Chọn B.
}
\end{ex}

% ===== Câu 254 =====
% ID: L10C1B3-74-TLN
% Mức: NB
\dangbt{Viết số dưới dạng kí hiệu khoa học với số chữ số có nghĩa cho trước}
\begin{ex}
% ID: L10C1B3-74-TLN
% Mức: NB
Trong số 6,25×10^4, số mũ của 10 bằng bao nhiêu? Chỉ ghi số.
\shortans{4}
\loigiai{
Số đã ở dạng a×10^n nên số mũ cần tìm là 4.
}
\end{ex}

% ===== Câu 255 =====
% ID: L10C1B3-74-TN
% Mức: VD
\dangbt{Tính giá trị trung bình của đại lượng đo}
\begin{ex}
% ID: L10C1B3-74-TN
% Mức: VD
Đo chiều dài ba lần được $17,0\,\text{cm}$; $18,0\,\text{cm}$; $19,0\,\text{cm}$. Giá trị trung bình bằng bao nhiêu?
\choice
{$17\,\text{cm}$}
{\True $18\,\text{cm}$}
{$19\,\text{cm}$}
{$20\,\text{cm}$}
\loigiai{
Giá trị trung bình: (17+18+19)/3 = $18\,\text{cm}$. Chọn B.
}
\end{ex}

% ===== Câu 256 =====
% ID: L10C1B3-75-TLN
% Mức: TH
\dangbt{Viết số dưới dạng kí hiệu khoa học với số chữ số có nghĩa cho trước}
\begin{ex}
% ID: L10C1B3-75-TLN
% Mức: TH
Trong số 6,25×10^5, số mũ của 10 bằng bao nhiêu? Chỉ ghi số.
\shortans{5}
\loigiai{
Số đã ở dạng a×10^n nên số mũ cần tìm là 5.
}
\end{ex}

% ===== Câu 257 =====
% ID: L10C1B3-75-TN
% Mức: VD
\dangbt{Tính giá trị trung bình của đại lượng đo}
\begin{ex}
% ID: L10C1B3-75-TN
% Mức: VD
Đo chiều dài ba lần được $18,0\,\text{cm}$; $19,0\,\text{cm}$; $20,0\,\text{cm}$. Giá trị trung bình bằng bao nhiêu?
\choice
{$18\,\text{cm}$}
{\True $19\,\text{cm}$}
{$20\,\text{cm}$}
{$21\,\text{cm}$}
\loigiai{
Giá trị trung bình: (18+19+20)/3 = $19\,\text{cm}$. Chọn B.
}
\end{ex}

% ===== Câu 258 =====
% ID: L10C1B3-76-TLN
% Mức: TH
\dangbt{Viết số dưới dạng kí hiệu khoa học với số chữ số có nghĩa cho trước}
\begin{ex}
% ID: L10C1B3-76-TLN
% Mức: TH
Trong số 6,25×10^6, số mũ của 10 bằng bao nhiêu? Chỉ ghi số.
\shortans{6}
\loigiai{
Số đã ở dạng a×10^n nên số mũ cần tìm là 6.
}
\end{ex}

% ===== Câu 259 =====
% ID: L10C1B3-76-TN
% Mức: VD
\dangbt{Tính giá trị trung bình của đại lượng đo}
\begin{ex}
% ID: L10C1B3-76-TN
% Mức: VD
Đo chiều dài ba lần được $19,0\,\text{cm}$; $20,0\,\text{cm}$; $21,0\,\text{cm}$. Giá trị trung bình bằng bao nhiêu?
\choice
{$19\,\text{cm}$}
{\True $20\,\text{cm}$}
{$21\,\text{cm}$}
{$22\,\text{cm}$}
\loigiai{
Giá trị trung bình: (19+20+21)/3 = $20\,\text{cm}$. Chọn B.
}
\end{ex}

% ===== Câu 260 =====
% ID: L10C1B3-77-TLN
% Mức: TH
\dangbt{Viết số dưới dạng kí hiệu khoa học với số chữ số có nghĩa cho trước}
\begin{ex}
% ID: L10C1B3-77-TLN
% Mức: TH
Viết 0,004728 dưới dạng a × 10^(-3) với 3 chữ số có nghĩa. Hệ số a bằng bao nhiêu?
\shortans{4,73}
\loigiai{
Làm tròn đến 3 chữ số có nghĩa: 0,004728 = 4,73 × 10^(-3), nên a = 4,73.
}
\end{ex}

% ===== Câu 261 =====
% ID: L10C1B3-77-TN
% Mức: NB
\begin{ex}
% ID: L10C1B3-77-TN
% Mức: NB
Viết số 123400 dưới dạng kí hiệu khoa học với 4 chữ số có nghĩa.
\choice
{1,234×10^4}
{\True 1,234×10^5}
{12,34×10^3}
{0,123×10^6}
\loigiai{
Dời dấu phẩy 5 chữ số sang trái: 123400 = 1,234×10^5. Chọn B.
}
\end{ex}

% ===== Câu 262 =====
% ID: L10C1B3-78-TLN
% Mức: TH
\begin{ex}
% ID: L10C1B3-78-TLN
% Mức: TH
Viết 0,004728 dưới dạng a × 10^(-3) với 3 chữ số có nghĩa. Hệ số a bằng bao nhiêu?
\shortans{4,73}
\loigiai{
Làm tròn đến 3 chữ số có nghĩa: 0,004728 = 4,73 × 10^(-3), nên a = 4,73.
}
\end{ex}

% ===== Câu 263 =====
% ID: L10C1B3-78-TN
% Mức: NB
\begin{ex}
% ID: L10C1B3-78-TN
% Mức: NB
Viết số 123500 dưới dạng kí hiệu khoa học với 4 chữ số có nghĩa.
\choice
{1,235×10^4}
{\True 1,235×10^5}
{12,35×10^3}
{0,124×10^6}
\loigiai{
Dời dấu phẩy 5 chữ số sang trái: 123500 = 1,235×10^5. Chọn B.
}
\end{ex}

% ===== Câu 264 =====
% ID: L10C1B3-79-TLN
% Mức: TH
\begin{ex}
% ID: L10C1B3-79-TLN
% Mức: TH
Viết 0,004728 dưới dạng a × 10^(-3) với 3 chữ số có nghĩa. Hệ số a bằng bao nhiêu?
\shortans{4,73}
\loigiai{
Làm tròn đến 3 chữ số có nghĩa: 0,004728 = 4,73 × 10^(-3), nên a = 4,73.
}
\end{ex}

% ===== Câu 265 =====
% ID: L10C1B3-79-TN
% Mức: NB
\begin{ex}
% ID: L10C1B3-79-TN
% Mức: NB
Viết số 123600 dưới dạng kí hiệu khoa học với 4 chữ số có nghĩa.
\choice
{1,236×10^4}
{\True 1,236×10^5}
{12,36×10^3}
{0,124×10^6}
\loigiai{
Dời dấu phẩy 5 chữ số sang trái: 123600 = 1,236×10^5. Chọn B.
}
\end{ex}

% ===== Câu 266 =====
% ID: L10C1B3-80-TLN
% Mức: TH
\begin{ex}
% ID: L10C1B3-80-TLN
% Mức: TH
Viết 0,004728 dưới dạng a × 10^(-3) với 3 chữ số có nghĩa. Hệ số a bằng bao nhiêu?
\shortans{4,73}
\loigiai{
Làm tròn đến 3 chữ số có nghĩa: 0,004728 = 4,73 × 10^(-3), nên a = 4,73.
}
\end{ex}

% ===== Câu 267 =====
% ID: L10C1B3-80-TN
% Mức: TH
\begin{ex}
% ID: L10C1B3-80-TN
% Mức: TH
Viết số 123700 dưới dạng kí hiệu khoa học với 4 chữ số có nghĩa.
\choice
{1,237×10^4}
{\True 1,237×10^5}
{12,37×10^3}
{0,124×10^6}
\loigiai{
Dời dấu phẩy 5 chữ số sang trái: 123700 = 1,237×10^5. Chọn B.
}
\end{ex}

% ===== Câu 268 =====
% ID: L10C1B3-81-TLN
% Mức: TH
\begin{ex}
% ID: L10C1B3-81-TLN
% Mức: TH
Viết 58200 dưới dạng a × 10^4 với 3 chữ số có nghĩa. Hệ số a bằng bao nhiêu?
\shortans{5,82}
\loigiai{
Làm tròn đến 3 chữ số có nghĩa: 58200 = 5,82 × 10^4, nên a = 5,82.
}
\end{ex}

% ===== Câu 269 =====
% ID: L10C1B3-81-TN
% Mức: TH
\begin{ex}
% ID: L10C1B3-81-TN
% Mức: TH
Viết số 123800 dưới dạng kí hiệu khoa học với 4 chữ số có nghĩa.
\choice
{1,238×10^4}
{\True 1,238×10^5}
{12,38×10^3}
{0,124×10^6}
\loigiai{
Dời dấu phẩy 5 chữ số sang trái: 123800 = 1,238×10^5. Chọn B.
}
\end{ex}

% ===== Câu 270 =====
% ID: L10C1B3-82-TLN
% Mức: TH
\begin{ex}
% ID: L10C1B3-82-TLN
% Mức: TH
Viết 58200 dưới dạng a × 10^4 với 3 chữ số có nghĩa. Hệ số a bằng bao nhiêu?
\shortans{5,82}
\loigiai{
Làm tròn đến 3 chữ số có nghĩa: 58200 = 5,82 × 10^4, nên a = 5,82.
}
\end{ex}

% ===== Câu 271 =====
% ID: L10C1B3-82-TN
% Mức: TH
\begin{ex}
% ID: L10C1B3-82-TN
% Mức: TH
Viết số 123900 dưới dạng kí hiệu khoa học với 4 chữ số có nghĩa.
\choice
{1,239×10^4}
{\True 1,239×10^5}
{12,39×10^3}
{0,124×10^6}
\loigiai{
Dời dấu phẩy 5 chữ số sang trái: 123900 = 1,239×10^5. Chọn B.
}
\end{ex}

% ===== Câu 272 =====
% ID: L10C1B3-83-TLN
% Mức: TH
\begin{ex}
% ID: L10C1B3-83-TLN
% Mức: TH
Viết 58200 dưới dạng a × 10^4 với 3 chữ số có nghĩa. Hệ số a bằng bao nhiêu?
\shortans{5,82}
\loigiai{
Làm tròn đến 3 chữ số có nghĩa: 58200 = 5,82 × 10^4, nên a = 5,82.
}
\end{ex}

% ===== Câu 273 =====
% ID: L10C1B3-83-TN
% Mức: TH
\begin{ex}
% ID: L10C1B3-83-TN
% Mức: TH
Viết số 124000 dưới dạng kí hiệu khoa học với 4 chữ số có nghĩa.
\choice
{1,24×10^4}
{\True 1,24×10^5}
{12,4×10^3}
{0,124×10^6}
\loigiai{
Dời dấu phẩy 5 chữ số sang trái: 124000 = 1,24×10^5. Chọn B.
}
\end{ex}

% ===== Câu 274 =====
% ID: L10C1B3-84-TLN
% Mức: TH
\begin{ex}
% ID: L10C1B3-84-TLN
% Mức: TH
Viết 58200 dưới dạng a × 10^4 với 3 chữ số có nghĩa. Hệ số a bằng bao nhiêu?
\shortans{5,82}
\loigiai{
Làm tròn đến 3 chữ số có nghĩa: 58200 = 5,82 × 10^4, nên a = 5,82.
}
\end{ex}

% ===== Câu 275 =====
% ID: L10C1B3-84-TN
% Mức: TH
\begin{ex}
% ID: L10C1B3-84-TN
% Mức: TH
Viết số $0{,}000708$ dưới dạng kí hiệu khoa học và giữ nguyên $3$ chữ số có nghĩa.
\choice
{$7,8\times 10^{-4}$}
{\True $7,08\times 10^{-4}$}
{$70,8\times 10^{-5}$}
{$0,708\times 10^{-3}$}
\loigiai{
Ta có:
 $0,000708=7,08\times 10^{-4}.$
 Số $7,08$ có đúng $3$ chữ số có nghĩa.
}
\end{ex}

% ===== Câu 276 =====
% ID: L10C1B3-85-TLN
% Mức: TH
\begin{ex}
% ID: L10C1B3-85-TLN
% Mức: TH
Viết 0,0007084 dưới dạng a × 10^(-4) với 3 chữ số có nghĩa. Hệ số a bằng bao nhiêu?
\shortans{7,08}
\loigiai{
Làm tròn đến 3 chữ số có nghĩa: 0,0007084 = 7,08 × 10^(-4), nên a = 7,08.
}
\end{ex}

% ===== Câu 277 =====
% ID: L10C1B3-85-TN
% Mức: VD
\begin{ex}
% ID: L10C1B3-85-TN
% Mức: VD
Viết số 124100 dưới dạng kí hiệu khoa học với 4 chữ số có nghĩa.
\choice
{1,241×10^4}
{\True 1,241×10^5}
{12,41×10^3}
{0,124×10^6}
\loigiai{
Dời dấu phẩy 5 chữ số sang trái: 124100 = 1,241×10^5. Chọn B.
}
\end{ex}

% ===== Câu 278 =====
% ID: L10C1B3-86-TLN
% Mức: TH
\begin{ex}
% ID: L10C1B3-86-TLN
% Mức: TH
Viết 0,0007084 dưới dạng a × 10^(-4) với 3 chữ số có nghĩa. Hệ số a bằng bao nhiêu?
\shortans{7,08}
\loigiai{
Làm tròn đến 3 chữ số có nghĩa: 0,0007084 = 7,08 × 10^(-4), nên a = 7,08.
}
\end{ex}

% ===== Câu 279 =====
% ID: L10C1B3-86-TN
% Mức: VD
\begin{ex}
% ID: L10C1B3-86-TN
% Mức: VD
Viết số 124200 dưới dạng kí hiệu khoa học với 4 chữ số có nghĩa.
\choice
{1,242×10^4}
{\True 1,242×10^5}
{12,42×10^3}
{0,124×10^6}
\loigiai{
Dời dấu phẩy 5 chữ số sang trái: 124200 = 1,242×10^5. Chọn B.
}
\end{ex}

% ===== Câu 280 =====
% ID: L10C1B3-87-TLN
% Mức: TH
\begin{ex}
% ID: L10C1B3-87-TLN
% Mức: TH
Viết 0,0007084 dưới dạng a × 10^(-4) với 3 chữ số có nghĩa. Hệ số a bằng bao nhiêu?
\shortans{7,08}
\loigiai{
Làm tròn đến 3 chữ số có nghĩa: 0,0007084 = 7,08 × 10^(-4), nên a = 7,08.
}
\end{ex}

% ===== Câu 281 =====
% ID: L10C1B3-87-TN
% Mức: VD
\begin{ex}
% ID: L10C1B3-87-TN
% Mức: VD
Viết số 124300 dưới dạng kí hiệu khoa học với 4 chữ số có nghĩa.
\choice
{1,243×10^4}
{\True 1,243×10^5}
{12,43×10^3}
{0,124×10^6}
\loigiai{
Dời dấu phẩy 5 chữ số sang trái: 124300 = 1,243×10^5. Chọn B.
}
\end{ex}

% ===== Câu 282 =====
% ID: L10C1B3-88-TLN
% Mức: TH
\begin{ex}
% ID: L10C1B3-88-TLN
% Mức: TH
Viết 0,0007084 dưới dạng a × 10^(-4) với 3 chữ số có nghĩa. Hệ số a bằng bao nhiêu?
\shortans{7,08}
\loigiai{
Làm tròn đến 3 chữ số có nghĩa: 0,0007084 = 7,08 × 10^(-4), nên a = 7,08.
}
\end{ex}

% ===== Câu 283 =====
% ID: L10C1B3-88-TN
% Mức: NB
\dangbt{Xác định chữ số có nghĩa trong giá trị đo}
\begin{ex}
% ID: L10C1B3-88-TN
% Mức: NB
Giá trị đo $12,02\,\text{cm}$ có bao nhiêu chữ số có nghĩa?
\choice
{2}
{3}
{\True 4}
{5}
\loigiai{
Mọi chữ số khác 0 và các chữ số 0 nằm giữa hoặc ở cuối phần thập phân đều có nghĩa; 12,02 có 4 chữ số có nghĩa. Chọn C.
}
\end{ex}

% ===== Câu 284 =====
% ID: L10C1B3-89-TLN
% Mức: VD
\dangbt{Viết số dưới dạng kí hiệu khoa học với số chữ số có nghĩa cho trước}
\begin{ex}
% ID: L10C1B3-89-TLN
% Mức: VD
Trong số 6,25×10^7, số mũ của 10 bằng bao nhiêu? Chỉ ghi số.
\shortans{7}
\loigiai{
Số đã ở dạng a×10^n nên số mũ cần tìm là 7.
}
\end{ex}

% ===== Câu 285 =====
% ID: L10C1B3-89-TN
% Mức: NB
\dangbt{Xác định chữ số có nghĩa trong giá trị đo}
\begin{ex}
% ID: L10C1B3-89-TN
% Mức: NB
Giá trị đo $13,03\,\text{cm}$ có bao nhiêu chữ số có nghĩa?
\choice
{2}
{3}
{\True 4}
{5}
\loigiai{
Mọi chữ số khác 0 và các chữ số 0 nằm giữa hoặc ở cuối phần thập phân đều có nghĩa; 13,03 có 4 chữ số có nghĩa. Chọn C.
}
\end{ex}

% ===== Câu 286 =====
% ID: L10C1B3-90-TLN
% Mức: VD
\dangbt{Viết số dưới dạng kí hiệu khoa học với số chữ số có nghĩa cho trước}
\begin{ex}
% ID: L10C1B3-90-TLN
% Mức: VD
Trong số 6,25×10^8, số mũ của 10 bằng bao nhiêu? Chỉ ghi số.
\shortans{8}
\loigiai{
Số đã ở dạng a×10^n nên số mũ cần tìm là 8.
}
\end{ex}

% ===== Câu 287 =====
% ID: L10C1B3-90-TN
% Mức: NB
\dangbt{Xác định chữ số có nghĩa trong giá trị đo}
\begin{ex}
% ID: L10C1B3-90-TN
% Mức: NB
Giá trị đo $14,04\,\text{cm}$ có bao nhiêu chữ số có nghĩa?
\choice
{2}
{3}
{\True 4}
{5}
\loigiai{
Mọi chữ số khác 0 và các chữ số 0 nằm giữa hoặc ở cuối phần thập phân đều có nghĩa; 14,04 có 4 chữ số có nghĩa. Chọn C.
}
\end{ex}

% ===== Câu 288 =====
% ID: L10C1B3-91-TLN
% Mức: NB
\begin{ex}
% ID: L10C1B3-91-TLN
% Mức: NB
Số 0,0205 có bao nhiêu chữ số có nghĩa? Chỉ ghi số.
\shortans{3}
\loigiai{
Các số 0 đứng đầu không có nghĩa; các chữ số 2, 0, 5 tạo 3 chữ số có nghĩa.
}
\end{ex}

% ===== Câu 289 =====
% ID: L10C1B3-91-TN
% Mức: NB
\begin{ex}
% ID: L10C1B3-91-TN
% Mức: NB
Một học sinh đo chiều dài của một vật và ghi kết quả là $2,50\ \mathrm{m}$. Kết quả này có bao nhiêu chữ số có nghĩa?
\choice
{2 chữ số có nghĩa}
{\True 3 chữ số có nghĩa}
{4 chữ số có nghĩa}
{1 chữ số có nghĩa}
\loigiai{
Để xác định chữ số có nghĩa trong số $2{,}50$:
- Chữ số $2$ là chữ số khác $0$ → có nghĩa
- Chữ số $5$ là chữ số khác $0$ → có nghĩa
- Chữ số $0$ đứng sau dấu phẩy và sau chữ số khác $0$ → có nghĩa

Vì vậy, các chữ số có nghĩa là: $2$, $5$, $0$.

Kết quả có **3 chữ số có nghĩa**.
}
\end{ex}

% ===== Câu 290 =====
% ID: L10C1B3-92-TLN
% Mức: NB
\begin{ex}
% ID: L10C1B3-92-TLN
% Mức: NB
Số 0,0305 có bao nhiêu chữ số có nghĩa? Chỉ ghi số.
\shortans{3}
\loigiai{
Các số 0 đứng đầu không có nghĩa; các chữ số 3, 0, 5 tạo 3 chữ số có nghĩa.
}
\end{ex}

% ===== Câu 291 =====
% ID: L10C1B3-92-TN
% Mức: NB
\begin{ex}
% ID: L10C1B3-92-TN
% Mức: NB
Số $1,200\times 10^3$ có bao nhiêu chữ số có nghĩa?
\choice
{$2$}
{$3$}
{\True $4$}
{$5$}
\loigiai{
Trong số $1{,}200$, các chữ số $1$, $2$, $0$, $0$ đều là chữ số có nghĩa. 
 Vậy số đã cho có $4$ chữ số có nghĩa.
}
\end{ex}

% ===== Câu 292 =====
% ID: L10C1B3-93-TLN
% Mức: TH
\begin{ex}
% ID: L10C1B3-93-TLN
% Mức: TH
Số 0,0405 có bao nhiêu chữ số có nghĩa? Chỉ ghi số.
\shortans{3}
\loigiai{
Các số 0 đứng đầu không có nghĩa; các chữ số 4, 0, 5 tạo 3 chữ số có nghĩa.
}
\end{ex}

% ===== Câu 293 =====
% ID: L10C1B3-93-TN
% Mức: TH
\begin{ex}
% ID: L10C1B3-93-TN
% Mức: TH
Giá trị đo $15,02\,\text{cm}$ có bao nhiêu chữ số có nghĩa?
\choice
{2}
{3}
{\True 4}
{5}
\loigiai{
Mọi chữ số khác 0 và các chữ số 0 nằm giữa hoặc ở cuối phần thập phân đều có nghĩa; 15,02 có 4 chữ số có nghĩa. Chọn C.
}
\end{ex}

% ===== Câu 294 =====
% ID: L10C1B3-94-TLN
% Mức: TH
\begin{ex}
% ID: L10C1B3-94-TLN
% Mức: TH
Số 0,0505 có bao nhiêu chữ số có nghĩa? Chỉ ghi số.
\shortans{3}
\loigiai{
Các số 0 đứng đầu không có nghĩa; các chữ số 5, 0, 5 tạo 3 chữ số có nghĩa.
}
\end{ex}

% ===== Câu 295 =====
% ID: L10C1B3-94-TN
% Mức: TH
\begin{ex}
% ID: L10C1B3-94-TN
% Mức: TH
Giá trị đo $16,03\,\text{cm}$ có bao nhiêu chữ số có nghĩa?
\choice
{2}
{3}
{\True 4}
{5}
\loigiai{
Mọi chữ số khác 0 và các chữ số 0 nằm giữa hoặc ở cuối phần thập phân đều có nghĩa; 16,03 có 4 chữ số có nghĩa. Chọn C.
}
\end{ex}

% ===== Câu 296 =====
% ID: L10C1B3-95-TLN
% Mức: TH
\begin{ex}
% ID: L10C1B3-95-TLN
% Mức: TH
Giá trị đo $0,00450\,\text{m}$ có bao nhiêu chữ số có nghĩa?
\shortans{3}
\loigiai{
Các chữ số có nghĩa là 4, 5 và chữ số 0 ở cuối sau dấu thập phân, nên có 3 chữ số có nghĩa.
}
\end{ex}

% ===== Câu 297 =====
% ID: L10C1B3-95-TN
% Mức: TH
\begin{ex}
% ID: L10C1B3-95-TN
% Mức: TH
Giá trị đo $17,04\,\text{cm}$ có bao nhiêu chữ số có nghĩa?
\choice
{2}
{3}
{\True 4}
{5}
\loigiai{
Mọi chữ số khác 0 và các chữ số 0 nằm giữa hoặc ở cuối phần thập phân đều có nghĩa; 17,04 có 4 chữ số có nghĩa. Chọn C.
}
\end{ex}

% ===== Câu 298 =====
% ID: L10C1B3-96-TLN
% Mức: TH
\begin{ex}
% ID: L10C1B3-96-TLN
% Mức: TH
Giá trị đo $0,00450\,\text{m}$ có bao nhiêu chữ số có nghĩa?
\shortans{3}
\loigiai{
Các chữ số có nghĩa là 4, 5 và chữ số 0 ở cuối sau dấu thập phân, nên có 3 chữ số có nghĩa.
}
\end{ex}

% ===== Câu 299 =====
% ID: L10C1B3-96-TN
% Mức: TH
\begin{ex}
% ID: L10C1B3-96-TN
% Mức: TH
Giá trị đo $18,02\,\text{cm}$ có bao nhiêu chữ số có nghĩa?
\choice
{2}
{3}
{\True 4}
{5}
\loigiai{
Mọi chữ số khác 0 và các chữ số 0 nằm giữa hoặc ở cuối phần thập phân đều có nghĩa; 18,02 có 4 chữ số có nghĩa. Chọn C.
}
\end{ex}

% ===== Câu 300 =====
% ID: L10C1B3-97-TLN
% Mức: TH
\begin{ex}
% ID: L10C1B3-97-TLN
% Mức: TH
Giá trị đo $0,00450\,\text{m}$ có bao nhiêu chữ số có nghĩa?
\shortans{3}
\loigiai{
Các chữ số có nghĩa là 4, 5 và chữ số 0 ở cuối sau dấu thập phân, nên có 3 chữ số có nghĩa.
}
\end{ex}

% ===== Câu 301 =====
% ID: L10C1B3-97-TN
% Mức: TH
\begin{ex}
% ID: L10C1B3-97-TN
% Mức: TH
Làm tròn số $0{,}003456$ đến $2$ chữ số có nghĩa, ta được
\choice
{$0{,}0034$}
{\True $0{,}0035$}
{$0{,}00346$}
{$0{,}0040$}
\loigiai{
Hai chữ số có nghĩa đầu tiên là $3$ và $4$. Chữ số tiếp theo là $5$ nên làm tròn $4$ thành $5$.
 
 Vậy:
 $0,003456\approx 0,0035.$
}
\end{ex}

% ===== Câu 302 =====
% ID: L10C1B3-98-TLN
% Mức: TH
\begin{ex}
% ID: L10C1B3-98-TLN
% Mức: TH
Giá trị đo $0,00450\,\text{m}$ có bao nhiêu chữ số có nghĩa?
\shortans{3}
\loigiai{
Các chữ số có nghĩa là 4, 5 và chữ số 0 ở cuối sau dấu thập phân, nên có 3 chữ số có nghĩa.
}
\end{ex}

% ===== Câu 303 =====
% ID: L10C1B3-98-TN
% Mức: TH
\begin{ex}
% ID: L10C1B3-98-TN
% Mức: TH
Giá trị nào sau đây có 4 chữ số có nghĩa (CSCN)?
\choice
{13,1}
{13,1}
{0,0013}
{\True 13,1}
\loigiai{
13,10 có 4 chữ số có nghĩa vì tất cả các chữ số, bao gồm cả số 0 ở cuối, đều là các chữ số có nghĩa.
}
\end{ex}

% ===== Câu 304 =====
% ID: L10C1B3-99-TLN
% Mức: TH
\begin{ex}
% ID: L10C1B3-99-TLN
% Mức: TH
Giá trị đo $12,030\,\text{s}$ có bao nhiêu chữ số có nghĩa?
\shortans{5}
\loigiai{
Các chữ số 1, 2, 0, 3, 0 đều có nghĩa, nên có 5 chữ số có nghĩa.
}
\end{ex}

% ===== Câu 305 =====
% ID: L10C1B3-99-TN
% Mức: VD
\begin{ex}
% ID: L10C1B3-99-TN
% Mức: VD
Giá trị đo $19,03\,\text{cm}$ có bao nhiêu chữ số có nghĩa?
\choice
{2}
{3}
{\True 4}
{5}
\loigiai{
Mọi chữ số khác 0 và các chữ số 0 nằm giữa hoặc ở cuối phần thập phân đều có nghĩa; 19,03 có 4 chữ số có nghĩa. Chọn C.
}
\end{ex}
\dangbt{Phép đo trực tiếp và phép đo gián tiếp}
\begin{ex}
Phép đo của một đại lượng vật lí
\choice
{là những sai xót gặp phải khi đo một đại lượng vật lí.}
{là sai số gặp phải khi dụng cụ đo một đại lương vật lí.}
{\True là phép so sánh nó với một đại lượng cùng loại được quy ước làm đơn vị.}
{là những công cụ đo các đại lượng vật lý như thước, cân.}
\loigiai{Phép đo của một đại lượng vật lí là phép so sánh nó với một đại lượng cùng loại được quy ước làm đơn vị.}
\end{ex}

\dangbt{Phép đo trực tiếp và phép đo gián tiếp}
\begin{ex}
Để đo lực kéo tác dụng lên vật m, chỉ cần dùng dụng cụ đo là
\choice
{\True lực kế.}
{thước mét.}
{cân.}
{đồng hồ.}
\loigiai{Để đo lực kéo tác dụng lên vật m, chỉ cần dùng dụng cụ đo là lực kế.}
\end{ex}

\dangbt{Phép đo trực tiếp và phép đo gián tiếp}
\begin{ex}
Phép đo nào sau đây là phép đo gián tiếp?
\choice
{Đo chiều cao của học sinh trong lớp.}
{Đo cân nặng của học sinh trong lớp.}
{Đo thời gian đi từ nhà đến trường.}
{\True Đo vận tốc đi xe đạp từ nhà đến trường.}
\loigiai{Đo vận tốc đi xe đạp từ nhà đến trường là phép đo gián tiếp vì vận tốc được tính thông qua quãng đường và thời gian, mà quãng đường và thời gian là các đại lượng đo trực tiếp.}
\end{ex}

\dangbt{Phép đo trực tiếp và phép đo gián tiếp}
\begin{ex}
Phát biểu nào sau đây là sai khi nói về phép đo?
\choice
{Phép đo trực tiếp là phép so sánh trực tiếp qua dụng cụ đo.}
{Phép đo gián tiếp là phép đo thông qua công thức liên hệ với các đại lượng có thể đo trực tiếp.}
{\True Các đại lượng vật lý luôn có thể đo trực tiếp.}
{Phép đo gián tiếp là phép đo thông qua từ hai phép đo trực tiếp trở lên}
\loigiai{Không phải tất cả các đại lượng vật lý đều có thể đo trực tiếp. Một số đại lượng phải đo gián tiếp thông qua các đại lượng khác.}
\end{ex}

\dangbt{Phép đo trực tiếp và phép đo gián tiếp}
\begin{ex}
Phép so sánh trực tiếp nhờ dụng cụ đo gọi là
\choice
{phép đo gián tiếp.}
{giá trị trung bình.}
{\True phép đo trực tiếp.}
{dụng cụ đo trực tiếp.}
\loigiai{Phép so sánh trực tiếp nhờ dụng cụ đo gọi là phép đo trực tiếp.}
\end{ex}

\dangbt{Phép đo trực tiếp và phép đo gián tiếp}
\begin{ex}
Trong các phép đo sau, phép đo nào là phép đo gián tiếp?
\choice
{\True Đo tốc độ trung bình của một vật bằng thước đo chiều dài và đồng hồ bấm giây.}
{Đo chiều dài chiếc bút chì bằng thước đo chiều dài.}
{Đo khối lượng của 5 quả táo bằng cân đồng hồ.}
{Đo thời gian đi từ nhà đến trường bằng đồng hồ bấm giây.}
\loigiai{Đo tốc độ trung bình của một vật bằng thước đo chiều dài và đồng hồ bấm giây là phép đo gián tiếp vì tốc độ được tính từ chiều dài và thời gian.}
\end{ex}

\dangbt{Phép đo trực tiếp và phép đo gián tiếp}
\begin{ex}
Dùng thước kẹp đo đường kính viên bi là
\choice
{phép đo gián tiếp.}
{không thuộc phép đo nào.}
{\True phép đo trực tiếp.}
{không thể đo được.}
\loigiai{Dùng thước kẹp đo đường kính viên bi là phép đo trực tiếp vì giá trị được đọc trực tiếp trên dụng cụ đo.}
\end{ex}

\dangbt{Phép đo trực tiếp và phép đo gián tiếp}
\begin{ex}
Chỉ dùng thước đo chiều dài và đồng hồ bấm giây để đo tốc độ trung bình của một chiếc xe đồ chơi chuyển động thẳng từ điểm A đến điểm B. Nhận định nào sau đây là sai?
\choice
{Dùng đồng hồ bấm giây đo thời gian t là phép đo trực tiếp.}
{\True Có thể đo trực tiếp được tốc độ trung bình của chuyển động.}
{Dùng công thức $v = s/t$ tính tốc độ trung bình là phép đo gián tiếp.}
{Dùng thước đo quãng đường s là phép đo trực tiếp.}
\loigiai{Tốc độ trung bình không thể đo trực tiếp bằng thước và đồng hồ mà phải tính toán từ quãng đường và thời gian, nên đây là phép đo gián tiếp.}
\end{ex}

\dangbt{Phép đo trực tiếp và phép đo gián tiếp}
\begin{ex}
Giá trị của đại lượng cần đo được đọc trực tiếp trên dụng cụ đo gọi là
\choice
{sai số ngẫu nhiên.}
{\True phép đo trực tiếp.}
{phép đo gián tiếp.}
{dự đoán kết quả đo.}
\loigiai{Giá trị của đại lượng cần đo được đọc trực tiếp trên dụng cụ đo gọi là phép đo trực tiếp.}
\end{ex}

\dangbt{Phép đo trực tiếp và phép đo gián tiếp}
\begin{ex}
Trong các phép đo dưới đây, đâu là phép đo trực tiếp?
(1) Dùng thước đo chiều cao.
(2) Dùng cân đo cân nặng.
(3) Dùng cân và ca đong đo khối lượng riêng của nước.
(4) Dùng đồng hồ và cột cây số đo tốc độ của người lái xe.
\choice
{(2), (3), (4).}
{(2), (4).}
{\True (1), (2).}
{(1), (2), (4).}
\loigiai{Phép đo trực tiếp là phép đo mà giá trị của đại lượng cần đo được đọc trực tiếp trên dụng cụ đo.
(1) Dùng thước đo chiều cao: Trực tiếp.
(2) Dùng cân đo cân nặng: Trực tiếp.
(3) Dùng cân và ca đong đo khối lượng riêng của nước: Gián tiếp (khối lượng riêng = khối lượng / thể tích).
(4) Dùng đồng hồ và cột cây số đo tốc độ của người lái xe: Gián tiếp (tốc độ = quãng đường / thời gian).
Vậy các phép đo trực tiếp là (1) và (2).}
\end{ex}

\dangbt{Phép đo trực tiếp và phép đo gián tiếp}
\begin{ex}
Đâu là một phép đo gián tiếp?
\choice
{\True Phép đo thể tích của một cái hộp hình chữ nhật.}
{Phép đo chiều cao của một cái hộp hình chữ nhật.}
{Phép đo chiều rộng của một cái hộp hình chữ nhật.}
{Phép đo chiều dài của một cái hộp hình chữ nhật.}
\loigiai{Phép đo thể tích của một cái hộp hình chữ nhật là phép đo gián tiếp vì thể tích được tính bằng công thức $V = dài \times rộng \times cao$, trong đó chiều dài, chiều rộng và chiều cao là các đại lượng đo trực tiếp.}
\end{ex}

\dangbt{Phép đo trực tiếp và phép đo gián tiếp}
\begin{ex}
Trong các phép đo dưới đây, đâu là phép đo gián tiếp?
(1) Dùng thước đo quãng đường đi được.
(2) Dùng đồng hồ đo thời gian chuyển động.
(3) Đo gia tốc rơi tự do.
(4) Đo vận tốc trung bình của vật chuyển động.
\choice
{(1), (2).}
{(1), (2), (4).}
{(2), (3), (4).}
{\True (3), (4).}
\loigiai{Phép đo gián tiếp là phép đo thông qua công thức liên hệ với các đại lượng có thể đo trực tiếp.
(1) Dùng thước đo quãng đường đi được: Trực tiếp.
(2) Dùng đồng hồ đo thời gian chuyển động: Trực tiếp.
(3) Đo gia tốc rơi tự do: Gián tiếp (thường tính từ quãng đường và thời gian).
(4) Đo vận tốc trung bình của vật chuyển động: Gián tiếp (tốc độ = quãng đường / thời gian).
Vậy các phép đo gián tiếp là (3) và (4).}
\end{ex}

\dangbt{Phép đo trực tiếp và phép đo gián tiếp}
\begin{ex}
Phép đo của một đại lượng vật lý
\choice
{là những sai xót gặp phải khi đo một đại lượng vật lý.}
{là sai số gặp phải khi dụng cụ đo một đại lương vật lý.}
{\True là phép so sánh nó với một đại lượng cùng loại được quy ước làm đơn vị.}
{là những công cụ đo các đại lượng vật lý như thước, cân…vv.}
\loigiai{Phép đo của một đại lượng vật lý là phép so sánh nó với một đại lượng cùng loại được quy ước làm đơn vị.}
\end{ex}

\dangbt{Phép đo trực tiếp và phép đo gián tiếp}
\begin{ex}
Để xác định tốc độ trung bình của một người đi xe đạp chuyển động trên đoạn đường từ A đến B, ta cần dùng dụng cụ đo là
\choice
{tốc kế.}
{\True đồng hồ và thước mét.}
{chỉ cần đồng hồ.}
{chỉ cần thước.}
\loigiai{Để xác định tốc độ trung bình, ta cần đo quãng đường (bằng thước mét) và thời gian (bằng đồng hồ), sau đó tính toán $v = s/t$. Đây là phép đo gián tiếp.}
\end{ex}

\dangbt{Phép đo trực tiếp và phép đo gián tiếp}
\begin{ex}
Đo tốc độ tức thời là phép đo
\choice
{gián tiếp, dùng đồng hồ bấm giây.}
{\True trực tiếp, dùng súng bắn tốc độ.}
{trực tiếp, dùng thước đo và đồng hồ đo thời gian.}
{gián tiếp, dùng súng bắn tốc độ.}
\loigiai{Đo tốc độ tức thời bằng súng bắn tốc độ là phép đo trực tiếp vì giá trị tốc độ được đọc trực tiếp trên thiết bị.}
\end{ex}

\dangbt{Phép đo trực tiếp và phép đo gián tiếp}
\begin{ex}
Phép đo tốc độ trong bài “Thực hành đo tốc độ của vật chuyển động thẳng” (Sách giáo khoa Vật lí 10 – Bộ sách chân trời sáng tạo) là phép đo gián tiếp thông qua các phép đo trực tiếp
\choice
{\True quãng đường và thời gian.}
{thời gian và gia tốc.}
{gia tốc và quãng đường.}
{gia tốc và độ dịch chuyển.}
\loigiai{Tốc độ được tính bằng quãng đường chia thời gian, nên đây là phép đo gián tiếp thông qua phép đo trực tiếp quãng đường và thời gian.}
\end{ex}

\dangbt{Sai số}
\begin{ex}
Khi đo lực kéo tác dụng lên vật m, kết quả thu được là $12,845 \pm 0,095$ N thì
\choice
{\True sai số tuyệt đối của phép đo là 0,095 N.}
{kết quả chính xác của phép đo là 12,845 N.}
{sai số tương tối của phép đo là 0,095%.}
{giá trị trung bình của phép đo là 0,095 N.}
\loigiai{Cách viết kết quả đo là $A = \bar{A} \pm \Delta A$ với $\Delta A$ là sai số tuyệt đối của phép đo. Vậy nên sai số tuyệt đối của phép đo bài toán này là 0,095 N.}
\end{ex}

\dangbt{Sai số}
\begin{ex}
Một phép đo đại lượng vật lí A thu được giá trị trung bình là $\bar{A}$, sai số của phép đo là $\Delta A$. Cách ghi đúng kết quả đo A là
\choice
{A = $\bar{A} - \Delta A$.}
{A = $\bar{A} + \Delta A$.}
{\True A = $\bar{A} \pm \Delta A$.}
{A= A $\pm \Delta A$.}
\loigiai{Cách ghi đúng kết quả đo A là $A = \bar{A} \pm \Delta A$.}
\end{ex}

\dangbt{Sai số}
\begin{ex}
Chọn ý sai? Sai số ngẫu nhiên
\choice
{không có nguyên nhân rõ ràng.}
{\True là những sai sót mắc phải khi đo.}
{có thể do khả năng giác quan của con người dẫn đến thao tác đo không chuẩn.}
{chịu tác động của các yếu tố ngẫu nhiên bên ngoài.}
\loigiai{Sai số ngẫu nhiên là những sai lệch không có nguyên nhân rõ ràng, không phải là những sai sót mắc phải khi đo (đó là sai số hệ thống hoặc sai sót chủ quan).}
\end{ex}

\dangbt{Sai số}
\begin{ex}
Với $\overline{\Delta A_{ng}}$ là sai số ngẫu nhiên và $\Delta A'$ là sai số dụng cụ, sai số tuyệt đối của phép đo được xác định theo biểu thức
\choice
{$\Delta A = \overline{\Delta A_{ng}} - \Delta A'$.}
{\True $\Delta A = \overline{\Delta A_{ng}} + \Delta A'$.}
{$\Delta A = \Delta A'$.}
{$\Delta A = \overline{\Delta A_{ng}} \cdot \Delta A'$}
\loigiai{Sai số tuyệt đối của phép đo là tổng của sai số ngẫu nhiên trung bình và sai số dụng cụ: $\Delta A = \overline{\Delta A_{ng}} + \Delta A'$.}
\end{ex}

\dangbt{Sai số}
\begin{ex}
Khi đo n lần cùng một đại lượng A, ta nhận được các giá trị $A_1, A_2, ..., A_n$. Giá trị trung bình của A là $\bar{A}$. Sai số ngẫu nhiên tuyệt đối ứng với lần đo thứ n được tính bằng công thức
\choice
{$|\Delta A_n| = A_n - \bar{A}$}
{\True $|\Delta A_n| = |A_n - \bar{A}|$}
{$|\Delta A_n| = \bar{A} - A_n$}
{$|\Delta A_n| = A_n + \bar{A}$}
\loigiai{Sai số ngẫu nhiên tuyệt đối ứng với lần đo thứ n được tính bằng công thức $|\Delta A_n| = |A_n - \bar{A}|$.}
\end{ex}

\dangbt{Sai số}
\begin{ex}
Sai số tỉ đối của phép đo là
\choice
{tỉ số giữa sai số ngẫu nhiên và sai số tuyệt đối.}
{tỉ số giữa sai ngẫu nhiên và sai số hệ thống.}
{tỉ số giữa sai số tuyệt đối và sai số ngẫu nhiên.}
{\True tỉ số giữa sai số tuyệt đối và giá trị trung bình của đại lượng cần đo.}
\loigiai{Sai số tỉ đối của phép đo là tỉ số giữa sai số tuyệt đối và giá trị trung bình của đại lượng cần đo.}
\end{ex}

\dangbt{Sai số}
\begin{ex}
Sai số phép đo phân thành mấy loại?
\choice
{1.}
{\True 2.}
{3.}
{4.}
\loigiai{Sai số phép đo phân thành 2 loại chính: sai số hệ thống và sai số ngẫu nhiên.}
\end{ex}

\dangbt{Sai số}
\begin{ex}
Một học sinh sử dụng Vôn kế để đo hiệu điện thế, tuy nhiên chưa hiệu chỉnh kim của Vôn kế về vạch số 0 dẫn đến phép đo gặp sai số. Loại sai số này gọi là
\choice
{sai số tuyệt đối.}
{\True sai số hệ thống.}
{sai số tương đối.}
{sai số ngẫu nhiên.}
\loigiai{Sai số do chưa hiệu chỉnh kim của Vôn kế về vạch số 0 là sai số hệ thống, vì nó có nguyên nhân rõ ràng và có thể được loại bỏ hoặc giảm thiểu bằng cách hiệu chỉnh dụng cụ.}
\end{ex}

\dangbt{Sai số}
\begin{ex}
Sai số tuyệt đối của phép đo cho biết phạm vi biến thiên của giá trị đo được và bằng tổng của . và sai số dụng cụ $\Delta x = \overline{\Delta x_{ng}} + \Delta x_{dc}$. Trong đó sai số dụng cụ $\Delta x_{dc}$ thường được xem có giá trị bằng một nữa độ chia nhỏ nhất với những dụng cụ đơn giản như thước kẻ, cân bàn, bình chia độ.
\choice
{\True sai số ngẫu nhiên.}
{sai số hệ thống.}
{sai số tỉ đối.}
{giá trị trung bình của đại lượng cần đo.}
\loigiai{Sai số tuyệt đối của phép đo bằng tổng của sai số ngẫu nhiên trung bình và sai số dụng cụ.}
\end{ex}

\dangbt{Sai số}
\begin{ex}
Một số phương án hạn chế sai số khi thực hiện phép đo sau
(1) Thường xuyên hiệu chỉnh dụng cụ đo.
(2) Sử dụng thiết bị đo có độ chính xác cao.
(3) Thực hiện phép đo đúng kĩ thuật, đúng qui trình.
(4) Thực hiện đo nhiều lần và lấy giá trị trung bình.
Phương án đúng nhất là
\choice
{(1), (2), (3).}
{(1), (2), (4).}
{(2), (3), (4).}
{\True (1), (2), (3), (4).}
\loigiai{Tất cả các phương án (1), (2), (3), (4) đều là các cách hiệu quả để hạn chế sai số trong phép đo.}
\end{ex}

\dangbt{Sai số}
\begin{ex}
Tỉ số giữa sai số tuyệt đối và giá trị trung bình của đại lượng cần đo là sai số
\choice
{ngẫu nhiên.}
{hệ thống.}
{\True tương đối.}
{tuyệt đối.}
\loigiai{Tỉ số giữa sai số tuyệt đối và giá trị trung bình của đại lượng cần đo là sai số tương đối (hay sai số tỉ đối).}
\end{ex}

\dangbt{Sai số}
\begin{ex}
Biểu thức nào sau đây là sai khi ghi kết quả phép đo và sai số phép đo?
\choice
{$A = \bar{A} \pm \Delta A$}
{$\Delta A = \overline{\Delta A_{ng}} + \Delta A_{dc}$}
{$\delta A = \frac{\Delta A}{\bar{A}}$}
{\True $\Delta A = \bar{A} - \Delta A_{dc}$}
\loigiai{Biểu thức $\Delta A = \bar{A} - \Delta A_{dc}$ là sai. Sai số tuyệt đối là tổng của sai số ngẫu nhiên trung bình và sai số dụng cụ, không phải là hiệu của giá trị trung bình và sai số dụng cụ.}
\end{ex}

\dangbt{Sai số}
\begin{ex}
Một học sinh sử dụng thước thẳng có độ chia nhỏ nhất là 1 mm để đo chiều dài của một đoạn thẳng. Sai số dụng cụ của phép đo bằng thước trên thường được lấy giá trị bằng
\choice
{5 mm.}
{\True 0,5 mm.}
{1 mm.}
{0,1 mm.}
\loigiai{Sai số dụng cụ $\Delta A_{dc}$ thường được lấy bằng một nửa độ chia nhỏ nhất của dụng cụ. Với độ chia nhỏ nhất là 1 mm, sai số dụng cụ là $1/2 \times 1 \text{ mm} = 0,5 \text{ mm}$.}
\end{ex}

\dangbt{Sai số}
\begin{ex}
Khi sử dụng Vôn kế có nhiều thang đo, để đảm bảo an toàn cho Vôn kế, trước tiên người ta phải điều chỉnh Vôn kế dùng ở thang đo
\choice
{gần bằng với giá trị đo.}
{đúng bằng với giá trị đo.}
{\True cao nhất.}
{thấp nhất.}
\loigiai{Để đảm bảo an toàn cho Vôn kế và tránh làm hỏng thiết bị, trước tiên người ta phải điều chỉnh Vôn kế dùng ở thang đo cao nhất, sau đó giảm dần nếu cần.}
\end{ex}

\dangbt{Sai số}
\begin{ex}
Chọn cách viết sai kết quả của phép đo?
\choice
{$A = 12,3 \pm 0,1 \text{ cm}$}
{$A = 12,3 \pm 0,12 \text{ cm}$}
{\True $A = 12,3 \pm 0,123 \text{ cm}$}
{$A = 12,30 \pm 0,12 \text{ cm}$}
\loigiai{Khi ghi kết quả phép đo, số chữ số thập phân của sai số tuyệt đối phải tương ứng với số chữ số thập phân của giá trị trung bình. Trong phương án C, sai số tuyệt đối có 3 chữ số thập phân trong khi giá trị trung bình có 1 chữ số thập phân, điều này là không phù hợp.}
\end{ex}

\dangbt{Sai số}
\begin{ex}
Kết quả đo gia tốc rơi tự do được viết dưới dạng $g = 9,81 \pm 0,02 \text{ m/s}^2$. Sai số tỉ đối của phép đo xấp xỉ là
\choice
{0,20%.}
{0,203%.}
{\True 0,204%.}
{0,205%.}
\loigiai{Cách viết kết quả đo là $A = \bar{A} \pm \Delta A$. Sai số tỉ đối của phép đo là $\delta A = \frac{\Delta A}{\bar{A}} = \frac{0,02}{9,81} \approx 0,0020387 \approx 0,204\%$.}
\end{ex}

\dangbt{Sai số}
\begin{ex}
Khi đo gia tốc rơi tự do, một học sinh tính được $\bar{g} = 9,786 \text{ (m/s)}$ và $\Delta g = 0,0259 \text{ (m/s)}$. Sai số tỉ đối của phép đo là
\choice
{0,59%.}
{\True 0,265%.}
{2,65%.}
{2%.}
\loigiai{Sai số tỉ đối được xác định bằng công thức $\delta g = \frac{\Delta g}{\bar{g}} = \frac{0,0259}{9,786} \approx 0,0026466 \approx 0,265\%$.}
\end{ex}

\dangbt{Sai số}
\begin{ex}
Một người đo chiều cao một người khác trong ba lần được là 160 cm, 161 cm, 162 cm. Sai số tỉ đối của phép đo gần đúng là
\choice
{72%.}
{7,2%.}
{0%.}
{\True 0,72%.}
\loigiai{Chiều cao trung bình $\bar{h} = \frac{160 + 161 + 162}{3} = 161 \text{ cm}$.
Sai số tuyệt đối ứng với mỗi lần đo:
$|\Delta h_1| = |160 - 161| = 1 \text{ cm}$
$|\Delta h_2| = |161 - 161| = 0 \text{ cm}$
$|\Delta h_3| = |162 - 161| = 1 \text{ cm}$
Sai số tuyệt đối trung bình là $\overline{\Delta h_{ng}} = \frac{1 + 0 + 1}{3} = \frac{2}{3} \approx 0,67 \text{ cm}$.
Giả sử sai số dụng cụ $\Delta h_{dc}$ là 0 (hoặc rất nhỏ).
Sai số tuyệt đối của phép đo $\Delta h = \overline{\Delta h_{ng}} + \Delta h_{dc} \approx 0,67 \text{ cm}$.
Sai số tỉ đối của phép đo là $\delta h = \frac{\Delta h}{\bar{h}} = \frac{0,67}{161} \approx 0,00416 \approx 0,416\%$.
(Lưu ý: Nếu sai số dụng cụ được lấy là 0.5 độ chia nhỏ nhất, và độ chia nhỏ nhất của thước là 1cm, thì $\Delta h_{dc} = 0.5$ cm. Khi đó $\Delta h = 0.67 + 0.5 = 1.17$ cm. $\delta h = 1.17/161 \approx 0.00726 \approx 0.726\%$. Đáp án D là 0.72% nên có thể giả định sai số dụng cụ được tính vào.)}
\end{ex}

\dangbt{Sai số}
\begin{ex}
Một em học sinh đo (3 lần) chiều dài một quyển sách được các kết quả là 22,20 cm, 22,31 cm, 22,62 cm. Giá trị sai số ngẫu nhiên tuyệt đối trung bình của 3 lần đo là
\choice
{22,49 cm.}
{0,49 cm.}
{\True 0,16 cm.}
{22,41 cm.}
\loigiai{Chiều dài trung bình $\bar{L} = \frac{22,20 + 22,31 + 22,62}{3} = \frac{67,13}{3} \approx 22,3767 \text{ cm}$.
Sai số tuyệt đối ứng với mỗi lần đo:
$|\Delta L_1| = |22,20 - 22,3767| \approx 0,1767 \text{ cm}$
$|\Delta L_2| = |22,31 - 22,3767| \approx 0,0667 \text{ cm}$
$|\Delta L_3| = |22,62 - 22,3767| \approx 0,2433 \text{ cm}$
Sai số ngẫu nhiên tuyệt đối trung bình là $\overline{\Delta L_{ng}} = \frac{0,1767 + 0,0667 + 0,2433}{3} = \frac{0,4867}{3} \approx 0,1622 \text{ cm}$. Làm tròn đến hai chữ số thập phân ta được 0,16 cm.}
\end{ex}

\dangbt{Sai số}
\begin{ex}
Phép đo độ dài đường đi cho giá trị trung bình $\bar{s} = 12,34 \text{ m}$. Sai số của phép đo tính được là $\Delta s = 0,12 \text{ m}$. Kết quả đo được viết là
\choice
{$s = (12,34 - 0,12) \text{ m}$.}
{$s = (12,34 + 0,12) \text{ m}$.}
{\True $s = (12,34 \pm 0,12) \text{ m}$.}
{$s = 12,34 \pm 0,12 \text{ m}$.}
\loigiai{Cách viết kết quả đo là $A = \bar{A} \pm \Delta A$. Vậy nên $s = (12,34 \pm 0,12) \text{ m}$.}
\end{ex}

\dangbt{Sai số}
\begin{ex}
Kết quả đo tốc độ của một vật chuyển động ghi $(5,2 \pm 0,1)$ m/s. Sai số tỉ đối của phép đo bằng
\choice
{\True 1,9%.}
{10,0%.}
{9,8%.}
{52,0%.}
\loigiai{Sai số tỉ đối của phép đo là $\delta v = \frac{\Delta v}{\bar{v}} = \frac{0,1}{5,2} \approx 0,01923 \approx 1,9\%$.}
\end{ex}

\dangbt{Sai số}
\begin{ex}
Khi đo quãng đường di chuyển của vật m, kết quả thu được là $s = 125,856 \pm 1,546 \text{ cm}$. Sai số tương đối của phép đo này là
\choice
{1,546%.}
{\True 1,228%.}
{0,012%.}
{1,213%.}
\loigiai{Cách viết kết quả đo là $A = \bar{A} \pm \Delta A$ với $\Delta A$ là sai số tuyệt đối của phép đo và $\bar{A}$ là giá trị trung bình.
Sai số tương đối được xác định bằng công thức $\delta A = \frac{\Delta A}{\bar{A}}$.
Vậy nên sai số tương đối của bài toán này là $\frac{1,546}{125,856} \approx 0,012284 \approx 1,228\%$.}
\end{ex}

\dangbt{Sai số}
\begin{ex}
Khi đo hiệu điện thế giữa hai đầu điện trở R, kết quả thu được là $U = 12,50 \pm 3,50 \text{ V}$. Sai số tuyệt đối của phép đo này là
\choice
{3,57 V.}
{12,5 V.}
{0,44 V.}
{\True 3,50 V.}
\loigiai{Cách viết kết quả đo là $A = \bar{A} \pm \Delta A$ với $\Delta A$ là sai số tuyệt đối của phép đo và $\bar{A}$ là giá trị trung bình. Vậy sai số tuyệt đối là 3,50 V.}
\end{ex}

\dangbt{Sai số}
\begin{ex}
Trong thí nghiệm tính vận tốc của vật chuyển động thẳng đều, kết quả đo quãng đường $s = 8,255 \pm 0,245 \text{ m}$ và thời gian $t = 4,025 \pm 0,120 \text{ s}$. Kết quả của phép tính vận tốc là
\choice
{\True $2,051 \pm 0,122 \text{ m/s}$}
{$2,510 \pm 0,122 \text{ m/s}$}
{$2,051 \pm 0,242 \text{ m/s}$}
{$2,510 \pm 0,242 \text{ m/s}$}
\loigiai{Cách viết kết quả đo là $A = \bar{A} \pm \Delta A$ với $\Delta A$ là sai số tuyệt đối của phép đo và $\bar{A}$ là giá trị trung bình.
Giá trị trung bình của vận tốc $\bar{v} = \frac{\bar{s}}{\bar{t}} = \frac{8,255}{4,025} \approx 2,0510 \text{ m/s}$.
Sai số tỉ đối được xác định $\delta v = \delta s + \delta t = \frac{\Delta s}{\bar{s}} + \frac{\Delta t}{\bar{t}} = \frac{0,245}{8,255} + \frac{0,120}{4,025} \approx 0,02968 + 0,02981 \approx 0,05949$.
Sai số tuyệt đối của vận tốc $\Delta v = \bar{v} \cdot \delta v = 2,0510 \cdot 0,05949 \approx 0,12199 \text{ m/s}$.
Vậy nên $v = 2,051 \pm 0,122 \text{ m/s}$.}
\end{ex}

\dangbt{Sai số}
\begin{ex}
Trong thí nghiệm tính công của lực F, kết quả đo lực $F = 15,045 \pm 0,625 \text{ N}$ và quãng đường $s = 5,125 \pm 0,095 \text{ m}$. Kết quả của phép tính công là
\choice
{\True $77,106 \pm 6,008 \text{ J}$.}
{$2,936 \pm 6,008 \text{ J}$.}
{$77,106 \pm 6,008 \text{ J}$.}
{$77,106 \pm 4,632\% \text{ J}$.}
\loigiai{Cách viết kết quả đo là $A = \bar{A} \pm \Delta A$ với $\Delta A$ là sai số tuyệt đối của phép đo và $\bar{A}$ là giá trị trung bình.
Ta có $\bar{A} = \bar{F} \cdot \bar{s} = 15,045 \cdot 5,125 \approx 77,105625 \text{ J}$.
Sai số tỉ đối được xác định $\delta A = \delta F + \delta s = \frac{\Delta F}{\bar{F}} + \frac{\Delta s}{\bar{s}} = \frac{0,625}{15,045} + \frac{0,095}{5,125} \approx 0,04154 + 0,01853 \approx 0,06007$.
Sai số tuyệt đối $\Delta A = \bar{A} \cdot \delta A = 77,105625 \cdot 0,06007 \approx 4,6317 \text{ J}$.
Vậy nên $A = 77,106 \pm 4,632 \text{ J}$. (Làm tròn đến 3 chữ số thập phân cho giá trị trung bình và 3 chữ số thập phân cho sai số tuyệt đối).
Trong các đáp án, đáp án A có sai số tuyệt đối là 6,008 J, có thể do làm tròn khác hoặc có sai số dụng cụ khác. Tuy nhiên, nếu tính theo công thức thì kết quả gần nhất là $77,106 \pm 4,632 \text{ J}$. Nếu chấp nhận làm tròn sai số tuyệt đối lên 6,008 J thì đáp án A là hợp lý nhất.}
\end{ex}

\dangbt{Tính giá trị trung bình của đại lượng đo}
\begin{ex}
Một học sinh làm thí nghiệm đo gia tốc rơi tự do. Năm lần đo cho kết quả lần lượt là 9,7805 m/s2, 9,7807 m/s2, 9,7806 m/s2, 9,7806 m/s2, 9,7807 m/s2. Gia tốc rơi tự do trung bình của năm lần đo này bằng
\choice
{\True 9,7806 m/s2.}
{9,7840 m/s2.}
{9,7802 m/s2.}
{9,7860 m/s2.}
\loigiai{Gia tốc rơi tự do trung bình là $\bar{g} = \frac{9,7805 + 9,7807 + 9,7806 + 9,7806 + 9,7807}{5} = \frac{48,9031}{5} = 9,78062 \text{ m/s}^2$. Làm tròn đến 4 chữ số thập phân ta được 9,7806 m/s2.}
\end{ex}

\dangbt{Tính giá trị trung bình của đại lượng đo}
\begin{ex}
Đo chiều dày của một cuốn sách, được kết quả 2,3 cm, 2,4 cm, 2,5 cm, 2,4 cm. Giá trị trung bình của các lần đo là
\choice
{2,3 cm.}
{\True 2,4 cm}
{2,5 cm}
{2,2 cm.}
\loigiai{Giá trị trung bình $\bar{L} = \frac{2,3 + 2,4 + 2,5 + 2,4}{4} = \frac{9,6}{4} = 2,4 \text{ cm}$.}
\end{ex}

\dangbt{Thứ nguyên}
\begin{ex}
Đơn vị đo khối lượng trong hệ thống đo lường SI là
\choice
{tấn.}
{gam.}
{\True kilôgam.}
{miligam.}
\loigiai{Đơn vị đo khối lượng trong hệ thống đo lường SI là kilôgam (kg).}
\end{ex}

\dangbt{Thứ nguyên}
\begin{ex}
Đại lượng không phải là đại lượng cơ bản của hệ SI là
\choice
{khối lượng.}
{thời gian.}
{quãng đường.}
{\True vận tốc.}
\loigiai{Các đại lượng cơ bản của hệ SI là chiều dài (mét), khối lượng (kilôgam), thời gian (giây), cường độ dòng điện (ampere), nhiệt độ nhiệt động lực học (kelvin), lượng chất (mol), cường độ sáng (candela). Vận tốc là đại lượng dẫn xuất.}
\end{ex}

\dangbt{Thứ nguyên}
\begin{ex}
Dụng cụ nào sau đây dùng để đo vận tốc tức thời?
\choice
{Tần số kế.}
{Nhiệt kế.}
{Ôm kế.}
{\True Tốc kế.}
\loigiai{Tốc kế là dụng cụ dùng để đo vận tốc tức thời.}
\end{ex}

\dangbt{Thứ nguyên}
\begin{ex}
Đơn vị của lực trong hệ SI là
\choice
{km/h.}
{\True N.}
{m/s.}
{m/s2.}
\loigiai{Đơn vị của lực trong hệ SI là Newton (N).}
\end{ex}

\dangbt{Thứ nguyên}
\begin{ex}
Trong các lựa chọn sau, lựa chọn tất cả các đơn vị đều có trong hệ thống đo lường SI là
\choice
{Nhiệt độ (K), hiệu điệu thế (V), Cường độ sáng (Cd).}
{Lượng chất (mol), Chiều dài (m), lực (N).}
{Cường độ dòng điện (A), thời gian (s), công (J).}
{\True Chiều dài (m), thời gian (s), khối lượng (kg).}
\loigiai{Các đơn vị cơ bản trong hệ SI là mét (m), kilôgam (kg), giây (s), ampere (A), kelvin (K), mol (mol), candela (Cd).
A. Hiệu điện thế (V) không phải đơn vị cơ bản.
B. Lực (N) không phải đơn vị cơ bản.
C. Công (J) không phải đơn vị cơ bản.
D. Chiều dài (m), thời gian (s), khối lượng (kg) đều là các đơn vị cơ bản của hệ SI.}
\end{ex}

\dangbt{Các phương pháp nghiên cứu vật lí}
\begin{ex}
Đối tượng nghiên cứu chủ yếu của Vật lí là
\choice
{các dạng của vật chất, động lượng.}
{các dạng của vật chất, hạt nhân nguyên tử.}
{\True các dạng của vật chất, năng lượng.}
{các dạng của vật chất, công suất.}
\loigiai{Đối tượng nghiên cứu chủ yếu của Vật lí là các dạng của vật chất và năng lượng, cùng với sự tương tác giữa chúng.}
\end{ex}

\dangbt{Các phương pháp nghiên cứu vật lí}
\begin{ex}
Các phương pháp nghiên cứu chủ yếu của Vật lí bao gồm
\choice
{phương pháp thực nghiệm và phương pháp điều tra.}
{phương pháp mô hình hóa và phương pháp điều tra.}
{phương pháp lí thuyết và phương pháp mô hình hóa.}
{\True phương pháp thực nghiệm và phương pháp lí thuyết.}
\loigiai{Các phương pháp nghiên cứu chủ yếu của Vật lí bao gồm phương pháp thực nghiệm và phương pháp lí thuyết.}
\end{ex}

\dangbt{Quy trình đo}
\begin{ex}
Dưới đây là các bước của một trong các cách đo tốc độ chuyển động của một vật
(a) Dùng đồng hồ bấm giây đo thời gian t từ khi vật bắt đầu chuyển động từ điểm xuất phát tới điểm kết thúc.
(b) Xác định điểm xuất phát, điểm kết thúc.
(c) Dùng công thức v = s/t để tính tốc độ.
(d) Dùng thước đo độ dài của quãng đường s (tính từ điểm xuất phát đến điểm kết thúc).
Thứ tự đúng của các bước tiến hành là
\choice
{a, b, d, c.}
{\True b, d, a, c.}
{b, a, c, d.}
{a, c, b, d.}
\loigiai{Thứ tự đúng của các bước tiến hành là:
(b) Xác định điểm xuất phát, điểm kết thúc.
(d) Dùng thước đo độ dài của quãng đường s (tính từ điểm xuất phát đến điểm kết thúc).
(a) Dùng đồng hồ bấm giây đo thời gian t từ khi vật bắt đầu chuyển động từ điểm xuất phát tới điểm kết thúc.
(c) Dùng công thức v = s/t để tính tốc độ.}
\end{ex}

\dangbt{Quy trình đo}
\begin{ex}
Trong thí nghiệm thực hành đo tốc độ của vật chuyển động, sử dụng hai cổng quang điện để đo
\choice
{\True thời gian chuyển động của viên bi thép.}
{tốc độ trung bình của viên bi thép.}
{đường kính của viên bi thép.}
{tốc độ tức thời của viên bi thép.}
\loigiai{Hai cổng quang điện được sử dụng để đo thời gian mà vật đi qua giữa hai cổng, từ đó xác định thời gian chuyển động của viên bi thép trên một quãng đường nhất định.}
\end{ex}

\dangbt{Quy trình đo}
\begin{ex}
Những dụng cụ chính để đo tốc độ trung bình của viên bi gồm
\choice
{\True đồng hồ đo thời gian hiện số, cổng quang điện, viên bi, máng và thước thẳng.}
{đồng hồ đo thời gian hiện số, cổng quang điện, viên bi, máng và thước kẹp.}
{đồng hồ đo thời gian hiện số, cần rung, viên bi, máng và thước kẹp.}
{đồng hồ đo thời gian hiện số, cần rung, viên bi, máng và thước thẳng.}
\loigiai{Để đo tốc độ trung bình của viên bi, cần có đồng hồ đo thời gian (thường là đồng hồ hiện số kết hợp cổng quang điện), viên bi, máng để viên bi lăn, và thước thẳng để đo quãng đường.}
\end{ex}
\dangbt{Phép đo trực tiếp và phép đo gián tiếp}
\begin{ex}
Xét các phát biểu về phép đo trực tiếp sau.
\choiceTF
{Đo tốc độ trung bình của một vật bằng thước đo chiều dài và đồng hồ bấm giây là phép đo trực tiếp.}
{\True Đo chiều dài chiếc bút chì bằng thước đo chiều dài là phép đo trực tiếp.}
{Đo khối lượng của 5 quả táo bằng khối lượng những quả tạ có trọng lượng 2 N là phép đo trực tiếp.}
{\True Đo thời gian đi từ nhà đến trường bằng đồng hồ bấm giây là phép đo trực tiếp.}
\loigiai{a sai. Vì phải cần tính toán vận tốc trung bình sau khi đo. b đúng. c sai. Vì phải cộng khối lượng của các quả cân sau khi cân. d đúng.}
\end{ex}

\dangbt{Quy trình đo}
\begin{ex}
Quan sát hình bên, phát biểu đúng về độ chia nhỏ nhất và giới hạn đo là
\choice
{Độ chia nhỏ nhất của hai thước như nhau.}
{\True Giới hạn đo của hai thước bằng nhau.}
{Thước ở hình 1 có độ chia nhỏ nhất là 0,2 cm.}
{Thước ở hình 2 có độ chia nhỏ nhất là 0,2 cm.}
\loigiai{a sai. Hình 2: Thước có độ chia nhỏ nhất là 0,2 cm, sai số dụng cụ là 0,2 cm, khác với hình 1. b đúng. c sai. Hình 1: Thước có độ chia nhỏ nhất là 0,1 cm. d sai. Hình 2: Thước có độ chia nhỏ nhất là 0,2 cm.}
\end{ex}

\dangbt{Sai số}
\begin{ex}
Hai người cùng đo chiều dài của cánh cửa sổ, kết quả thu được như sau:
Người thứ nhất: $L_1 = (1,20 \pm 0,02)$ m
Người thứ hai: $L_2 = (1,21 \pm 0,03)$ m
\choiceTF
{Cả hai người đều dùng phép đo gián tiếp.}
{\True Sai số tỉ đối của phép đo của người thứ nhất là $1,67\%$.}
{\True Sai số tỉ đối của phép đo của người thứ hai là $2,48\%$.}
{Người thứ hai đo chính xác hơn người thứ nhất.}
\loigiai{a sai. Cả hai người đều dùng phép đo trực tiếp. b đúng. Sai số tỉ đối của phép đo của người thứ nhất $\delta_1 = \frac{0,02}{1,20} \times 100\% \approx 1,67\%$. c đúng. Sai số tỉ đối của phép đo của người thứ hai $\delta_2 = \frac{0,03}{1,21} \times 100\% \approx 2,48\%$. d sai. Do $\delta_1 < \delta_2$ nên người thứ nhất đo chính xác hơn người thứ hai.}
\end{ex}

\dangbt{Tính giá trị trung bình của đại lượng đo}
\begin{ex}
Bố trí một thí nghiệm dùng con lắc đơn để xác định gia tốc trọng trường. Các số liệu đo được cho ở bảng dưới đây:
\begin{tabular}{|c|c|c|c|}
\hline
Lần đo & Chiều dài dây treo (m) & Chu kỳ dao động (s) & Gia tốc trọng trường (m/s$^2$) \\
\hline
1 & 1,2 & 2,19 & 9,8776 \\
2 & 0,9 & 1,90 & 9,8423 \\
3 & 1,3 & 2,29 & 9,7866 \\
\hline
\end{tabular}
\choiceTF
{Chu kì dao động trung bình xấp xỉ bằng 2,13 s.}
{\True Giá trị trung bình của gia tốc là 9,8355 m/s$^2$.}
{Sai số của gia tốc trọng trường là 0,045 m/s$^2$.}
{\True Kết quả phép đo được viết là $g = (9,8355 \pm 0,045)$ m/s$^2$.}
\loigiai{a sai. Chu kì dao động trung bình $\bar{T} = \frac{2,19 + 1,90 + 2,29}{3} \approx 2,127$ s. b đúng. Giá trị trung bình gia tốc $\bar{g} = \frac{9,8776 + 9,8423 + 9,7866}{3} \approx 9,8355$ m/s$^2$. c sai. Sai số của gia tốc trọng trường $\Delta g = \frac{|9,8776 - 9,8355| + |9,8423 - 9,8355| + |9,7866 - 9,8355|}{3} \approx \frac{0,0421 + 0,0068 + 0,0489}{3} \approx 0,0326$ m/s$^2$. d đúng. Kết quả phép đo được viết là $g = (9,8355 \pm 0,0326)$ m/s$^2$.}
\end{ex}

\dangbt{Sai số}
\begin{ex}
Tiến hành đo thời gian chuyển động của một viên bi ta thu được số liệu như bảng sau:
\begin{tabular}{|c|c|c|c|}
\hline
Lần đo & Lần 1 & Lần 2 & Lần 3 \\
\hline
Thời gian t (s) & 1,553 & 1,549 & 1,556 \\
\hline
\end{tabular}
\choiceTF
{\True Thời gian trung bình viên bi chuyển động là 1,553 s.}
{Sai số tuyệt đối ứng với lần đo thứ 2 là 0,003 s.}
{\True Sai số tuyệt đối trung bình là 0,0024 s.}
{Kết quả phép đo là $t = (1,553 \pm 0,0024)$ s.}
\loigiai{a đúng. Giá trị thời gian trung bình $\bar{t} = \frac{1,553 + 1,549 + 1,556}{3} = 1,55266... \approx 1,553$ s. b sai. Sai số tuyệt đối của lần đo thứ 2 là $|1,549 - 1,553| = 0,004$ s. c đúng. Sai số tuyệt đối trung bình là $\overline{\Delta t} = \frac{|1,553 - 1,553| + |1,549 - 1,553| + |1,556 - 1,553|}{3} = \frac{0 + 0,004 + 0,003}{3} \approx 0,00233... \approx 0,0024$ s. d sai. Kết quả phép đo là $t = (1,553 \pm 0,002)$ s (làm tròn sai số đến 1 chữ số có nghĩa).}
\end{ex}

\dangbt{Sai số}
\begin{ex}
Để xác định tốc độ của một vật chuyển động thẳng đều, một người đã đo quãng đường vật đi được bằng $(18 \pm 0,3)$ m trong khoảng thời gian là $(3,0 \pm 0,2)$ s.
\choiceTF
{\True Vận tốc trung bình của xe là 6 m/s.}
{Sai số tỉ đối của phép đo là $0,083$.}
{Tốc độ lớn nhất của vật là 6,1 m/s.}
{\True Kết quả phép đo vận tốc của vật là $v = (6,0 \pm 0,5)$ m/s.}
\loigiai{a đúng. Vận tốc trung bình $\bar{v} = \frac{\bar{s}}{\bar{t}} = \frac{18}{3,0} = 6$ m/s. b sai. Sai số tỉ đối được xác định $\delta_v = \delta_s + \delta_t = \frac{0,3}{18} + \frac{0,2}{3,0} \approx 0,0167 + 0,0667 \approx 0,0834$. c sai. Tốc độ lớn nhất của vật là $v_{max} = \frac{s_{max}}{t_{min}} = \frac{18 + 0,3}{3,0 - 0,2} = \frac{18,3}{2,8} \approx 6,536$ m/s. d đúng. Cách viết kết quả đo là $v = (\bar{v} \pm \Delta v)$ với $\Delta v = \bar{v} \times \delta_v = 6 \times 0,0834 \approx 0,5004$ m/s. Vậy nên $v = (6,0 \pm 0,5)$ m/s.}
\end{ex}

\dangbt{Sai số}
\begin{ex}
Một người đo chiều cao một người khác trong ba lần được là 160 cm, 161 cm, 162 cm.
\choiceTF
{\True Chiều cao trung bình của người đó là 161 cm.}
{Sai số tuyệt đối trung bình là 0,05 m.}
{\True Sai số tỉ đối của phép đo là $0,62\%$.}
{Kết quả phép đo là $h = (161 \pm 1)$ cm.}
\loigiai{a đúng. Chiều cao trung bình $\bar{h} = \frac{160 + 161 + 162}{3} = 161$ cm. b sai. Sai số tuyệt đối ứng với mỗi lần đo: $\Delta h_1 = |160 - 161| = 1$ cm, $\Delta h_2 = |161 - 161| = 0$ cm, $\Delta h_3 = |162 - 161| = 1$ cm. Sai số tuyệt đối trung bình là $\overline{\Delta h} = \frac{1 + 0 + 1}{3} = \frac{2}{3} \approx 0,67$ cm $\approx 0,0067$ m. c đúng. Sai số tuyệt đối của phép đo $\Delta h = \overline{\Delta h} \approx 0,67$ cm. Sai số tỉ đối của phép đo là $\delta_h = \frac{\Delta h}{\bar{h}} = \frac{0,67}{161} \times 100\% \approx 0,42\%$. d sai. Kết quả của phép đo là: $h = (161 \pm 0,7)$ cm.}
\end{ex}

\dangbt{Sai số}
\begin{ex}
Một em học sinh đo (3 lần) chiều dài một quyển sách được các kết quả là 22,20 cm, 22,31 cm, 22,62 cm.
\choiceTF
{Sai số tuyệt đối ứng với lần đo thứ nhất là 5,20 cm.}
{\True Chiều dài trung bình của quyển sách là 22,377 cm.}
{Sai số tuyệt đối trung bình là 0,18 cm.}
{\True Thực hiện đo chiều dài một quyển sách 3 lần để giảm sai số.}
\loigiai{a sai. Chiều dài trung bình $\bar{L} = \frac{22,20 + 22,31 + 22,62}{3} = 22,3766... \approx 22,377$ cm. Sai số tuyệt đối ứng với lần đo thứ nhất là $|22,20 - 22,377| \approx 0,177$ cm. b đúng. c sai. Sai số tuyệt đối trung bình là $\overline{\Delta L} = \frac{|22,20 - 22,377| + |22,31 - 22,377| + |22,62 - 22,377|}{3} \approx \frac{0,177 + 0,067 + 0,243}{3} \approx 0,162$ cm. d đúng.}
\end{ex}

\dangbt{Sai số}
\begin{ex}
Một học sinh đo chiều dài $d$ của một bàn học, kết quả thu được là $d = (120 \pm 1)$ cm.
\choiceTF
{\True Chiều dài trung bình của một bàn học là 120 cm.}
{Sai số của phép đo là 1 m.}
{\True Sai số tỉ đối của phép đo là 0,83\%.}
{Kết quả phép đo có thể được viết là $d = (1,2 \pm 0,01)$ m.}
\loigiai{a đúng. Giá trị gần đúng nhất của đại lượng cần đo là $\bar{d} = 120$ cm. b sai. Sai số của phép đo là $\Delta d = 1$ cm. c đúng. Sai số tỉ đối trong phép đo $\delta_d = \frac{\Delta d}{\bar{d}} = \frac{1}{120} \times 100\% \approx 0,83\%$. d sai. Kết quả phép đo có thể được viết là $d = (1,20 \pm 0,01)$ m.}
\end{ex}

\dangbt{Sai số}
\begin{ex}
Dùng thước có độ chia nhỏ nhất 0,02 cm để đo 4 lần chiều dài của một vật hình trụ. Kết quả đo được cho ở bảng sau:
\begin{tabular}{|c|c|c|c|c|}
\hline
Lần đo & 1 & 2 & 3 & 4 \\
\hline
Chiều dài $l$ (cm) & 3,29 & 3,36 & 3,32 & 3,27 \\
\hline
\end{tabular}
\choiceTF
{\True Chiều dài trung bình của vật là 3,31 cm.}
{Sai số tuyệt đối của phép đo là 0,02 cm.}
{\True Sai số tương đối của phép đo là 1,21\%.}
{Biểu diễn kết quả đo: $l = (3,31 \pm 0,02)$ cm.}
\loigiai{a đúng. Chiều dài trung bình $\bar{l} = \frac{3,29 + 3,36 + 3,32 + 3,27}{4} = 3,31$ cm. b sai. Sai số tuyệt đối ứng với mỗi lần đo: $\Delta l_1 = |3,29 - 3,31| = 0,02$ cm, $\Delta l_2 = |3,36 - 3,31| = 0,05$ cm, $\Delta l_3 = |3,32 - 3,31| = 0,01$ cm, $\Delta l_4 = |3,27 - 3,31| = 0,04$ cm. Sai số tuyệt đối trung bình $\overline{\Delta l} = \frac{0,02 + 0,05 + 0,01 + 0,04}{4} = \frac{0,12}{4} = 0,03$ cm. Sai số dụng cụ $\Delta l_{dc} = \frac{0,02}{2} = 0,01$ cm. Sai số tuyệt đối $\Delta l = \overline{\Delta l} + \Delta l_{dc} = 0,03 + 0,01 = 0,04$ cm. c đúng. Sai số tương đối $\delta_l = \frac{\Delta l}{\bar{l}} = \frac{0,04}{3,31} \times 100\% \approx 1,21\%$. d sai. Biểu diễn kết quả đo $l = (3,31 \pm 0,04)$ cm.}
\end{ex}

\dangbt{Sai số}
\begin{ex}
Để đo tốc độ tức thời của một vật, người ta dùng một thước có độ chia nhỏ nhất là 2 mm, thực hiện đo quãng đường mà vật đi thì luôn thu được kết quả 40 cm. Dùng một đồng hồ đo thời gian có độ chia nhỏ nhất là 0,002 s để đo thời gian vật chuyển động được quãng đường trên thì thu được bảng số liệu sau:
\begin{tabular}{|c|c|c|c|c|c|}
\hline
Lần đo thứ & 1 & 2 & 3 & 4 & 5 \\
\hline
t (s) & 0,401 & 0,399 & 0,403 & 0,395 & 0,402 \\
\hline
\end{tabular}
\choiceTF
{\True Giá trị trung bình của thời gian là 0,4 s.}
{Vận tốc trung bình là 1,0 m/s.}
{\True Kết quả đo tốc độ trung bình là $v = (1,0 \pm 0,02)$ m/s.}
{Phép đo này là phép đo trực tiếp.}
\loigiai{a đúng. Giá trị trung bình của thời gian $\bar{t} = \frac{0,401 + 0,399 + 0,403 + 0,395 + 0,402}{5} = \frac{2,0}{5} = 0,4$ s. b sai. Giá trị trung bình của tốc độ tức thời là $\bar{v} = \frac{\bar{s}}{\bar{t}} = \frac{40 \text{ cm}}{0,4 \text{ s}} = \frac{0,4 \text{ m}}{0,4 \text{ s}} = 1,0$ m/s. c đúng. Sai số tuyệt đối trung bình của 5 lần đo thời gian: $\overline{\Delta t} = \frac{|0,401 - 0,4| + |0,399 - 0,4| + |0,403 - 0,4| + |0,395 - 0,4| + |0,402 - 0,4|}{5} = \frac{0,001 + 0,001 + 0,003 + 0,005 + 0,002}{5} = \frac{0,012}{5} = 0,0024$ s. Sai số dụng cụ $\Delta t_{dc} = \frac{0,002}{2} = 0,001$ s. Sai số tuyệt đối của phép đo thời gian $\Delta t = \overline{\Delta t} + \Delta t_{dc} = 0,0024 + 0,001 = 0,0034$ s. Sai số tuyệt đối của quãng đường $\Delta s = \frac{2 \text{ mm}}{2} = 1 \text{ mm} = 0,001$ m. Sai số tỉ đối của tốc độ $\delta_v = \delta_s + \delta_t = \frac{\Delta s}{\bar{s}} + \frac{\Delta t}{\bar{t}} = \frac{0,001}{0,4} + \frac{0,0034}{0,4} = 0,0025 + 0,0085 = 0,011$. Sai số tuyệt đối của tốc độ tức thời là $\Delta v = \bar{v} \times \delta_v = 1,0 \times 0,011 = 0,011$ m/s. Kết quả đo tốc độ trung bình là $v = (1,0 \pm 0,01)$ m/s. d sai. Đây là phép đo gián tiếp thông qua việc đo trực tiếp quãng đường đi được và thời gian đi của vật.}
\end{ex}

\dangbt{Sai số}
\begin{ex}
Cho bảng số liệu thể hiện kết quả đo khối lượng của một túi trái cây bằng cân đồng hồ. Em hãy xác định sai số tuyệt đối ứng với từng lần đo, sai số tuyệt đối và sai số tương đối của phép đo. Biết sai số dụng cụ là 0,1 kg.
\begin{tabular}{|c|c|c|}
\hline
Lần đo & m (Kg) & $\Delta$m (kg) \\
\hline
1 & 4,2 & \\
2 & 4,4 & \\
3 & 4,4 & \\
4 & 4,2 & \\
\hline
Trung bình & $\bar{m} = ?$ & $\overline{\Delta m} = ?$ \\
\hline
\end{tabular}
\choiceTF
{\True Giá trị trung bình khối lượng của túi trái cây là 4,3 kg.}
{Sai số tuyệt đối trung bình của phép đo là 0,05 kg.}
{Sai số tuyệt đối của phép đo là 0,1 kg.}
{\True Sai số tương đối của phép đo là 3,49\%.}
\loigiai{a đúng. Giá trị trung bình khối lượng của túi trái cây là $\bar{m} = \frac{4,2 + 4,4 + 4,4 + 4,2}{4} = \frac{17,2}{4} = 4,3$ kg. b sai. Sai số tuyệt đối ứng với mỗi lần đo: $\Delta m_1 = |4,2 - 4,3| = 0,1$ kg, $\Delta m_2 = |4,4 - 4,3| = 0,1$ kg, $\Delta m_3 = |4,4 - 4,3| = 0,1$ kg, $\Delta m_4 = |4,2 - 4,3| = 0,1$ kg. Sai số tuyệt đối trung bình của phép đo $\overline{\Delta m} = \frac{0,1 + 0,1 + 0,1 + 0,1}{4} = 0,1$ kg. c sai. Sai số dụng cụ $\Delta m_{dc} = \frac{0,1}{2} = 0,05$ kg. Sai số tuyệt đối của phép đo $\Delta m = \overline{\Delta m} + \Delta m_{dc} = 0,1 + 0,05 = 0,15$ kg. d đúng. Sai số tương đối của phép đo là $\delta_m = \frac{\Delta m}{\bar{m}} = \frac{0,15}{4,3} \times 100\% \approx 3,49\%$.}
\end{ex}

\dangbt{Sai số}
\begin{ex}
Trong giờ thực hành, một học sinh đo chu kì dao động của con lắc đơn bằng đồng hồ bấm giây. Kết quả 5 lần đo được cho ở bảng sau.
\begin{tabular}{|c|c|c|c|c|c|}
\hline
Lần đo & 1 & 2 & 3 & 4 & 5 \\
\hline
Chu kì T (s) & 2,01 & 2,11 & 2,05 & 2,03 & 2,00 \\
\hline
\end{tabular}
Cho biết thang chia nhỏ nhất của đồng hồ là 0,02 s.
\choiceTF
{\True Giá trị trung bình của chu kì dao động là 2,04 s.}
{Sai số tuyệt đối trung bình của phép đo là 0,036 s.}
{Sai số tuyệt đối của phép đo là 0,046 s.}
{\True Kết quả đo chu kì là $T = (2,04 \pm 0,04)$ s.}
\loigiai{a đúng. Giá trị trung bình của chu kì dao động $\bar{T} = \frac{2,01 + 2,11 + 2,05 + 2,03 + 2,00}{5} = \frac{10,2}{5} = 2,04$ s. b sai. Sai số tuyệt đối của mỗi lần đo: $\Delta T_1 = |2,01 - 2,04| = 0,03$ s, $\Delta T_2 = |2,11 - 2,04| = 0,07$ s, $\Delta T_3 = |2,05 - 2,04| = 0,01$ s, $\Delta T_4 = |2,03 - 2,04| = 0,01$ s, $\Delta T_5 = |2,00 - 2,04| = 0,04$ s. Sai số tuyệt đối trung bình của phép đo: $\overline{\Delta T} = \frac{0,03 + 0,07 + 0,01 + 0,01 + 0,04}{5} = \frac{0,16}{5} = 0,032$ s. c sai. Sai số dụng cụ $\Delta T_{dc} = \frac{0,02}{2} = 0,01$ s. Sai số tuyệt đối của phép đo $\Delta T = \overline{\Delta T} + \Delta T_{dc} = 0,032 + 0,01 = 0,042$ s. d đúng. Kết quả đo chu kì $T = (2,04 \pm 0,04)$ s.}
\end{ex}

\dangbt{Sai số}
\begin{ex}
Dùng thước kẹp có độ chia nhỏ nhất là 1mm để đo đường kính của viên bi thu được kết quả là $D = (15,0 \pm 0,5)$ mm. Sai số tỉ đối của phép đo trên là bao nhiêu phần trăm?
\shortans{3.33}
\loigiai{Sai số tỉ đối của phép đo trên $\delta_D = \frac{\Delta D}{\bar{D}} = \frac{0,5}{15,0} \times 100\% \approx 3,33\%$.}
\end{ex}

\dangbt{Sai số}
\begin{ex}
Sử dụng đồng hồ đa năng để đo cường độ dòng điện qua điện trở R, kết quả của 5 lần đo lần lượt là: 4,988 A, 5,012 A, 4,910 A, 5,025 A, 5,057 A. Sai số tuyệt đối của phép đo là bao nhiêu ampe?
\shortans{0.05}
\loigiai{Giá trị trung bình $\bar{I} = \frac{4,988 + 5,012 + 4,910 + 5,025 + 5,057}{5} = \frac{25,092}{5} = 5,0184$ A.
Sai số tuyệt đối ứng với mỗi lần đo:
$\Delta I_1 = |4,988 - 5,0184| = 0,0304$ A
$\Delta I_2 = |5,012 - 5,0184| = 0,0064$ A
$\Delta I_3 = |4,910 - 5,0184| = 0,1084$ A
$\Delta I_4 = |5,025 - 5,0184| = 0,0066$ A
$\Delta I_5 = |5,057 - 5,0184| = 0,0386$ A
Sai số tuyệt đối trung bình là $\overline{\Delta I} = \frac{0,0304 + 0,0064 + 0,1084 + 0,0066 + 0,0386}{5} = \frac{0,1904}{5} = 0,03808$ A.
Sai số tuyệt đối của phép đo $\Delta I = \overline{\Delta I} \approx 0,04$ A (làm tròn đến 1 chữ số có nghĩa).
(Nếu không có sai số dụng cụ, thì sai số tuyệt đối của phép đo là sai số tuyệt đối trung bình, làm tròn đến 1 chữ số có nghĩa, tức là 0,04 A. Nếu làm tròn đến 2 chữ số thập phân thì là 0,05 A.)}
\end{ex}

\dangbt{Sai số}
\begin{ex}
Trong thí nghiệm tính vận tốc của vật chuyển động thẳng đều, kết quả đo quãng đường $s = (8,255 \pm 0,245)$ m và thời gian $t = (4,025 \pm 0,120)$ s. Sai số tỉ đối của phép đo là bao nhiêu m/s? (kết quả làm tròn đến hai chữ số thập phân).
\shortans{0.06}
\loigiai{Giá trị trung bình của vận tốc $\bar{v} = \frac{\bar{s}}{\bar{t}} = \frac{8,255}{4,025} \approx 2,051$ m/s.
Sai số tỉ đối được xác định $\delta_v = \delta_s + \delta_t = \frac{0,245}{8,255} + \frac{0,120}{4,025} \approx 0,02968 + 0,02981 \approx 0,05949$.
Sai số tuyệt đối của vận tốc $\Delta v = \bar{v} \times \delta_v = 2,051 \times 0,05949 \approx 0,122$ m/s.
Kết quả làm tròn đến hai chữ số thập phân là 0,12 m/s.}
\end{ex}

\dangbt{Sai số}
\begin{ex}
Dùng thước có độ chia nhỏ nhất 0,02 cm để đo 4 lần chiều dài của một vật hình trụ. Kết quả đo được cho ở bảng sau:
\begin{tabular}{|c|c|c|c|c|}
\hline
Lần đo & 1 & 2 & 3 & 4 \\
\hline
Chiều dài $l$ (cm) & 3,29 & 3,36 & 3,32 & 3,27 \\
\hline
\end{tabular}
Sai số tương đối của phép đo này là bao nhiêu phần trăm? (kết quả làm tròn đến hai chữ số thập phân).
\shortans{1.21}
\loigiai{Chiều dài trung bình $\bar{l} = \frac{3,29 + 3,36 + 3,32 + 3,27}{4} = 3,31$ cm.
Sai số tuyệt đối ứng với mỗi lần đo:
$\Delta l_1 = |3,29 - 3,31| = 0,02$ cm
$\Delta l_2 = |3,36 - 3,31| = 0,05$ cm
$\Delta l_3 = |3,32 - 3,31| = 0,01$ cm
$\Delta l_4 = |3,27 - 3,31| = 0,04$ cm
Sai số tuyệt đối trung bình $\overline{\Delta l} = \frac{0,02 + 0,05 + 0,01 + 0,04}{4} = \frac{0,12}{4} = 0,03$ cm.
Sai số dụng cụ $\Delta l_{dc} = \frac{0,02}{2} = 0,01$ cm.
Sai số tuyệt đối $\Delta l = \overline{\Delta l} + \Delta l_{dc} = 0,03 + 0,01 = 0,04$ cm.
Sai số tương đối $\delta_l = \frac{\Delta l}{\bar{l}} = \frac{0,04}{3,31} \times 100\% \approx 1,208\% \approx 1,21\%$.}
\end{ex}

\dangbt{Sai số}
\begin{ex}
Tiến hành đo chuyển động của một viên bi khi được bắn ra xa (bỏ qua sai số dụng cụ), ta thu được số liệu như bảng sau:
\begin{tabular}{|c|c|c|c|c|}
\hline
N & s (m) & $\Delta$s (m) & t (s) & $\Delta$t (s) \\
\hline
1 & 0,1 & & 0,02 & \\
2 & 0,12 & & 0,023 & \\
3 & 0,11 & & 0,022 & \\
4 & 0,123 & & 0,021 & \\
5 & 0,1 & & 0,02 & \\
\hline
Trung bình & & & & \\
\hline
\end{tabular}
Sai số tỉ đối của phép đo là bao nhiêu phần trăm? (kết quả làm tròn đến hai chữ số thập phân).
\shortans{8.11}
\loigiai{
\begin{tabular}{|c|c|c|c|c|}
\hline
N & s (m) & $\Delta$s (m) & t (s) & $\Delta$t (s) \\
\hline
1 & 0,1 & 0,011 & 0,02 & 0,001 \\
2 & 0,12 & 0,009 & 0,023 & 0,002 \\
3 & 0,11 & 0,001 & 0,022 & 0,001 \\
4 & 0,123 & 0,012 & 0,021 & 0,000 \\
5 & 0,1 & 0,011 & 0,02 & 0,001 \\
\hline
Trung bình & 0,111 & 0,009 & 0,021 & 0,001 \\
\hline
\end{tabular}
Giá trị trung bình của vận tốc $\bar{v} = \frac{\bar{s}}{\bar{t}} = \frac{0,111}{0,021} \approx 5,2857$ m/s.
Sai số tỉ đối được xác định $\delta_v = \delta_s + \delta_t = \frac{\overline{\Delta s}}{\bar{s}} + \frac{\overline{\Delta t}}{\bar{t}} = \frac{0,009}{0,111} + \frac{0,001}{0,021} \approx 0,08108 + 0,04762 \approx 0,1287$.
Sai số tỉ đối của phép đo là $\delta_v \times 100\% \approx 12,87\%$.
(Lưu ý: Lời giải gốc có vẻ tính sai số tỉ đối của từng đại lượng rồi cộng lại, nhưng kết quả cuối cùng không khớp với cách tính đó. Nếu tính $\Delta v = \bar{v} \times (\frac{\overline{\Delta s}}{\bar{s}} + \frac{\overline{\Delta t}}{\bar{t}})$ thì $\Delta v = 5,2857 \times (\frac{0,009}{0,111} + \frac{0,001}{0,021}) \approx 5,2857 \times (0,08108 + 0,04762) \approx 5,2857 \times 0,1287 \approx 0,680$ m/s.
Sai số tỉ đối của phép đo là $\frac{\Delta v}{\bar{v}} \times 100\% \approx \frac{0,680}{5,2857} \times 100\% \approx 12,86\%$.
Nếu tính theo lời giải gốc: $\delta_v = \frac{0,009}{0,111} \times 100\% + \frac{0,001}{0,021} \times 100\% \approx 8,11\% + 4,76\% = 12,87\%$.
Có vẻ lời giải gốc đã tính sai số tỉ đối của từng đại lượng rồi cộng lại, nhưng kết quả cuối cùng lại là 8.11%.)
Theo lời giải gốc, $\delta_s = \frac{0,009}{0,111} \approx 0,08108$, $\delta_t = \frac{0,001}{0,021} \approx 0,04762$.
$\delta_v = \delta_s + \delta_t \approx 0,08108 + 0,04762 \approx 0,1287$.
Nếu kết quả là 8.11% thì có thể là chỉ lấy $\delta_s$.
Giả sử đề hỏi sai số tỉ đối của quãng đường: $\frac{0,009}{0,111} \times 100\% \approx 8,11\%$.
}
\end{ex}

\dangbt{Sai số}
\begin{ex}
Một người đo chiều cao một người khác trong ba lần được là 160 cm, 161 cm, 162 cm. Sai số tỉ đối của phép đo gần đúng là bao nhiêu phần trăm? (kết quả làm tròn đến hai chữ số thập phân).
\shortans{0.42}
\loigiai{Chiều cao trung bình $\bar{h} = \frac{160 + 161 + 162}{3} = 161$ cm.
Sai số tuyệt đối ứng với mỗi lần đo:
$\Delta h_1 = |160 - 161| = 1$ cm
$\Delta h_2 = |161 - 161| = 0$ cm
$\Delta h_3 = |162 - 161| = 1$ cm
Sai số tuyệt đối trung bình là $\overline{\Delta h} = \frac{1 + 0 + 1}{3} = \frac{2}{3} \approx 0,6667$ cm.
Sai số tuyệt đối của phép đo $\Delta h = \overline{\Delta h} \approx 0,67$ cm (làm tròn đến 2 chữ số thập phân).
Sai số tỉ đối của phép đo là $\delta_h = \frac{\Delta h}{\bar{h}} = \frac{0,67}{161} \times 100\% \approx 0,416\% \approx 0,42\%$.}
\end{ex}

\dangbt{Sai số}
\begin{ex}
Tiến hành đo thời gian chuyển động của một viên bi ta thu được số liệu như bảng sau:
\begin{tabular}{|c|c|c|c|c|c|}
\hline
Lần đo & Lần 1 & Lần 2 & Lần 3 & Giá trị t trung bình & Sai số $\Delta t$ \\
\hline
Thời gian t (s) & 1,553 & 1,549 & 1,556 & & \\
\hline
\end{tabular}
Sai số tuyệt đối trung bình của viên bi là bao nhiêu giây? (kết quả làm tròn đến ba chữ số thập phân).
\shortans{0.002}
\loigiai{Giá trị thời gian trung bình $\bar{t} = \frac{1,553 + 1,549 + 1,556}{3} = 1,55266... \approx 1,553$ s.
Sai số tuyệt đối ứng với mỗi lần đo:
$\Delta t_1 = |1,553 - 1,553| = 0$ s
$\Delta t_2 = |1,549 - 1,553| = 0,004$ s
$\Delta t_3 = |1,556 - 1,553| = 0,003$ s
Sai số tuyệt đối trung bình là $\overline{\Delta t} = \frac{0 + 0,004 + 0,003}{3} = \frac{0,007}{3} \approx 0,00233... \approx 0,002$ s.}
\end{ex}

\dangbt{Sai số}
\begin{ex}
Biết rằng đại lượng vật lí A được đo gián tiếp thông qua 2 đại lượng vật lí B và C bằng công thức toán học $A = B^2 C$. Kết quả đo đại lượng B và C lần lượt là $B = (2,0 \pm 0,1)$ và $C = (4,0 \pm 0,2)$. Sai số tương đối của phép đo là bao nhiêu? (kết quả làm tròn đến hai chữ số thập phân).
\shortans{0.10}
\loigiai{Giá trị A trung bình $\bar{A} = \bar{B}^2 \bar{C} = (2,0)^2 \times 4,0 = 4 \times 4 = 16$.
Sai số tương đối $\delta_A = 2\delta_B + \delta_C = 2 \times \frac{0,1}{2,0} + \frac{0,2}{4,0} = 2 \times 0,05 + 0,05 = 0,10 + 0,05 = 0,15$.
Sai số tuyệt đối $\Delta A = \bar{A} \times \delta_A = 16 \times 0,15 = 2,4$.
Kết quả phép đo $A = (16 \pm 2,4)$.
Sai số tương đối của phép đo là $\delta_A = 0,15$.
(Lời giải gốc đưa ra 0.10, có thể là sai số tương đối của B hoặc C, hoặc là một cách làm tròn khác.)
Nếu làm tròn đến hai chữ số thập phân thì 0,15.}
\end{ex}
